
% !TEX root = lin-prog-main.tex

\section{Introduction}
\label{sec:intro}
Applications of (integer) linear programming, and linear algebra in general, 
are ubiquitous in computer science (see e.g.~\cite{SilvaTC96,Cormenbook,Colcombet15}),
including recent and potential future applications to analysis of data-enriched 
models~\cite{HLT17,HL18,GSAH19,BKM21}.
%\slawek{any other good references?}
%\arka{Maybe we don't need to give a reference here}
%\PH{linear programing, computational geometry, 
%computer graphics, I think we can add Cormen Introduction to algorithms, on }
%
Whenever finite (integer) linear programs arise in analysis of finite models of computation, 
   orbit-finite (integer) linear programs arise naturally in data-enriched versions of these models.
   For example, in decision problems for data Petri nets, such as reachability \cite{Nets-with-Tokens-which-Carry-Data} or
   continuous reachability \cite{GSAH19}; or in process mining \cite{WDHS08}.
Similar approach seems applicable also to
   structural properties or termination time of data Petri nets; and to  
   learning of probabilistic automata with registers.
%    \slawek{we need more citations here}. 
   
This paper is a continuation of the study of \emph{orbit-finite} systems of linear equations~\cite{GHL22},
i.e., systems which are infinite but finite up to permutations.
In this setting one fixes a countably infinite set $\A$, %=\set{\at 1, \at 2, \at 3, \ldots}$, 
whose elements are called \emph{atoms} (or data values)~\cite{atombook,Pitts:book}, 
assuming that atoms can only be accessed in a very limited way, namely can only be tested for equality.
Starting from atoms one builds a hierarchy of sets which are \emph{orbit-finite}: they are infinite, but finite up to permutations of atoms.
Along these lines, we study orbit-finite sets of linear inequalities, over an orbit-finite set of unknowns.
%, assuming that the only operations on atoms are (dis)equality tests.

%\para{Contribution}

%\piotrek{There is some logical gap here. How are orbit finite sets and their internal structures connected to orbit finite system of equations?}

The main result of~\cite{GHL22} is a decision procedure to check if a given orbit-finite system
of equations is solvable.
This result is general and applies to solvability over a wide range of commutative rings, 
in particular to real and integer solvability.
%
%\arka{I think we should say it works for commutative rings}
%
%
In this paper we do a next step and extend the setting from equations to inequalities.
Our goal is algorithmic solvability of orbit-finite systems
of inequalities, but also optimisation 
%(minimisation/maximisation) \arka{\st{(minimisation/maximisation)}}
of linear objective functions
over solution sets of such systems.
%(the former one reduces to the latter one, in a standard way).
%
%\arka{I am not sure if we should call it a reduction.
%The latter one simply has the former one as a subproblem.}
%
We call this problem \emph{orbit-finite (integer) linear programming}
(depending on whether the considered solutions are real or integer).


\begin{slexample} \label{ex:lp} %\rm
For illustration, consider the set $\A$ as unknowns,
and the infinite system of constraints given by 
an infinite matrix whose rows and columns are indexed by $\A$:
%
\begin{align} \label{eq:matr}
%& \qquad \begin{matrix} \ \cdots &  \a  & \cdots\ \end{matrix} \\
& 
%\begin{matrix}
%\vdots \\
%\b \\
%\vdots \\
%\end{matrix}
%\
\begin{bmatrix}
\ 0 & 1 & 1 & \cdots \ \\
\ 1 & 0 & 1 & \cdots \ \\
\ 1 & 1 & 0 & \cdots \ \\
 \vdots     & \vdots & \vdots & \ddots \   
\end{bmatrix}
\cdot
%\begin{bmatrix}
%\vdots \\
%\a \\
%\vdots \\
%\\
\ \vr x %\\
%\\
%\end{bmatrix}
\quad
\geq
\quad
\begin{bmatrix}
\, 1 \, \\
\, 1 \, \\
\, 1 \, \\
 \vdots 
\end{bmatrix}
\end{align}
%
Alternatively, one can write the infinite set of non-strict inequalities over unknowns $\a\in\A$, indexed by atoms $\b\in\A$:
%
\begin{align} \label{eq:constr}
\sum_{\a \in \A\setminus\set{\b}} \a \ \geq \ 1 \quad (\b \in \A).
\end{align}
%
Any permutation $\pi: \A \to \A$ induces a permutation of the inequalities by sending 
\[
\sum_{\a \in \A\setminus\set{\b}} \a \ \geq \ 1\qquad \stackrel{\pi}{\longmapsto}
%\text{ to}
\quad 
\sum_{\a \in \A\setminus\set{\pi(\b)}} \a  \ \geq \ 1,
\]
but the whole system~\eqref{eq:constr} is invariant under permutations of atoms.
Furthermore, up to permutations of atoms the system consists of just one equation -- it is one \emph{orbit};
in the sequel we consider orbit-finite systems (finite unions of orbits).
Likewise, the matrix $\A\times\A\to\R$ in \eqref{eq:matr} consists, up to permutation of atoms, of just two entries.
Indeed, its domain $\A\times\A$ is a union of two orbits: 
$\setof{\pair \a\b}{\a=\b}$ and $\setof{\pair \a \b}{\a\neq \b}$, and the matrix is constant
inside each orbit.
It is therefore invariant under permutations of atoms.


%For the sake of example, suppose that we seek for solutions $\vr x : \A\to\R$ assigning non-zero values to finitely many unknowns only
%(otherwise the sums on left-hand sides of inequalities could be undefined).
%
%\arka{I think we can skip the above sentence}

The system~\eqref{eq:constr} is solvable. 
%\red{By this we mane that there is a valuation of the set of unknowns, such that in each inequality the infinite sum 
%is well defined and the inequality holds.} 
For example, given $n>1$  atoms $S = \set{\a_1, \ldots, \a_n} \subseteq \A$, 
the vector $\vr{x}_n : \A \to \R$ defined by:
\[
\vr x_n(\a) = \frac 1 {n-1} \text{\ \ if \ }\a\in S, \qquad %and
\vr x_n(\a) = 0 \text{\ \ if \ }\a\notin S,
%\vr{x}(\g) =
%\begin{cases}
%1 & \text{ if }\g \in \{\a,\b\}, \\
%0 & \text{ otherwise.} 
%\end{cases}
\]
is a solution since the left-hand side of~\eqref{eq:constr} sums up to
$1$ if $\b\in S$, and to $\frac n {n-1}$ if $\b\notin S$.
%\[
%\sum_{\a \neq \b} \vr{x}(\a)  =
%\begin{cases}
%n-1 & \text{ if }\b \in S, \\
%n & \text{ otherwise.}
%\end{cases}
%\]
\end{slexample}

%Orbit-finite system of equations are solvable by a reduction to solvability of a finite number of classical finite systems~\cite{GHL22}.
%This strategy does not seem to extend to inequalities, and to linear programming, i.e., 
%the  optimisation of a linear objective function subject to 
%%over the solution set of 
%an orbit-finite system of inequality constraints.
%\arka{Maybe the previous statement is not true with the recent changes}
%
Orbit-finite linear programming, i.e., 
optimisation of a linear objective function subject to an orbit-finite system of inequality constraints,
faces phenomena not present
in the classical setting.
For instance, as illustrated in the next example, 
the objective function may not achieve its optimum over solutions of non-strict inequalities.

\begin{slexample} \label{ex:lpmin} %\rm
Suppose that we aim at \emph{minimization}, with respect to the constraints~\eqref{eq:constr},
of the value of the objective function:
\begin{align} \label{eq:objf}
S(\vr x) \ = \ 2\cdot \sum_{\a\in\A} \vr x(\a).
\end{align}
%(note that the infinite sum is well-defined assuming $\vr x$ assigns zero to almost all unknowns).
%
%\arka{The above statement is not crucial for this example.
%Moreover it does not make sense unless we say that we look for finitely supported solutions.
%But mentioning that will probably decrease the readability of the introduction.}
%
The function is invariant under permutations of atoms, and
%
its value is always greater than $2$.
Indeed, for every solution $\vr{x} : \A \to \R$ there is necessarily some $\b \in \A$ such that $\vr{x}(\b) > 0$,
and hence
\begin{align} \label{eq:gre1}
S(\vr x)  \ 
%\sum_{\a \in \A} \vr{x}(\a) =
%= \left(\sum_{\a \neq \b} \vr{x}(a)\right) + \vr{x}(b) 
> \ 2 \cdot \!\! \sum_{\a \in \A\setminus\set{\b}} \vr{x}(\a) \ \geq \ 2.
\end{align}
%
What is the minimal value of the objective function? 
For solutions $\vr x_ n$ defined in Example~\ref{ex:lp}, the value $S(\vr x_n) = \frac {2n} {n-1}$
may be arbitrarily close to 2 but, according to \eqref{eq:gre1}, $S$ never achieves 2. 
Surprisingly, this is in contrast with classical linear programming where, whenever constraints are specified by non-strict inequalities and are solvable,
%the solution set is closed and hence 
a linear objective function always achieves its minimum (or is unbounded from below).
%
%\arka{Also $+\infty$ if it is not solvable?}
%\piotrek{If it is not solvable then it is not solvable, you can say that in this case it is $+\infty$ but it is only a convention. I think.}
%\piotrek{I think the theorems says that the 
%image of a compact set is compact. If the set is closed and bounded then it is compact so the continuous function achieves its minimum. But the 
%compactness is crucial so only form the fact that it is closed we can not get that it achieves minimum.}
\end{slexample}

Inequalities and unknowns can be indexed by more than one atom, as illustrated in the next example.
%
\begin{slexample}
\label{ex:Kirchoff}
Let $\otu \A 2 = \setof{\a\b\in\A^2}{a\neq\b}$ be the set of pairs of distinct atoms.
Consider a system whose inequalities and unknowns are indexed by
%
%\begin{align*} %\label{eq:BCex}
$\A \uplus \set{*}$ and
%\qquad\text{and}\qquad
$\A  \uplus \otu \A 2$,
%\end{align*}
%
respectively.
Intuitively, unknowns correspond to vertices $\a$ and edges $\a\b$ of an infinite directed clique.
Let the system contain an inequality
%
\begin{align}  \label{eq:K1}
\sum_{\a \in \A} \a \ \geq \  1
\end{align}
%
enforcing the sum of values assigned to all vertices to be at least 1,
and the following inequalities
%
\begin{align}  \label{eq:K2}
\sum_{\b \in \A\setminus \set{\a}} \a\b \ - \ \a \ - \ \sum_{\b \in \A\setminus \set{\a}} \b\a  \ \geq \ 0 \qquad\quad (\a \in\A)
\end{align}
%
enforcing, for each vertex $\a\in\A$, the sum of values assigned to all outgoing edges
to be larger or equal to the sum of values assigned to all incoming edges, plus the value assigned
to the vertex $\a$.
In matrix form (0 entries of the matrix are omitted):

\vspace{-3mm}
\begin{align}
\label{eq:Kirchoff-matr}
\begin{aligned}
%\vr A \ = \ 
&\qquad\qquad\quad \begin{matrix} \A & \qquad\qquad \otu \A 2 \end{matrix}\\
%
&\begin{matrix}
\\
\A \\
\phantom{\ddots} \\
%\hline
\set{*}
\end{matrix}
%
\left[
\begin{matrix}
  -1 & &  \\  
  & -1 &     \\     
  &  & {\ddots} \ \\ 
  \hline
   \ \;1 & \ \; 1 & {\cdots}  \ 
\end{matrix}
\right\rvert
\hspace{-1.25mm}
\left\rvert
\begin{matrix}
\ \\
\qquad \vr A  \qquad \\
\phantom{\ddots} \\
\hline    
\ 
\end{matrix}
\right]
\cdot
\ \vr x 
\quad
\geq
\quad
\begin{bmatrix}
\ 0 \ \\ \vdots \\ 0 \\ \hline 1
\end{bmatrix}
\end{aligned}
\end{align}
\vspace{1mm}

\noindent
where $\vr A$ is the oriented incidence matrix, namely for every distinct atoms $\a,\b\in\A$, 
\[
\vr A(\a, \a \b) = 1 \qquad
\vr A(\a, \b \a) = -1,
\]
and all other entries of $\vr A$ are 0.
As in previous examples, the system is invariant under permutations of atoms.
Solutions of the system correspond to directed graphs, 
whose vertices and edges are labeled
in accordance with constraints \eqref{eq:K1} and \eqref{eq:K2}.
We return to this system in Example \ref{ex:Kirchoff-cont} below, and 
in Sections \ref{sec:weak} and \ref{sec:max}.
%\piotrek{We should write a sentence that later we will answer if it has a solution. If it is like this then the reader may be confused, if this is the end of the example.}
\end{slexample}

The systems appearing in the above examples are invariant under all permutations of atoms.
As usual when working in \emph{sets with atoms} \cite{atombook} (also known as \emph{nominal sets} \cite{Pitts:book}),
we allow systems which, for some finite subset $S\subseteqfin \A$ (called a \emph{support}), 
are only invariant under permutations that fix $S$.
In the standard terminology of orbit-finite sets \cite{atombook}, we allow for \emph{finitely supported}
systems.
Likewise, 
we allow for finitely supported objective functions, and seek for finitely supported solutions. 
Equivalently, using terminology of sets with atoms,
we allow for \emph{orbit-finite} systems and objective functions, and seek for
orbit-finite solutions.
Note that solutions appearing in Examples \ref{ex:lp} and \ref{ex:lpmin} are \emph{finitary}, 
i.e. assign non-zero to only finitely
many unknowns, and are therefore finitely supported.
Each finite system is finitely supported, and thus orbit-finite linear programming is a generalisation
of the classical one.


\para{Contribution}

As our main contribution, we provide decision procedures for orbit-finite linear programming, both for
the decision problem of solvability of systems of inequalities, and for the optimisation problem.

The core ingredient of our approach is to reduce solvability (resp.~optimisation) of an orbit-finite system of inequalities to the analogous question
on a finite system which is \emph{polynomially parametrised}, i.e., where coefficients are 
univariate polynomials in an integer variable $n$.
The parameter $n$ corresponds, intuitively, to the number of atoms appearing in (a support of) a solution.
In this parametrised setting we ask for solvability for \emph{some} $n\in\Nat$, or for optimisation
when ranging over \emph{all} $n\in\Nat$.
We can compute an answer by encoding the problem into first-order
real arithmetic~\cite{Tar51,ra-book}, %A Canny88,Ben-OrKR86,ra-book}, 
which yields decidability.
We provide also an efficient \PTIME algorithm for a relevant subclass of instances, 
resorting to polynomially many calls of 
classical linear programming.
%The algorithm is however completely different than the approach of \cite{GHL22}.

\begin{slexample} %\rm
\label{ex:poly}
%\piotrek{I think this example is pointless}
%, as the reader does not know what he/she should appreciate. It is like look here it is some magic, you can not really understand what is going on.} 
%\slawek{I would keep it, it illustrates polynomially-parametrised stuff}
For instance, the system \eqref{eq:matr} % in Example~\ref{ex:lp} 
is transformed to the following two inequalities with one unknown $x$,
which are polynomially (actually, linearly) parametrised in a parameter $n$ (the details are exposed in Example \ref{ex:cont} in Section \ref{sec:red2poly}): 
\begin{align} \label{eq:ppar}
\begin{bmatrix}
n-1 \\
n
\end{bmatrix}
\cdot
x
\ \geq \ 
\begin{bmatrix}
n \\
n
\end{bmatrix}
\end{align}
%
The objective function $S$~\eqref{eq:objf}, on the other hand, is transformed to a 
non-parametrised linear map $x\mapsto 2\cdot x$.
For every $n > 1$, the system \eqref{eq:ppar} is solvable,
the minimal solution is $x=\frac{n}{n-1}$, and
the minimum of the objective function is $\frac{2n}{n-1}$.
Ranging over all $n\in\Nat$, the minimum can get arbitrarily close to 2, but never reaches 2.
\end{slexample}

Our reduction to a polynomially-parametrised linear program involves a size blowup
which is exponential in \emph{atom dimension}
(i.e., the maximal number of atoms appearing in an index of inequality or unknown of a system)
but polynomial in the number of its orbits.
In consequence, orbit-finite linear programming is solvable in \EXPTIME, and in \PTIME when
atom dimension is fixed.
In the setting of orbit-finite sets this means that the problem is feasible \cite{BT18}.
Therefore, in the latter case the complexity is not worse than in case of
classical linear programming.

%\arka{Also cite Tarski?}

\smallskip

One of cornerstones of the reduction is an observation that every system of inequalities that admits
a finitary solution, 
admits also a solution which is invariant under all permutations of its support, i.e., of atoms it uses.
%\piotrek{Not clear what is uses means?}
%In short words, solvable invariant systems admit invariant solutions.

\begin{slexample}
\label{ex:Kirchoff-cont}
As an illustration, let's prove that the system in Example \ref{ex:Kirchoff} has no finitary solution.
Indeed, if there existed a finite labeled  directed graph 
satisfying constraints \eqref{eq:K1} and \eqref{eq:K2}, 
by the above observation there would also exist a finite labeled directed
clique, where labels of all vertices are pairwise equal, and labels of all edges are pairwise equal as well.
In particular each vertex would carry the same value, 
necessarily positive due to constraint \eqref{eq:K1},
and
all edges incoming to a vertex would carry the same value as all outgoing edges.
These requirements are clearly contradictory with constraint \eqref{eq:K2}.
In conclusion, the system has no solutions.
%\piotrek{I think it is not clear what kind of solutions we are looking for, so it can be not clear what we mean by saying that there are no solutions.}
\end{slexample}


As our second main result we prove undecidability of orbit-finite \emph{integer} linear programming,
already for the decision problem of solvability.
While the classical linear programming and integer linear programming are on the opposite sides of the feasibility border,
%Therefore the \PTIME vs \NP gap between classical linear and integer linear programming extends,
in case of orbit-finite systems the two problems are on the opposite sides of the decidability border.
One of key reasons behind undecidability of integer linear programming is that it can
express existence of a finite path. Specifically, if integer solutions are sought,
the system of inequalities of the form
\[
0 \ \leq \ \sum_{\a \in \A\setminus\set{\b}} \b \a \ \leq \ 1 \quad (\b \in \A)
\]
   allow us to say that every node $\beta$ has either no successors, or just one successor. 
   This clearly fails if real solutions are sought.
   
This article is an improved and extended version of the conference submission \cite{GHL23}.
The major improvement is lowering the complexity of orbit-finite linear programming from
\twoEXPTIME to \EXPTIME, and from \EXPTIME to \PTIME for fixed atom dimension.
This is achieved by identifying a suitable subclass of polynomially-parametrised 
linear programs which we show to be
solvable in \PTIME, and a reduction from orbit-finite linear programming to this subclass.



\para{Related research} 

This paper belongs to a wider research program that aims at lifting different aspects of 
theory of computation from finite to orbit-finite sets (essentially equivalent to first-order definable sets)
\cite{lics11,datamonoids,program-atoms,BKLT13,lmcs14,KLOT14,CL15,KKOT15,KLOT16,BT18,KKLO19}.

Our findings generalise, or are closely related to, some earlier % (in terms of decidability but not complexity) 
results %~\cite{HLT17,HR21,KKOT15,GHL22} 
on systems of linear equations.
%\piotrek{No these results ware not about inequalities but solving equations over $\Nat$ which in case of orbit finite systems is weaker.} 
%\arka{Should we add reference to paper by BKM, since they proved $\GLin{B}$ has orbit-finite spanning set} 
%
Systems studied in~\cite{HLT17} have row indexes of atom dimension 1.
In a more general but still restricted case studied in~\cite{HR21},
all row indexes are assumed to have the same atom dimension. 
%
%\piotrek{We ware able to consider only columns indexed by sets of atoms but not by a more complicated structures like sets of sets with atoms, which are used to encode tuples. Maybe this is a good way to classify these results, by the level of the hierarchy of sets with atoms.}
%\slawek{This is not a good classification criterion. Is my text ok or should I modify it?}
%\piotrek{Your text is ok. In my first comment I have written wrong. In~\cite{HLT17} rows were indexed with atoms so level 0. In~\cite{HR21} rows were indexed 
%with sets of size $n$, so level $1$ not full level $1$ because sizes of sets were equal. Now rows may be on arbitrary level of the hierarchy. Maybe there is 
%no point in emphasising this but it is what actually happened.}  
%
Furthermore, columns of a matrix are assumed to be finitary in~\cite{HLT17,HR21}.
Both the papers investigate (nonnegative) integer solvability and are subsumed by
\cite{GHL22} (in terms of decidability, but not in terms of complexity), 
a starting point for our investigations.
%
Orbit-finite systems of equations are a special case of our present setting, as long as
   the solution domain is reals or integers. However, the setting of \cite{GHL22} is not 
   restricted to reals or integers, and allows for an arbitrary commutative ring as a solution domain
   (under some mild effectiveness assumption). This larger generality makes the results of 
   \cite{GHL22} and our results incomparable. Consequently, our methods are different than those
   of \cite{GHL22}.
   
The work~\cite{HL18} goes beyond \cite{HLT17} and investigates linear equations, in atom dimension 1,
over ordered atoms. 
Nonnegative integer solvability is decidable and equivalent to VAS reachability  
(and hence \ackermann-complete~\cite{CO,L,Lasota22}).

%\TODO{Piotrek: check if the above short info is correct and sufficient?}
%\piotrek{Done}

Systems in another related work~\cite{KKOT15} are over a finite field, contain only 
finite equations, and are studied as a special case of orbit-finite constraint satisfaction problems.
Furthermore, solutions sought are not restricted to be finitely-supported.

Orbit-finitely generated vector spaces were recently investigated in \cite{BKM21} and \cite{GHL22}.
The former paper shows that every chain of vector subspaces which are invariant under permutations of atoms eventually stabilises,
and apply this observation to prove decidability of zero-ness for orbit-finite \emph{weighted} automata.
The two papers jointly show that dual of an orbit-finitely generated vector space has an orbit-finite base.
%
In \cite{MinkowskiWeyl2021} the authors study cones in such spaces which are invariant under permutations of atoms, 
and extend accordingly theorems of Carath\'eodory and Minkowski-Weyl.

Finally, our technique discussed in Example \ref{ex:Kirchoff-cont}, and
developed formally in Section \ref{sec:weak},
seems to be reminiscent of (but independent from) 
the techniques in the recent work \cite{Asia}.




%\arka{In (https://arxiv.org/pdf/2110.10657.pdf) the authors study equivariant cones inside $\flin{\A}$,
%where they extend theorems of Carath\'eodory and Minkowski-Weyl to this setting (theorem 3.1 and 4.4).}

\para{Outline} 
After preliminaries on orbit-finite sets in Section \ref{sec:atoms},
in Section \ref{sec: linear eq} we introduce the setting of orbit-finite linear inequalities and
in Section \ref{sec:results} we state our results.
The rest of the paper is devoted to proofs.
In Sections \ref{sec:weak} and \ref{sec:poly}  we develop tools that are later used in 
decision procedures for linear programming in Sections \ref{sec:lp-decid-proof} and \ref{sec:max}.
Finally, Section \ref{sec:ilp-undecid-proof} contains the proof of undecidability of integer linear programming.
We conclude in Section \ref{sec:conc}.
%\piotrek{Do we plan some conclusions/future work.}
Some routine or lengthy arguments are moved to Appendix. 

% !TEX root = lin-prog-main.tex

\section{Preliminaries on orbit-finite sets} \label{sec:atoms}

%\para{Sets with atoms} 
%
%\arka{I think there are several places where the mentioned fact is true for specific
%class of atoms including equality atoms, but not true for atoms in general.
%But I don't see anywhere where it's mentioned that we are working with equality atoms}
Our definitions rely on basic notions and results of the theory of \emph{sets with atoms}~\cite{atombook}, 
also known as nominal sets~\cite{Pitts:book,lmcs14}. 
%This paper is a 
%%In this section we recall what is necessary to follow our arguments; this is 
%part of the theory
%of computing in the realm of orbit-finite sets, developed e.g.~in~\cite{lics11,lmcs14,atombook}.
We only work with \emph{equality atoms} which have no additional structure except for
equality.

\para{Sets with atoms}

We fix a countably infinite set $\A$ %=\set{\at 1, \at 2, \at 3, \ldots}$, 
whose elements we call \emph{atoms}.
Greek letters $\a,\b,\g,  \ldots$ are reserved to range over atoms.
%
%Informally speaking, a set with atoms is a set that can have atoms, or other sets with atoms, as elements.
The universe of sets with atoms is defined formally by a suitably adapted cumulative hierarchy of sets,
by transfinite induction: % on ordinals $\gamma$:
the only set of \emph{rank} 0 is the empty set; and for a cardinal $i$, a set of rank $i$ may contain, as elements, 
sets of rank smaller than $i$ as well as atoms.
%\arka{Should we use some other letters since we reserved greek letters for atoms?}
In particular, nonempty subsets $X\subseteq\A$ have rank 1.

The group $\Aut{}$ of all permutations of $\A$, called in this paper \emph{atom automorphisms},
acts on sets with atoms by consistently renaming all atoms in a given set.
%\piotrek{in the set $\A$. I think, if we write in a given set then this may suggest some parameter like $\Aut{}(\A\setminus\{\alpha,\beta\})$}. 
Formally, by another transfinite induction, for $\pi\in\Aut{}$ we define $\pi(X) = \setof{\pi(x)}{x\in X}$.
Via standard set-theoretic encodings of pairs or finite sequences we obtain, in particular, 
the pointwise action on pairs $\pi (xy)=\pi(x)\pi(y)$,
and likewise on finite sequences.  %: $\pi\cdot (x_1, \ldots, x_n)=(\pi\cdot x_1, \ldots, \pi\cdot x_n)$ .
Relations and functions from $X$ to $Y$ are considered as subsets of $X\times Y$.
%for instance, in case of $f:\A\to\A$, we have $\pi(f)(\a)=\pi(f(\pi^{-1}(\a)))$, or
%$\pi(f)(\pi(\a))=\pi(f(\a))$.
%If a set $X$ is \emph{pure}, i.e., no atom belongs to $X$, its elements, elements of its elements, etc., 
%then $\pi(X) = X$ for every $\pi\in \Gr$; for instance $\pi(n) = n$ for $n\in\Nat$.

We restrict to sets with atoms $X$
% \piotrek{add:``$x\in X$'' otherwise one thinks about $S$ as about the set, ant it is confusing.} 
that only depend on finitely many atoms, in the following sense. 
For $T\subseteq \A$, let $\Aut T = \setof{\pi\in\Aut{}}{\pi(\a) = \a \text{ for every } \a \in T}$ be the set of atom automorphisms
that \emph{fix} $T$%
\footnote{
$\Aut T$ is often called the \emph{pointwise stabilizer} of $T$.
};
they are called $T$-automorphisms.
%, called \emph{\red{the pointwise stabilizer of $T$}}.
A finite set $T\subseteqfin\A$ (we use the symbol $\subseteqfin$ for finite subsets) is a \emph{support} of $X$ if for all
$\pi\in\Aut T$ it holds $\pi(X)=X$.
We also say: $T$ \emph{supports} $X$, or $X$ is \emph{$T$-supported}.
%\end{quote}
Thus a set is $T$-supported if and only if it is invariant under all $\pi \in \Aut T$.
% \red{automorphisms in the pointwise stabilizer of $T$}.
As an important special case, a function $f:X\to Y$, 
understood as its diagram $\setof{(x,f(x))}{x\in X}$, 
is $T$-supported if  $f(\pi(x)) = \pi(f(x))$ for every argument $x$ and $\pi\in\Aut{T}$.
In particular, whenever $f$ is $T$-supported, its domain $X$ is necessarily $T$-supported too.
A $T$-supported set is also $T'$-supported, assuming $T\subseteq T'$.

A set $x$ is \emph{finitely supported} if it has some finite support; in this case
%Finite supports of a fixed element are closed under intersections, and hence every finitely supported 
$x$ always has the least (inclusion-wise) support, denoted $\supp x$, called \emph{the support} of $x$ 
(cf.~\cite[Sect.~6]{atombook}).
%(c.f.~\cite[Prop.~2.3]{Pitts:book}, \cite[Cor.~9.4]{lmcs14}).
Thus $x$ is $T$-supported if and only if $\supp x \subseteq T$.
Sets supported by $\emptyset$ (i.e., invariant under all atom automorphisms) we call \emph{equivariant}.

\begin{slexample}
Given $\a,\b\in\A$, the support of the set $\A\setminus\set{\a,\b}$ is $\set{\a,\b}$.
The set $\A^2$ and the projection function $\pi_1 : \A^2 \to \A : (\a, \b) \mapsto \a$
are both equivariant; 
% (indeed, every $\pi \in \Gr$ preserves the graph $\setof{\pair{\pair a b} a}{a,b\in\A}$ of $\pi_1$);
and the support of a tuple $\tuple{\a_1, \ldots, \a_n}\in\A^n$, encoded as a set in a standard way, is
the set of atoms $\set{\a_1, \ldots, \a_n}$ appearing in it.
\end{slexample}

From now on, we shall only consider sets  that are hereditarily finitely
supported, %(called briefly \emph{legal}), 
i.e., ones that have a finite support, whose every element has some finite support,
and so on.
%A set $X$ in the cumulative hierarchy is called \emph{legal} (or \emph{set with atoms})
%if it is hereditarily finitely-supported: it is finitely-supported itself, all its elements are finitely supported, etc.
%From now on all sets are silently assumed to be legal.

%\vspace{-0.8mm}

\para{Orbit-finite sets} 

Let $T\subseteqfin\A$. Two atoms or sets $x, y$ % of an $S$-supported set  
are \emph{in the same $T$-orbit} if $\pi(x) = y$ for some $\pi\in\Aut T$.
This equivalence relation splits atoms and sets 
%every $S$-supported set $X$ 
into equivalence classes, which we call \emph{$T$-orbits};
$\emptyset$-orbits we call equivariant orbits, or simply \emph{orbits}.
%\slawek{these 2 sentences added:}
By the very definition, every $T$-orbit $\O$ is $T$-supported: $\supp \O \subseteq T$.%
\footnote{The inclusion may be strict, for singleton $T$-orbits $O$.
For instance, the singleton $\set{\a}\subseteq\A$ is a $\set{\a}$-orbit, but also a $\set{\a,\b}$-orbit for $\b\neq \a$.
}


%and $O$ is also a $\supp O$-orbit.
%When the set $S$ is irrelevant, we simply speak of an \emph{orbit}, 
%\slawek{this is not a good idea, 'orbit' should mean 'equivariant orbit'} having in mind an $S$-orbit for some $S\subseteqfin \A$.
%
$T$-supported sets are exactly unions of (necessarily disjoint) $T$-orbits.
Finite unions of $T$-orbits, for any $T\subseteqfin\A$, are called  \emph{orbit-finite} sets.
Orbit-finiteness is stable under orbit-refinement: if $T\subseteq T' \subseteqfin \A$, a finite union of $T$-orbits is also a finite union
of $T'$-orbits (but the number of orbits may increase, cf.~\cite[Theorem~3.16]{atombook}). 

\begin{slexample}
Examples of orbit-finite sets are: 
\begin{itemize}
\item the set of atoms $\A$ (1 orbit); 
\item $\A \setminus \set{\a}$ for some $\a\in\A$ (1 $\set{\a}$-orbit);
\item pairs of atoms $\A^2$ (2 orbits: diagonal $\setof{\a\a}{\a\in\A}$ and off-diagonal $\otu \A 2 = \setof{\a\b\in\A^2}{\a\neq \b}$);
\item $n$-tuples of atoms $\A^n$ for $n\in \Nat$; each orbit $U\subseteq \A^n$ contains all $n$-tuples of the same equality type,
where by the \emph{equality type} of an $n$-tuple $a_1 \ldots a_n\in\A^n$ we mean
the set $\setof{(i,j)}{a_i = a_j}$;
%\item
%$\A^3$ (5 orbits, corresponding to equality types of triples, defined below;
%\piotrek{It is 1/3 page later, some reference would be good});
%%atom permutations fix elements of $H$);
\item non-repeating $n$-tuples of atoms 
$
\otu \A n = \setof{\!\a_1\ldots\a_n\in\A^n\!}{\a_i\neq\a_j \text{ for }i\neq j}
$  (1 orbit);
\item $n$-sets of atoms $\utu \A n = \setof{X\subseteq \A}{\size X = n}$ (1 orbit).
%
\end{itemize}
All of them are equivariant, except $\A\setminus\set{\a}$.
On the other hand,
the set  $\pow {\text{fin}} \A$ of all finite subsets of atoms is orbit-infinite as cardinality is an invariant of each orbit.
\end{slexample}

We now state few properties to be used in the sequel.
For $T\subseteqfin\A$,
each $T$-orbit $\O\subseteq \otu \A n$ is determined by fixing pairwise distinct atoms
from $T$ on a subset $I\subseteq \setto n$ of positions, while allowing arbitrary atoms from $\A\setminus T$ on
remaining positions $\setto n\setminus I$:
%
\begin{lemma} \label{lem:tupleorbit}
Let $T\subseteqfin \A$.
$T$-orbits $\O\subseteq \otu \A n$ are exactly sets of the form
%
\begin{align} \label{eq:Torb}
\begin{aligned}
\setof{a \in \otu \A n}{\proj n I(a) = u, \ \ \proj n {\setto n \setminus I} (a) \in \otu {(\A\setminus T)} {n-\ell}},
\end{aligned}
\end{align}
%
where $I \subseteq \setto n$, $\size I = \ell$, and $u\in\otu T \ell$.
The projection $\proj n I : \otu \A n \to \otu \A \ell$ is defined in the expected way.
\end{lemma}
%
%\noindent
%\begin{proof}
Indeed, the set \eqref{eq:Torb} is invariant under all $T$-automorphisms,
%\red{automorphism in the pointwise stabilizer of $T$}, 
and each two of its elements
are related by some $T$-automorphism.
% \red{automorphism in the pointwise stabilizer of $T$}.
%\end{proof}
%\piotrek{Missing proof.}
%
%
%\smallskip
%
The orbit \eqref{eq:Torb} is in $T$-supported bijection with $\otu {(\A\setminus T)} {n-\ell}$.
This is a special case of a general property of
every $T$-orbit, not necessarily included in $\otu \A n$.
The following lemma is proved exactly as \cite[Theorem 6.3]{atombook}
and provides finite representations of $T$-orbits:
%
\begin{lemma}\label{lem:repn}
Let $T \subseteqfin \A$.
Every $T$-orbit admits a $T$-supported bijection to a set of the form 
$
\otu {(\A \setminus T)} n/_G,
$
for some $n \in \Nat$ and some subgroup $G$ of $S_n$.
\end{lemma}
%

%\smallskip

%By the \emph{equality type} of an $n$-tuple $a_1 \ldots a_n\in\A^n$ we mean the set $\setof{(i,j)}{a_i = a_j}$.
Recall that each orbit $U\subseteq \A^n$ contains all $n$-tuples of the same equality type.
In particular,  each orbit included in $\otu \A n \times \otu \A m \subseteq \A^{n+m}$
is induced by a partial injection $\iota$ from $\setto n$ to $\setto m$:
%
\begin{lemma} \label{lem:orbitpair}
Orbits $\O\subseteq \otu \A n \times \otu \A m$
are exactly sets of the form
$
\setof{ (a, b) \in \otu \A n \times \otu \A m
}{\prettyforall{i,j}{a(i) = b(j) \iff \iota(i) = j}},
$
where $\iota$ is a partial injection from $\setto n$ to $\setto m$.
\end{lemma}
%\piotrek{If we do not add the proof in the appendix, or state that it is trivial. Lemmas without proofs are evil.}
%
%\slawek{Lemma \ref{lem:tupleorbit} to app? Notation used in Sec \ref{sec:weak}}
%\arka{I think it is good to have the lemma here}
%

Atom automorphisms preserve the size of the support: $\size{\supp X} = \size{\supp{\pi(X)}}$ for every set $X$ and
$\pi\in\Aut{}$.
We define \emph{atom dimension} of an orbit as the size of the support of its elements.
For instance, atom dimension of $\otu \A n$ is $n$.

%\slawek{introduce atom dimension}



% !TEX root = lin-prog-main.tex


\section{Orbit-finite (integer) linear programming} 
\label{sec:problems}
%\section{Orbit-finite systems of linear inequalities}
\label{sec: linear eq}



%For the sake of concreteness we work in this paper with vector spaces over the rational field $\R$.

%Furthermore, since we never use division,
%the results are still valid when a field is replaced by a commutative ring,
%e.g.~the integer ring $\Int$, as long as the above effectiveness assumptions are satisfied.
%\slawek{is this remark fine?}
%\slawek{since we can use a ring, do we need to keep the subgroup $\X$ as a parameter?}
%\arka{I don't see why not}\PH{the problem is that in our assumptions we have effective inverse, which is false for rings. It has to be changed to we can solve 
%equations over a ring.} 

%\para{Vectors}

We introduce now the setting of linear inequalities we work with, and formulate our main results. 
%
We are working in vector spaces over the real%
\footnote{
All the results of the paper still hold if reals $\R$ are replaced by rationals $\Rat$
in all the subsequent definitions and results.
}
field $\R$,
%\footnote{
%Formally, in case when $\R$ is not a field, we should use the term \emph{module}.
%Since modules/vector spaces studied in this paper are of particularly simple kind,
%we prefer to stick to a widely known term \emph{vector space}.
%} 
where vectors are indexed by a fixed orbit-finite set $B$, i.e., are functions $\vr v : B\to\R$.
Observe that such a function $\vr v$, understood as its diagram $\setof{(b, \vr v(b))}{b\in B}$,
is orbit-finite exactly when it is finitely supported
(according to definitions in Section \ref{sec:atoms}).

\begin{definition}
By a \emph{vector} over $B$ we mean any orbit-finite (i.e., finitely-supported) function
%\footnote{
%That is, a function which is finitely-supported as a set of pairs (according to definitions in Section \ref{sec:atoms}), not necessarily a function  
%which is almost everywhere 0.} 
$\vr v$ from $B$ to $\R$, 
written $\vr v : B\tofs\R$ (vectors are written using boldface). 
%\red{Note that by finitely-supported function we mean a finitely supported set of pairs, not 
%function that is equal $0$ for all but finitely many elements.}
Vectors with integer range, $\vr v : B \tofs \Int$, we call \emph{integer} vectors.
\end{definition}

The set of all vectors over $B$ we denote by $\GLin B = B \tofs \R$. 
It is a vector space, 
with pointwise addition and scalar multiplication:
for $\vr v, \vr v'\in\GLin B$, $b\in B$ and $q\in \R$, we have
$(\vr v+ \vr v')(b) = \vr v(b) + \vr v'(b)$ and $(q \cdot \vr v)(b) = q\cdot \vr v(b)$.
These operation preserve the property of being finitely-supported, e.g., 
$\supp{\vr{v} + \vr{v'}} \subseteq \supp{\vr{v}} \cup \supp{\vr{v}'}$.
%
%The generating set $B$ can be considered as \emph{the dimension} of $\GLin B$.
%\TODO{do we want to call B dimension?}
%$\GLin B$ may be considered as  the vector space \emph{generated} by $B$.
%and $B$ as its \emph{dimension}%
%\footnote{
%Not to be confused with \emph{atom dimension} introduced in Section~\ref{sec:atoms}.
%}.
%
%\arka{Probably better to not say $B$ to be the dimension of $\lin{B}$}
%
We define the \emph{domain} of a vector $\vr v\in\GLin B$ as $\dom {\vr v} = \setof{b\in B}{\vr v(b)\neq 0}$.
%
%\para{Finitary vectors}
%
A vector $\vr v$ over $B$ is \emph{finitary}, written $\vr v : B\tofin \R$, 
if $\dom {\vr v}$ is finite, i.e., $\vr v(b) = 0$ for almost all  $b\in B$.

\begin{slexample}
\label{ex:vect}
Let $B = \otu \A 2$. Let $\a, \b\in \A$ be two fixed atoms.
The function $\vr v : B \to \R$ defined, for  $\x,\g\in\A\setminus\set{\a,\b}$, by
\begin{align*}
\vr v(\a \x) &= \vr v(\x \a) = -1 &
\vr v(\a \b) &= \vr v(\b \a) = 3  \\ 
\vr v(\b \x) &= \vr v(\x \b) = -2
&
\vr v(\x \g) &= 0
\end{align*}
is an $\set{\a,\b}$-supported integer vector over $B$.
It is not finitary, as $\dom {\vr v} = \setof{\d\s\in\otu \A 2}{\set{\d,\s}\cap\set{\a,\b} \neq \emptyset}$ is infinite.
%\piotrek{can we use different names for atoms $\x\g$? It is confusing if once they are in $\A\setminus\{\a,\b\}$ and next time they are any elements of $\A$.} .
Finitary $\set{\a,\b}$-supported vectors over $B$ assign $0$ to all elements of $B$ except for $\a\b$ and $\b\a$.
\end{slexample}


A finitary vector $\vr v$ with domain $\dom {\vr v} = \set{b_1, \ldots, b_k}$ such that
$\vr v(b_1) = q_1, \ldots, \vr v(b_k) = q_k$, 
may be identified with a formal linear combination of elements 
of $B$:
%
\begin{align} \label{eq:formalsum}
\vr v \ = \ q_1 \cdot b_1 + \ldots + q_k \cdot b_k.  %, \qquad  q_1,\ldots,q_k\in\R\setminus\set{0}.
\end{align}
%
The subspace of $\GLin B$ consisting of all finitary vectors we denote by $\Lin B = B\tofin\R$.
%The space $\GLin B$ can be considered as the dual space of $\Lin B$\PH{Why?, is it a definition or a lemma?}, since we
%restrict our considerations to finitely-supported vectors only.
For finite $B$ of size $\size B = n$, $\GLin B = \Lin B$ is isomorphic to $\R^{n}$.
%
%\toremove{
%\begin{slexample}
%Let $B = \otu \A 2$.
%The domain $\dom v$ of a finitary vector $v\in \Lin B$
%can be seen as edges of a directed graph without self-loops whose nodes are atoms, and the vector
%$v$ is just a $\R$-labelling of edges of this graph.
%%
%%\TODO{finish}
%%\PH{What do we want to present in this example?}
%\end{slexample}
%}
%\para{Characteristic vectors}

For a subset $X\subseteq B$, we denote by $\constvr 1 X\in\GLin B$ the characteristic function of $X$,
i.e., the vector that maps each element of  $X$ to $1$ and all elements of $B\setminus X$ to $0$:
\[
\constvr 1 X : b \mapsto \begin{cases}
1 & \text{ if } b \in X \\
0 & \text{ otherwise.}
\end{cases} 
\]
We write $\constvr 1 b$ instead of $\constvr 1 {\set{b}}$, and $\vr 1$ instead of $\constvr 1 B$.

%We sometimes want to treat $B$ itself as a subset of $\Lin B$, identifying every $b\in B$ with
%the vector $\constvr 1 b$, or equivalently with
%the trivial linear combination $1\cdot b$ as in~\eqref{eq:formalsum}.

%Let $\emb B = \setof{\constvr 1 b}{b\in B} \subseteq \Lin B$.
%Note that the mapping
%$
%b  \mapsto  \constvr 1 b : B \to \emb B
%$
%allows us to treat elements of $B$ as vectors in $\Lin B$.
%

%\para{Finitary spans}

\begin{lemma} \label{lem:1O}
Let $T\subseteqfin \A$ and $\vr v \in \GLin {B}$ such that $\supp {\vr v}\subseteq T$. Then
\begin{enumerate}
\item[(i)] $\vr v$ is constant, when restricted to every $T$-orbit $\O\subseteq B$;
\item[(ii)] $\vr v$ is a linear combination of characteristic vectors $\constvr 1 \O$ of $T$-orbits $\O\subseteq B$.
\end{enumerate}
\end{lemma}
%
\begin{proof}
%
%Consider any vector $\vr v \in \GLin {B}$, and let $S = \supp{\vr v}$.
The first part follows immediately as $T$ supports ${\vr v}$.
%This allows us to write $\vr v(\O) \in \R$ in place of $\vr v(x)$ for  $x \in \O$. 
As required in the second part, we have:
\begin{align} \label{eq:1O}
\vr v = \sum_\O \vr v(b_\O) \cdot \constvr 1 \O,
\end{align}
where $\O$ ranges over finitely many %pairwise-disjoint 
$T$-orbits $\O\subseteq B$, and 
$b_{\O}\in\O$ are arbitrarily chosen representatives of $T$-orbits. 
%Therefore $\vr v\in  \Span {} {\setof{\constvr 1 O}{O\in\orbits B}}$, as required.
\end{proof}
%
%\begin{corollary} \label{cor:1O}
%$\GLin B = \Span {} {\setof{\constvr 1 O}{O\subseteq B \text{ an orbit 
%\PH{This notation look strange, maybe write the sentence $O$ are a finitely supported 
%orbits of $B$}}}}$.
%\end{corollary}
%
\begin{notation} \label{not:orbval}
In the sequel, whenever we know that a vector $\vr v : B\tofs \R$ is constant over a $T$-orbit
$\O \subseteq B$,
we may write $\orbval{\vr v}(\O)$ instead of $\vr v(b)$, where $b\in\O$.
In particular, when $\vr v$ is equivariant, we have the \emph{orbit-value} vector
\[
\orbval{\vr v} : \orbits B \to \R,
\]
where $\orbits B$ stands for the set of all equivariant orbits $\O$ included in $B$.
\end{notation}

%\para{Inner product}

We note that the inner product of  vectors $\vr x, \vr y\in\GLin B$, 
\[
\innerprod {\vr x} {\vr y} \ = \ \sum_{b\in B} \vr x(b)\, \vr y(b),
\]
is not always well-defined.
We consider the right-hand side sum as well-defined when there are only finitely many $b\in B$
for which both $\vr x(b)$ and $\vr y(b)$ are non-zero
%, i.e., finitely many $b \in B$ such that $\vr x(b)$ and $\vr y(b)$ are both non-zero
(equivalently, the intersection $\dom {\vr x} \cap \dom {\vr y}$ is finite).%
\footnote{In particular, $\innerprod {\vr x} {\vr y}$ is always well-defined when one of $\vr x, \vr y$ is finitary.}

%\arka{Maybe (in footnotes?) we can mention that the fact that $\vr{x}$ and $\vr{y}$ are finitely-supported implies this.
%As we did in our LICS paper.}
%

%\begin{slremark}
%%Consider $\R = \Rat$. % or $\R = \Int$
%%(the remark remains valid for some other rings $\R$ without nonzero zero-divisors). % \slawek{added}
%%\arka{I think these rings are called \emph{integral domains}. Probably we should mention it in the abstract as well}
%%\slawek{no, as this assumption matters only in this remark}
%Since vectors are finitely supported and hence (c.f.~Lemma \ref{lem:1O}) contain only finitely many different numbers, 
%$\dom {\vr x} \cap \dom {\vr y}$ is finite exactly when 
%%\slawek{should we prove this?}\PH{I think it is trivial, we don't need to}
%%\arka{I don't think that is necessary}
%the right-hand side sum is \emph{unconditionally convergent},
%i.e., convergent to the same value irrespectively of the order in which
%the elements $b\in B$ are enumerated%
%\footnote{We are grateful to Szymon Toru{\'n}czyk for attracting our attention to unconditional convergence.}.
%%\arka{\emph{convergence} is not an well defined concept for rings in general}
%%\slawek{right - I've adapted the remark accordingly}
%\slawek{do we remove this remark?}
%\end{slremark}

\para{Orbit-finite systems of linear inequalities}

Fix an orbit-finite set $C$ (it can be thought of as the set of unknowns).
%
By a linear inequality over $C$ we mean a pair $e = (\vr a, t)$ where 
$\vr a : C \tofs \Int$ is an integer vector of left-hand side coefficients and $t\in\Z$ is a right-hand side target value%
\footnote{Rational coefficients and target are easily scaled up to integers.}.
An \emph{$\R$-solution} (real solution) of $e$ is any vector $\vr x : C\tofs \R$ such that the inner product
$\innerprod {\vr a} {\vr x}$ is well-defined and
\[
\innerprod {\vr a} {\vr x} \geq t;
\]
$\vr x$ is an \emph{$\Int$-solution} (integer solution) if $\vr x: C\tofs \Int$.
We may also consider constrained solutions, e.g., 
%\emph{integer} ones (where $\vr x(c)\in\Int$ for all $c\in C$),
%\emph{nonnegative} ones (where $\vr x(c)\geq 0$ for all $c\in C$), or 
\emph{finitary} ones.
A linear \emph{equality} $\innerprod {\vr a} {\vr x} = t$ may be encoded by two opposite inequalities:
\[
\innerprod {\vr a} {\vr x} \geq t \qquad
- \innerprod {\vr a} {\vr x} \geq -t.
\] 
%\slawek{should we change the inequalities to the opposite one $A\cdot x \leq b$ in the whole paper from now on?}
%\arka{ok. We should also swap max and min then}

In this paper we investigate sets of inequalities indexed by an orbit-finite set.
Formally, an orbit-finite system of linear inequalities (over $C$)
is the pair $(\vr A, \vr t)$, where $\vr A : B \times C \tofs \Z$ is an integer \emph{matrix} 
with row index $B$ and column index $C$, and $\vr t : B \tofs \Z$ is an integer
\emph{target} vector:
%
%\vspace{-4mm}
\begin{align*}
& \quad\ \ \ \,  \begin{matrix} \ \cdots \quad &  c  & \quad \ \  \cdots\ \end{matrix} \\
& \begin{matrix} \vdots \\ b \\ \vdots \end{matrix} \ \ 
\begin{bmatrix}
\  & \vdots & \ \\
\ \cdots & \vr A(b,c) & \cdots\ \  \\
      & \vdots & \  
\end{bmatrix}
\qquad
\begin{bmatrix}
\vdots  \\
\vr t(b)\\
\vdots   
\end{bmatrix}
\end{align*}
%
%\smallskip

\noindent
For $b\in B$ we denote by $\vr A(b, \_) \in \GLin C$ the corresponding (row) vector.
%and symmetrically, for $c\in C$ we denote by $\vr A(\_, c) \in \GLin B$ the corresponding (column) vector.
A solution of a system $(\vr A, \vr t)$ is any vector $\vr x\in\GLin C$ which is a solution of all inequalities 
$(\vr A(b, \_), \vr t(b))$, $b\in B$.
%
Equivalently, $\vr x$ is a solution if $\mult{\vr A}{\vr x} \geq \vr t$, where $\geq$ is the pointwise order on vectors,
and the (partial) operation of multiplication of a matrix $\vr A$ by a vector $\vr x$ is defined in an expected way:
\[
(\mult {\vr A}{\vr x})(b) = \innerprod {\vr A(b, \_)} {\vr x}
\]
for every $b\in B$.
The result $\mult {\vr A} {\vr x} \in \GLin B$ is well-defined if $\innerprod {\vr A(b, \_)}{\vr x}$ is well-defined for all $b\in B$.


%\smallskip
By the following examples, restricting to equivariant, finitary or integer solutions only
has impact on solvability:

\begin{slexample}\label{ex:nofin1}
Let columns be indexed by $C = \A$,
and consider the system consisting of just one infinitary inequality
$(\constvr 1 {\A}, 1)$ ($B$ is thus a singleton).
Identifying column indexes $\a \in \A$ with unknowns, the inequality may be written as:
\[
\sum_{\a\in\A} \! \a \ \geq \ 1.
\]
The inequality has an integer (finitary) solution, i.e., $\vr x = \constvr 1 \a$ for any $\a\in\A$, but no 
equivariant one. 
Indeed, equivariant vectors $\vr x : \A \tofs \R$ are necessarily 
constant ones $\vr x = r\cdot \constvr 1 {\A}$ (cf.~Lemma~\ref{lem:1O}), and then
the inner product
\[
\innerprod {\constvr 1 {\A}} {\vr x} 
= \sum_{\a \in \A} \vr x(\a)
= \sum_{\a \in \A} r
\]
is well-defined only if $r=0$, i.e. $\vr x(\a) = 0$ for all $\a \in \A$.
%
\end{slexample}

\begin{slexample}
Let columns be indexed by $C = \otu \A 2$ and rows by $B = \utu \A 2$.
Consider the system containing, 
for every $\set{\a,\b}\in B$, the inequality
$(\constvr 1 {\a\b} + \constvr 1 {\b\a}, 1)$.
Using the formal-sum notation as in~\eqref{eq:formalsum} it may be written as
$(\a\b + \b\a, 1)$ or,
% (we use the formal-sum notation as in~\eqref{eq:formalsum}).
identifying column indexes $\a \b\in C$ with unknowns, as: % $\vr x(\a \b)$, as:
\[
\a\b + \b\a \ \geq \ 1 \qquad\qquad (\a,\b\in\A, \a \neq \b).
\]
%\piotrek{I think we should write example with equality instead of $\geq$, splitting it into $\geq$ and $\leq$ part makes is pointless.}
All the equations are thus finitary, and the target is $\vr t = \constvr 1 B$.
The constant vector $\vr x = \constvr {\frac 1 2} {} : \a\b \mapsto \frac 1 2$ is a solution,
even if we extend the system with symmetric inequalities 
\[
\a\b + \b\a \ \leq \ 1 \qquad\qquad (\a,\b\in\A, \a \neq \b).
\]
The extended system has no finitary solution.
%as such a solution is in contradiction with the infinitary target $\vr t = \constvr 1 B$.
It has no integer % (infinitary) 
solution either.
%\arka{maybe here we can also say that the assumption that solutions are finitely supported is important?}
Indeed, since we restrict to finitely supported solutions only, any such solution $\vr x$ 
necessarily satisfies, for every distinct atoms $\a,\b \in \A\setminus \supp{\vr x}$, 
the equality $\vr x(\a\b) = \vr x(\b\a)$, 
which is incompatible with $\vr x(\a\b) + \vr x(\b\a) = 1$.
%
\end{slexample}

%\begin{slexample} \label{ex:nofin}
%For $C = \otu \A 2$ and $B = \A$, consider the system of inequalities containing, for every $\a\in \A$, the equation
%$(\constvr 1 {\a\_}, 1)$.
%Identifying column indexes $\a \b\in C$ with unknowns, the system may be written as:
%\[
%\sum_{\b\in\A\setminus\set{\a}} \! \a\b \ \geq \ 1 \qquad (\a\in \A).
%\]
%The system is thus equivariant, and all its inequalities are infinitary.
%The system has an integer solution. 
%Take any two fixed atoms $\c,\d\in\A$ and consider the vector
%\[
%\vr x \ = \ \constvr 1 {\_ \c} + \constvr 1 {\c \d}.
%\]
%Indeed, for $\a\neq \c$ we have
%$\innerprod {\constvr 1 {\a\_}}  {\vr x} = \constvr 1 {\a\c} \cdot \constvr 1 {\a\c} = 1$ as required.
%Furthermore, for $\a=\c$ we have
%$\innerprod {\constvr 1 {\a\_}}  {\vr x} = \constvr 1 {\c\d} \cdot \constvr 1 {\c\d} = 1$ as required.
%The system has no finitary solution (essentially for the same reason as in the previous example), and no equivariant one\piotrek{Why do we speak about 
%equivariant solutions? Does it illustrate something?}
%(as the only equivariant vectors over $C$ are constant ones $q\cdot\constvr 1 {}$, and the inner product
%$\constvr 1 {\a\_} \cdot \constvr 1 {}$ is ill-defined for every $\a\in\A$ as long as $q\neq 0$).
%\end{slexample}


%\PHr{Due to orbit-finiteness of $B$ and $C$, both $\vr A$ and $\vr t$ contain only finitely many distinct numbers/coefficients
%(cf. Lemma~\ref{lem:1O}).
%By scaling one may thus obtain an equivalent system with integer coefficients.}{This fragment seems to be not connected with other part of the text and is 
%distracting.}


%
%\para{Injective column-indexed representation}
%
%One can effectively translate between row-indexed, bi-indexed and column-indexed representations of the matrix $\vr A$.
%%
%\begin{lemma} \label{lem:set}
%In case of the column-indexed representation $\vr A : C\to \GLin B$ one
%can assume w.l.o.g.~that the mapping $\vr A : C\to \GLin B$ is injective.
%\end{lemma}
%%
%Intuitively speaking, one can assume that column vectors are pairwise different.
%This allows us to identify a matrix with the (non-indexed) set of its column vectors.
%%
%\begin{proof}
%Let $C' = C/\sim$ be the the quotient of $C$ by the kernel $\sim$ of the mapping:
%\[
%c \sim c' \iff \vr A(c) = \vr A(c'),
%\]
%and let $\vr A' : C'\to\GLin B$ be the 
%column-indexed representation of the induced matrix with columns $C'$. 
%Both $C'$ and $\vr A'$ can be effectively computed from $C$ and $\vr A$.
%Lemma~\ref{lem:set} follows once we show:
%%
%\begin{claim}
%$\GSpan {}{\vr A} = \GSpan {}{\vr A'}$.
%\end{claim}
%%
%Fir simplicity of presentation we prove the claim for equivariant $C$.
%In one direction, suppose $\mult {\vr A} {\vr x} = \vr t$ is well-defined for some $\vr x\in\GLin C$ and consider any
%equivalence class $E\subseteq C$ of $\sim$.
%Let $\vr A(c) = \vr z$ for all $c\in E$. 
%If $\vr z = \zerovec$ we can assume w.l.o.g.~that $\vr x(c) = 0$.
%We observe that  $\vr x$ is necessarily finitary when projected to $E$,
%$\restr {\vr x} E \in \Lin E$, as otherwise $\mult {\vr A} {\vr x}$ would be ill-defined on coordinate $b\in B$ for which $\vr z(b)\neq 0$.
%Therefore, the vector $\vr x'\in\GLin{C'}$ defined by $\vr x'(E) = \sum_{c\in E} x(c)$ satisfies
%$\mult {\vr A'} {\vr x'} = \vr t$.
%
%In the reverse direction, 
%suppose $\mult {\vr A'} {\vr x'} = \vr t$ is well-defined for some $\vr x'\in\GLin {C'}$.
%The set $C/\sim$ of equivalence classes of $\sim$ is equivariant and orbit-finite.
%Let $\cal O$ be the (finite) set of equivariant orbits in $C/\sim$.
%For each orbit $O\in\cal O$  choose one equivalence class $E_O \in O$ and one element $c_O\in E_O$.
%Thus $\setof{c_O}{O\in {\cal O}}$ is a finite subset of $C$, its size is the same as the size of $\cal O$.
%
%Fix a sufficiently large $S\subseteqfin\A$.
%For every equivalence class $E\in C/\sim$ consider the finite subset $\overline E \subseteq E$ of elements supported by
%$\supp{E} \cup S$.
%We claim that the union $\overline C = \bigcup_{E\in C/\sim} \overline E$ is supported by $S$.
%
%The finite set of pairs $\setof{(c_O, E_O)}{O\in\cal O}$ extends to an equivariant bijection $\iota$ between 
%$\overline C = \bigcup_{O\in {\cal O}} \orbit \emptyset {c_O}$ and $\cal E$, given by the membership relation.
%The vector $\vr x \in \GLin C$ defined by
%\[
%\vr x(c) = \begin{cases}
%\vr x'(\iota(c)) & \text{ if }c \in \overline C\\
%0 & \text{ otherwise.}
%\end{cases}
%\]
%satisfies $\mult {\vr A}{\vr x} = \vr t$.
%\end{proof}



% !TEX root = lin-prog-main.tex


%\section{Orbit-finite (integer) linear programming} \label{sec:problems}

\section{Results} \label{sec:results}

\para{Solvability problems}

We investigate decision problems of solvability of orbit-finite systems of inequalities over the ring of reals
%\footnote{Equivalently, one could consider rationals instead of reals.} 
or integers.
Consequently, we use $\F$ to stand either for $\R$ or $\Int$.
%
We identify a couple of variants. 
In the first one we ask about existence of a finitely-supported solution:

\prob{\ineqsolvname$(\F)$}
{An orbit-finite system of linear inequalities}
{Does it have an \xspace $\F$-solution}
%

\noindent
We recall that solutions are finitely-supported (equivalently, orbit-finite), by definition.
A closely related variant is solvability of equalities, under the restriction to nonnegative solutions only:

\prob{\nonnegsolvname$(\F)$}
{An orbit-finite system of linear equations}
{Does it have a nonnegative \xspace $\F$-solution}
%
\noindent
Furthermore, both the problems have "finitary" versions, where one seeks for finitary solutions only,
denoted as \finineqsolvname$(\F)$ and \finnonnegsolvname$(\F)$, respectively.

Three out of the four problems are inter-reducible, and hence equi-decidable,
both for $\F = \R$ and $\F = \Int$:
%(the proof is in Appendix):
%
\begin{theorem} \label{thm:equiv-problems}
Let $\F\in\set{\R,\Int}$. 
The problems \ineqsolvname$(\F)$, \finineqsolvname$(\F)$ and \nonnegsolvname$(\F)$ are inter-reducible. 
All reductions are in \PTIME, except the one from \ineqsolvname$(\F)$ to \finineqsolvname$(\F)$ 
which is in \EXPTIME, but in \PTIME for fixed atom dimension.
%\piotrek{complexity of reductions, are they polynomial? or linear?}
\end{theorem}
%
%\noindent
%The proofs of Theorem~\ref{thm:equiv-problems}, to be found in Appendix,
%amounts to lifting of standard arguments from finite to orbit-finite systems, 
%and checking that all constructed objects are finitely supported. 
%\piotrek{It can not be like this as fin-solv vs solve is not present in standard case.}
%
(The proof is in Section \ref{sec:problems-proofs-one}.)
The three problems listed in Theorem~\ref{thm:equiv-problems} deserve a shared name \emph{orbit-finite linear programming} 
(in case of $\F=\R$) and \emph{orbit-finite integer linear programming} (in case of $\F=\Int$).
Figure \ref{fig:diagred} shows the reductions of Theorem~\ref{thm:equiv-problems} using dashed arrows.
% shown in ,
%with  (integer) linear programming placed in frames.
%Dashed arrows depict reductions given in the proof .
% while other reductions are shown similarly:
%
% as those in the proof of Theorem~\ref{thm:equiv-problems}:
% except for the reduction from \solvname$(\F)$ to \finsolvname$(\F)$ which is shown in~\cite{GHL22}:
%
%\[
%\xymatrix@R=12pt@C=8pt{
%*+[F]{\text{\small \finineqsolvname}(\F)} \ar@{-->}@/^/[rrr]_{\text{Thm~\ref{thm:equiv-problems}}} &&& *+[F]{\text{\small \ineqsolvname}(\F)} \ar@{-->}@/^/[lll] \ar@{-->}@/^1pc/[dd] \\
%&&& \\
%\text{\small \finnonnegsolvname}(\F) \ar[rrr] \ar[uu] &&& *+[F]{\text{\small \nonnegsolvname}(\F)} \ar@{-->}@/^1pc/[uu]_{\text{Thm~\ref{thm:equiv-problems}}} 
%}
%\]

As our two main results we prove that the linear programming is decidable, while the integer one is not.
Furthermore, the complexity of the decision procedure is exponential in atom dimension of the input system,
but polynomial in the number of orbits. 
This yields \EXPTIME complexity in general, and \PTIME complexity for any fixed atom dimension
of input. 
%
\begin{theorem} \label{thm:lp-decid}
\finineqsolvname$(\R)$
%Orbit-finite linear programming 
is decidable in \EXPTIME.
For fixed atom dimension, it is decidable in \PTIME.
\end{theorem} 
%
\begin{theorem} \label{thm:ilp-undecid}
\finineqsolvname$(\Z)$
%Orbit-finite integer linear programming 
is undecidable.
\end{theorem} 
%
(The proofs % of Theorems \ref{thm:lp-decid} and \ref{thm:ilp-undecid} 
occupy Sections \ref{sec:lp-decid-proof} and
\ref{sec:ilp-undecid-proof}, respectively.)
%
Additionally, we settle the status of the last variant, \finnonnegsolvname$(\F)$.
The problem is decidable for $\F=\R$, as it reduces to both \finineqsolvname$(\F)$ and 
\nonnegsolvname$(\F)$ via reductions analogous to those of Theorem 
\ref{thm:equiv-problems}. We derive decidability also in case $\F = \Z$:
%(the proof, relying of the recent result of \cite{HR21}, is in Appendix):

%\arka{I think the result of \cite{HR21} only works for $\Int$.
%For $\Rat$ we neet S\l awek's idea/ simply adding an inequality forcing finiteness.}
%\piotrek{Yes, for $\Rat$ we reduce to linear programming. } 
%
\begin{theorem} \label{thm:finnegsolv-decid}
\finnonnegsolvname$(\F)$ is decidable, for $\F\in\set{\R,\Int}$. 
%\slawek{complexity?}\piotrek{If we want to speak about complexity then the encoding should be first.}
\end{theorem}
%
\noindent
(The proof is in Section \ref{sec:problems-proofs-two}.)
In consequence of  Theorems~\ref{thm:ilp-undecid} and \ref{thm:finnegsolv-decid}, in case $\F = \Int$
the two arrows outgoing from \finnonnegsolvname$(\Int)$ in Figure \ref{fig:diagred} can
not be completed by the reverse arrows.



%\vspace{1mm}
\begin{figure}
\begin{align*} 
%\label{eq:diagred}
%\begin{aligned}
\xymatrix@R=12pt@C=4pt{
*+[F]{\text{\small \finineqsolvname}(\F)} \ar@{-->}@/^/[rrr]_{\text{Thm~\ref{thm:equiv-problems}}} &&& *+[F]{\text{\small \ineqsolvname}(\F)} \ar@{-->}@/^/[lll] \ar@{-->}@/^1pc/[dd] \\
&&& \\
\text{\small \finnonnegsolvname}(\F) \ar[rrr] \ar[uu] 
&&& *+[F]{\text{\small \nonnegsolvname}(\F)} \ar@{-->}@/^1pc/[uu]_{\hspace{-0.6mm}\text{Thm~\ref{thm:equiv-problems}}} \\
&&& \\
\text{\small \finsolvname}(\F) \ar@/^/[rrr] \ar[uu] &&& \text{\small \solvname}(\F) \ar@/^/[lll]_{\text{\cite{GHL22}}}%^{\text{\cite{GHL22} Thm~7.1}} 
\ar[uu] 
}
%\end{aligned}
\end{align*}
\caption{Diagram of reductions between solvability problems.}
\label{fig:diagred}
\end{figure}
%\vspace{0.5mm}


\para{Linear equations vs inequalities}

Solvability of orbit-finite systems of \emph{equations} 
(\solvname$(\F)$)
easily reduces to \ineqsolvname$(\F)$,
by replacing each equation with two opposite inequalities, 
but also to \nonnegsolvname$(\F)$, by replacing each unknown with a difference of two unknowns.
Likewise does the variant \finsolvname$(\F)$, where one only seeks for finitary solutions.
%Here is the diagram of all reductions: 
%

\begin{theorem}[\cite{GHL22} Thms 4.4 and 6.1]
\solvname$(\F)$ and \finsolvname$(\F)$ are inter-reducible and decidable%
% in \EXPTIME%
\footnote{The results of~\cite{GHL22} apply to systems of equations where coefficients and solutions are from
any fixed commutative and effective ring $\F$. 
This includes integers $\Int$ or rationals $\Rat$ (and hence applies also to real solutions).}.
\end{theorem}

In summary, for each choice of $\F$ one may distinguish three different decision problems:
solving of systems of linear equations (two bottom nodes in Figure \ref{fig:diagred}),
solving of system of linear inequalities (three upper nodes in Figure \ref{fig:diagred}),
and the intermediate problem \finnonnegsolvname$(\F)$.


\para{Optimisation problems}

We consider $\F = \R$, due to the undecidability of Theorem~\ref{thm:ilp-undecid}.
All variants of linear programming mentioned above have corresponding \emph{maximisation} problems.
In each variant the input contains, except for a system $(\vr A, \vr t)$, 
an integer vector $\vr s : C \tofs \Z$ that represents
a (partial) linear \emph{objective} function $S : \GLin C \tofs \R$, defined by
\[
S(\vr x) = \innerprod{\vr s}{\vr x}.
\]
%
%\arka{We should write $\vr{s} \cdot \vr{x}$ to maintain consistency of representation}
%
The maximisation problem asks to compute the supremum of the objective function over all (finitary, nonnegative)
solutions of $(\vr A, \vr t)$.
A symmetrical \emph{minimisation} problem is easily transformed to a maximisation one by replacing 
$\vr s$ with $-\vr s$.
This yields three optimisation problems 
\ineqmaxname$(\R)$, \finineqmaxname$(\R)$ and \nonnegmaxname$(\R)$ which are,
as before, inter-reducible:
%(the proof is in Appendix):
%
%
\begin{theorem} \label{thm:equiv-max}
The problems \ineqmaxname$(\R)$, \finineqmaxname$(\R)$ and \nonnegmaxname$(\R)$
are inter-reducible, with the same complexity as in Theorem \ref{thm:equiv-problems}. 
\end{theorem}
%
(The proof is in Section \ref{sec:problems-proofs-one}.)
As our last main result we strengthen Theorem~\ref{thm:lp-decid} to the optimisation setting:
%
\begin{theorem} \label{thm:lp-max}
\finineqmaxname$(\R)$ is computable in \EXPTIME.
For fixed atom dimension, it is computable in \PTIME.
\end{theorem}
%
(The proof is in Section \ref{sec:max}.)
Hence, for every fixed atom dimension, orbit-finite linear programming is not more costly
than the classical finite linear programming.


\para{Representation of input}

There are several possible ways of representing input $(\vr A, \vr t, \vr s)$ to our algorithms.
One possibility is to rely on the equivalence between (hereditary) orbit-finite sets and
\emph{definable} sets \cite[Sect.~4]{atombook}.
We choose another standard possibility, as specified in items (1)--(3) below.
First, the representation includes:
\begin{enumerate}
\item[(1)]  a common support $T\subseteqfin\A$ of $\vr A$, $\vr t$ and $\vr s$.
\end{enumerate}
Second, knowing that $B$ and $C$ are disjoint unions of $T$-orbits, 
and relying on Lemma \ref{lem:repn}, the representation includes also:
\begin{enumerate}
\item[(2)] a list of all $T$-orbits included in $B$ and $C$, each one represented by some tuple $a\in\otu \A n$
and $G\leq S_n$; and a list of $T$-orbits included in $B\times C$, each one represented by some its element.
% a pair $(n, G)$ as in Lemma \ref{lem:repn}.  
\end{enumerate}
%We note that representations of $T$-orbits included in $B\times C$ are then computable using Lemmas
%\ref{lem:repn} and \ref{lem:orbitpair}.
Finally, relying on Lemma~\ref{lem:1O}, we assume that the representation includes also:
\begin{enumerate}
\item[(3)] a list of integer values $\orbval{\vr t}(\O)$, $\orbval{\vr s}(\O)$, and $\orbval{\vr A}(\O)$, respectively,  
for all $T$-orbits $\O$ included in $B$, $C$, and $B\times C$, respectively
(we apply Notation \ref{not:orbval}).
Integers are represented in binary.
\end{enumerate}
%
%
% we represent each $T$-orbit included in $B$ or $C$ by a pair
%and 
%assume that a vector $\vr v : D \tofs \Z$ over an orbit-finite set $D$
%is represented by:
%%
%\begin{enumerate}
%\item[(1)] its support $T = \supp {\vr v}$, 
%\item[(2)] list $L$ of all $T$-orbits $\O\subseteq D$,  
%\item[(3)] list of values $\vr v(\O) \in \Z$ for all $\O\in L$ (cf.~Lemma~\ref{lem:1O}).
%\end{enumerate}
%%
%We apply this representation to $\vr A$, $\vr t$, and $\vr s$. 
%Concerning point (2), each $T$-orbit is finitely representable, cf.~\cite[Thm 6.3]{atombook}.
%
%\arka{Suggestion start here}
%
%There are several possible ways of representing input $(\vr A, \vr t, \vr s)$ to our algorithms,
%such as [\cite[Definition 4.1]{atombook}].
%We choose another standard possibility.
%We assume $B$ and $C$ are given as finite union of $S$-orbits (for some finite $S\subseteq \A$ which supports the system) of the form $(\A \setminus S)^{(n)}/G$ where $n \in \Nat$ and $G$ is a subgroup of $S_n$,
%the set of permutations of the set $\set{1,\dots,n}$.
%$(\A \setminus S)^{(n)}/G$ is the set of non-repeating tuples of atoms in $(\A \setminus S)$ modulo coordinate permutations from the group $G$.
%This is not restrictive since using the same idea as \cite[Theorem 6.3]{atombook} we can prove the following lemma.
%%
%\begin{lemma}\label{lem:repn}
%For every finite $S \subseteq \A$,
%every $S$-orbit admits an $S$-supported effective bijection to a set of the form $(\A \setminus S)^{(n)}/G$
%for some $n \in \Nat$ and some subgroup $G$ of $S_n$.
%\end{lemma}
%%
%$\vr{A},\vr{s},\vr{t}$ are represented as vectors over the corresponding orbit-finite sets which are represented in the following way.
%A vector $\vr v : D \tofs \Z$ over an orbit-finite set $D$
%is represented by:
%%
%\begin{enumerate}
%\item[(1)] a finite subset $T \subseteq \A$,
%\item[(2)] list $L$ of all $T$-orbits $\O\subseteq D$,
%           where every $T$-orbit $X$ is represented by an element $x \in X$, and
%\item[(3)] a list of values $\vr v(\O) \in \Z$ for all $\O\in L$ (cf.~Lemma~\ref{lem:1O}).
%\end{enumerate}
%We assume the size of the above representation of a vector $\vr{v}$ is the cardinality of $L$.
%So the size of the input $(\vr{A},\vr{s},\vr{t})$ represented as system supported by $T$, with rows and columns indexed by respectively by $B$ and $C$, is total the number of $T$-orbits in $(B\times C)\uplus B \uplus C$. 
%\arka{Suggestion ends here}

\para{Strict inequalities}
In this paper we consider system of \emph{non-strict} inequalities, for the sake of presentation.
The decision procedures of 
Theorems~\ref{thm:lp-decid} and \ref{thm:lp-max}, work equally well if both \emph{non-strict and strict} inequalities are
allowed.
Reductions between \finineqsolvname$(\F)$ and \ineqsolvname$(\F)$ work as well, but not the reductions
from (\finname)\nonnegsolvname$(\F)$ to (\finname)\ineqsolvname$(\F)$
as we can not simulate equalities with strict inequalities.
%
%Nevertheless, our method works for \nonnegsolvname$(\F)$ with strict inequalities as well.

%All the reductions of Theorems~\ref{thm:equiv-problems} and \ref{thm:equiv-max}, as well as 
%\piotrek{Agree.}


\para{Proviso}
When investigating different systems of inequalities in the following sections, we
implicitly consider their \emph{real} solutions, unless specified otherwise.




 
% !TEX root = lin-prog-main.tex


\section{Polynomially-parametrised inequalities}
\label{sec:poly}


%\arka{We can shorten this section by refining the parametrised problem.
%Instead of asking if there exists $n$ for which the system is solvable,
%we can ask if there exists $\vr{x}$ which solves the system for all but finitely many $n$.}
%\slawek{problem \allpolyineqsolvname is added}

We now introduce a core problem that will serve as a target of reductions in the proofs of Theorems \ref{thm:lp-decid} and \ref{thm:lp-max}  in Sections \ref{sec:lp-decid-proof} and \ref{sec:max}.
Consider a finite inequality $\ineqal$ of the form:
%
\begin{align} \label{eq:polyeq}
p_1(n) \cdot x_1 + \ldots + p_k(n) \cdot x_k \ \geq \ q(n),
\end{align}
%
where $p_1, \ldots, p_n$ and $q$ are univariate polynomials with integer coefficients, 
and $x_1, \ldots, x_n$ are unknowns.
The special unknown $n$ plays a role of a nonnegative integer parameter, and that is why we call such an inequality
\emph{polynomially-parametrised}.
For every fixed value $n\in\Nat$, by evaluating all polynomials in $n$ we get an ordinary
inequality  $\ineqal(n)$ with integer coefficients.
Also, if $n$ does not appear in $\ineqal$, i.e., all polynomials are constants, $\ineqal$ is an ordinary inequality.

In the sequel we study solvability of a finite system $P$ of such inequalities~\eqref{eq:polyeq} 
with the same unknowns $x_1, \ldots, x_k$.
Again,
by evaluating all polynomials in $n$ we get an ordinary system $P(n)$.
We use the matrix form $P(n) = (\vr A(n), \vr t(n))$ when convenient.
A fundamental problem is to check if for some value $n\in\Nat$, the system $P(n)$ has a real solution:

\prob{\polyineqsolvname}
{A finite system of polynomially-parametrised inequalities $P$} 
{Does $P(n)$ have a real solution for some $n\in \Nat$}
%
%\begin{slremark}
%For every fixed value of $n\in\Nat$, the problem amounts to solvability of a finite system of linear
%equations, and is thus decidable in polynomial time.
%%Also, for every fixed rational values of variables $x_1, \ldots, x_k$, the problem reduces to nonnegative-integer
%%solvability of a system of univariate polynomial Diophantine inequalities, which is also decidable in polynomial time~\cite{CuckerKS99}. \slawek{does this paper solve inequalities too?} 
%\end{slremark}

\begin{theorem} \label{thm:polydecid}
\polyineqsolvname is decidable. %\TODO{\EXPSPACE}
\end{theorem}
%
\noindent
(The proof is in Section \ref{sec:poly-proofs}.)
In the sequel we will not use the decision procedure of Theorem \ref{thm:polydecid}, but rather the algorithm of Theorem \ref{thm:polyptime} stated below, since our
later applications only use \emph{monotonic} instances of \polyineqsolvname.


\subsection{Monotonic polynomially-parametrised inequalities}
\label{sec:poly-mon}

A system $P$ is \emph{monotonic} if there is some $n_0\in\Nat$ such that
every solution of $P(n)$, for an integer $n\geq n_0$, 
is also a solution of $P(n+1)$.
%\slawek{Does a monotonic system $P$ contain only monotonic inequalities?}
Note that monotonicity vacuously holds (with any value of $n_0$) 
if $n$ does not appear in $P$, i.e., when $P$ is an ordinary (non-parametrised) system.
When $P$ is monotonic, 
a solution of $P(n)$ for $n\geq n_0$ is also a solution of $P(n')$ for all integers $n'\geq n$.
%, i.e., is a solution of $P(n)$ for \emph{almost all} $n\in\Nat$ (all sufficiently large $n\in\Nat$).
%
%sufficiently large $n\in\Nat$, namely for some $n_0\in\Nat$, it is a solution of $P(n)$ for all $n > n_0$.
%When $P$ is monotonic and $P(n_0)$ has a solution for some $n_0$, 
%then there is $n_0\in\Nat$, such that every solution of $P(n_0)$ is a solution of $P(n)$ for all $n > n_0$.
%
%Any such $\vr x$ we call simple a \emph{solution of $P$}.
%Likewise we define a solution of a single polynomially-parametrised inequality $\ineqal$.
Therefore, instead of \polyineqsolvname we prefer to use in the sequel
the following core problem, where we do not assume monotonicity but
seek for a solution of $P(n)$ for \emph{almost all} 
(all sufficiently large) values of the parameter $n\in\Nat$:

\prob{\allpolyineqsolvname}
{A finite system of polynomially-parametrised inequalities $P$} 
{Is there $n_0 \in \Nat$ and a vector $(x_1, \ldots, x_k)$ which is a solution of $P(n)$ for every integer $n\geq n_0$}
%
%\arka{$\exists$solution which works $\forall n\geq n_0$.
%Not $\forall n \geq n_0 \exists$ a solution.}\piotrek{Arka , i think this is what is written every solution for $n$ is also a solution for $n+1$}
%
\noindent
From now on, a vector $(x_1, \ldots, x_k)$ which is a solution of an inequality $\ineqal(n)$ 
(resp.~a system $P(n)$) for almost all $n\in\Nat$ we call \emph{\aasol} of 
%a polynomially-parametrised inequality 
$\ineqal$ (resp.~$P$).
The rest of this section is devoted to designing a \PTIME algorithm for \allpolyineqsolvname:
% and thus also for monotonic \polyineqsolvname:

\begin{theorem} \label{thm:polyptime}
\allpolyineqsolvname is %decidable 
in \PTIME.
%(and hence so is monotonic \polyineqsolvname).
\end{theorem}
%
\begin{proof}
Consider a polynomially-parametrised 
inequality $\ineqal$ of the form:
%
\begin{align} \label{eq:polyeqpowt}
p_1(n) \cdot x_1 + \ldots + p_k(n) \cdot x_k \ \geq \ q(n).
\end{align}
%
Let $d$  be the maximal degree of polynomials $p_1, \ldots, p_k, q$ appearing
in $\ineqal$.
We call $d$ \emph{the degree} of $\ineqal$, and denote it also as $\deg \ineqal$.
Let $a_1, \ldots, a_k, b$ be (integer) coefficients of the monomial $n^{d}$ in 
$p_1,  \ldots, p_k, q$,
respectively.
Therefore
\begin{align} \label{eq:d}
\begin{aligned}
p_1(n) & = a_1 \cdot n^{d} + p'_1(n) \qquad \ldots \qquad 
p_k(n) = a_k \cdot n^{d} + p'_k(n) \qquad\qquad
q(n) & = b \cdot n^{d} + q'(n)
\end{aligned}
\end{align}
%\arka{better to write $p(n)$}
for some polynomials $p'_1, \ldots, p'_k, q$ of degree strictly smaller than $d$.
The ordinary inequality with integer coefficients
\begin{align} \label{eq:head}
a_1 \cdot x_1 + \ldots + a_k \cdot x_k \geq b,
\end{align}
we call \emph{the head inequality} of $\ineqal$, and denote by
$\li{\ineqal}$.
Furthermore, the polynomially-parametrised inequality 
\begin{align} \label{eq:tl}
p'_1(n) \cdot x_1 + \ldots + p'_k(n) \cdot x_k \ \geq \ q'(n),
\end{align}
obtained by removing all appearances of the monomial $n^d$,
we call the \emph{tail} of $\ineqal$, and denote it by $\tl \ineqal$.
We also consider below the strict strengthening of the head inequality \eqref{eq:head}, denoted as $\strictli \ineqal$, and the equality, denoted as $\lipar \ineqal =$:
\begin{align} \label{eq:strictli}
%p_1(n) \cdot x_1 + \ldots + p_k(n) \cdot x_k \ > \ q(n)
%\qquad\qquad
%p_1(n) \cdot x_1 + \ldots + p_k(n) \cdot x_k \ = \ q(n),
%
a_1 \cdot x_1 + \ldots + a_k \cdot x_k > b,
\qquad\qquad\qquad
%\end{align}
%\begin{align} \label{eq:eqli}
a_1 \cdot x_1 + \ldots + a_k \cdot x_k = b.
\end{align}
%
%and their heads $\li {\strict \ineqal}$ and $\li {\eqal \ineqal}$:
%
As $\ineqal$ is equal to the sum of its head $\li \ineqal$ multiplied by $n^d$, 
and its tail $\tl \ineqal$, we immediately deduce:
%
\begin{claim} \label{claim:hdtl}
For every $n\in\Nat$, every solution of $\lipar {\ineqal} =$ is either a solution of both $\ineqal(n)$ and 
${\tl \ineqal}(n)$, 
or of none of them.
\end{claim}

We now provide under- and over-approximations
of the solution set of $\ineqal$ (in Claims \ref{claim:strictlibase} and \ref{claim:libase}).
%will serve as a core of our algorithm.

\begin{claim} \label{claim:libase}
Every \aasol of $\ineqal$ is also a solution of $\li \ineqal$.
\end{claim}
\begin{proof}
%Let $\ineqal$ be as in \eqref{eq:polyeqpowt}.
%The claim follows immediately by application of Claim \ref{claim:strictlibase} to the reverse 
%$\ineqal^{-1}$ of $\ineqal$,
%\[
%p_1(n) \cdot x_1 + \ldots + p_k(n) \cdot x_k \ \leq \ q(n),
%\]
%obtained by reversing the direction of inequality.
%Indeed, suppose $\vr x$ is not a solution of $\li \ineqal$ \eqref{eq:head}, i.e., it is a solution
%of ${\strictli \ineqal}^{-1}$, which is the same as $\strictli {\ineqal^{-1}}$.
%By Claim \ref{claim:strictlibase}, $\vr x$ satisfies $\ineqal^{-1}$, which means that
%for all sufficiently large it satisfies $\ineqal^{-1}(n)$.
%is does not
%satisfy $\ineqal(n)$. 
%
Consider an inequality $\ineqal$ \eqref{eq:polyeqpowt}
and  its \aasol  $\vr x = (x_1, \ldots, x_k)$.
Let $d = \deg \ineqal$.
We thus have
\[
\frac{p_1(n)}{n^d} \cdot x_1 + \ldots + \frac{p_k(n)}{n^d} \cdot x_k \ \geq \ \frac{q(n)}{n^d}
\]
for all sufficiently large $n\in\Nat$. 
Using the decomposition \eqref{eq:d}, we rewrite the above inequality to
\[
\Big(a_1 + \frac{p'_1(n)}{n^d}\Big) \cdot x_1 + \ldots + 
\Big(a_k + \frac{p'_k(n)}{n^d} \Big) \cdot x_k \ \geq \ 
b + \frac{q'(n)}{n^d}.
\]
As the degrees of all polynomials $p'_1, \ldots, p'_k, q'$ are smaller than $d$, 
all the fractions tend to 0 when $n$ tends to $\infty$, and we may deduce
\[
a_1 \cdot x_1 + \ldots + a_k \cdot x_k \ \geq \  b,
\]
i.e., $\vr x$ is a solution of $\li \ineqal$, as required.
%
%\emph{not} a solution of
%$\li \ineqal$,
%%This means that for some inequality $\ineqal$ \eqref{eq:polyeq} in $P$, 
%%its head inequality \eqref{eq:head} is not satisfied by $\vr x$, 
%namely:
%\begin{align} \label{eq:notsat}
%a_1 \cdot x_1 + \ldots + a_k \cdot x_k < b.
%\end{align}
%\arka{the contradiction is unneccessary}
%We need to show that for all sufficiently large $n\in\Nat$, $\vr x$ is not a solution of $\ineqal(n)$.
%%(In consequence, $\vr x$ is not a solution of $P$.)
%%
%Recall the decomposition \eqref{eq:d} of polynomials appearing in $\ineqal$.
%From \eqref{eq:notsat} we deduce that
%\[
%\Big(a_1 + \frac{p'_1(n)}{n^d}\Big) \cdot x_1 + \ldots + 
%\Big(a_k + \frac{p'_k(n)}{n^d} \Big) \cdot x_k <
%b + \frac{q'(n)}{n^d}
%\]
%holds for all sufficiently large $n\in\Nat$.
%Equivalently,
%\[
%\frac{p_1(n)}{n^d} \cdot x_1 + \ldots + \frac{p_k(n)}{n^d} \cdot x_k \ < \ \frac{q(n)}{n^d}
%\]
%holds for all sufficiently large $n\in\Nat$.
%Multiplying both sides by $n^d$ yields
%\begin{align*} %\label{eq:polyeqtail}
%p_1(n) \cdot x_1 + \ldots + p_k(n) \cdot x_k \ < \ q(n),
%\end{align*}
%i.e., for all sufficiently large $n\in\Nat$,
%$\vr x$ is not a solution of $\ineqal(n)$, as required.
\end{proof}

\begin{claim} \label{claim:strictlibase}
Every solution of $\strictli {\ineqal}$ is also an \aasol of $\ineqal$.
\end{claim}
%
\begin{proof}
Let $d = \deg \ineqal$. 
Consider any vector $\vr x = (x_1, \ldots, x_k)$  satisfying
the strict inequality $\strictli {\ineqal}$ (in \eqref{eq:strictli} on the left). 
Therefore for any polynomials $p'_1, \ldots, p'_k, q'$ of degree strictly smaller than $d$,
the inequality
\[
\Big(a_1 + \frac{p'_1(n)}{n^d}\Big) \cdot x_1 + \ldots + 
\Big(a_k + \frac{p'_k(n)}{n^d} \Big) \cdot x_k > 
b + \frac{q'(n)}{n^d}
\]
is satisfied for all sufficiently large $n\in\Nat$.
Applying the above inequality to polynomials appearing in \eqref{eq:d}, we obtain:
\[
\frac{p_1(n)}{n^d} \cdot x_1 + \ldots + \frac{p_k(n)}{n^d} \cdot x_k \ > \ \frac{q(n)}{n^d}
\]
for all sufficiently large $n\in\Nat$.
We multiply both sides by $n^d$ in order to derive that 
$\vr x$ is a solution of $\ineqal(n)$ for all sufficiently large $n\in\Nat$,
as required.
\end{proof}

%\smallskip

Consider an instance $P$ of \polyineqsolvname, i.e., a finite system of 
polynomially-parametrised inequalities of the form \eqref{eq:polyeqpowt}.
%and assume that $P$ is monotonic.
Let $\li{P} := \setof{\li{\ineqal}}{\ineqal\in P}$ be the system of head inequalities
(note that degrees of different inequalities in $P$ may differ), and let
$\strictli{P} := \setof{\strictli{\ineqal}}{\ineqal\in P}$.
Using Claims \ref{claim:strictlibase} and \ref{claim:libase} we derive:

\begin{claim} \label{claim:li}
Every \aasol of $P$ is also a solution of \xspace $\li P$.
\end{claim}

\begin{claim} \label{claim:strictli}
Every solution of $\strictli P$ is also an \aasol of $P$.
\end{claim}

%

For time estimation, as the size measure $\size \ineqal$ of an inequality $\ineqal$ we take the total number of monomials
appearing in $\ineqal$.
In particular, $\size \ineqal > \size {\tl \ineqal}$.
The size of a system $P$ is the sum of sizes of all its inequalities.
% as the length of representation of all  coefficients and exponents of monomials appearing in $P$.
%Let $\deg P$ denote the sum of degrees of all inequalities in $P$:
%\[
%\deg P \ := \ \sum_{\ineqal \in P} \deg \ineqal.
%\]
For two systems $P'$, $P''$ of inequalities, we denote their union by $P'\conj P''$
(clearly, union of systems corresponds to conjunction of constraints).
We write $P\conj\ineqal$ instead of $P\conj\set{\ineqal}$.
By $P\setminus \ineqal$ we denote the system obtained from $P$ by removing an inequality $\ineqal$.

\para{The algorithm}
A decision procedure for \allpolyineqsolvname iteratively transforms an instance
of the form $P \conj \G$, where $P$ is a system
of polynomially-parametrised inequalities, and $\G$ is a system of ordinary (non-parametrised)
equalities over the same unknowns.
% (each equality can be represented as two opposite inequalities). 
Initially, $\G$ is empty.
We define a transformation step that given such an instance $P \conj\G$, 
either confirms its solvability (existence of an \aasol), 
or confirms its non-solvability (non-existence of an \aasol), or outputs an instance $P' \conj \G'$ 
which has the same {\aasol}s as $P\conj \G$, and 
such that $\size {P'} < \size P$.
\allpolyineqsolvname is solved by iterating the transformation step until it confirms
either solvability or non-solvability.
Termination after a polynomial number of iterations is guaranteed, as
$\size P$, while being nonnegative, strictly decreases in each iteration.
The transformation step invokes a \PTIME procedure 
%\slawek{what should we cite?}
%\piotrek{Why current citation is not good enough? The step trom $\leq$ to $<$ can be done like this, I think. $\leq$ is solvable in $poly$-time so every corner of the convex polyhidron has poly lenght description, so if you want to solve $< x$ it is sufficient to solve $\leq x-\epsilon$ for small epsilon and epsilon has polynomialy many bits. I don't know precise formula for epsilon but it can be calculated, it has to be so small that it will not remove any corner, except ones on the hyperplane.}
%\ordineqsolvname 
for 
%real solvability of ordinary (non-parametric) systems of inequalities with integer coefficients or 
ordinary linear programming
(as detailed in \eqref{eq:>=} and \eqref{eq:=} below).
Here is a pseudo-code of the algorithm:

%\begin{minipage}{#1\linewidth}
\begin{algorithm}[H]
\caption{(\allpolyineqsolvname)}%[H]                                                                                                                                                             
\label{alg:1}
\begin{algorithmic}[1]
\State \textbf{Input:} \ A polynomially-parametrised system $P$.
\smallskip
\State $\G \gets \emptyset$
\Repeat
\smallskip
\If{\ $\li P  \conj  \G$ \ \eqref{eq:>=} \  is non-solvable} \label{c:if1}
\Comment see Claim \ref{claim:li}
%\State answer \noanswer
\State report non-solvability of $P$
\Else
\If{$\strictli \ineqal \conj \li {P\setminus \ineqal} \conj \G$ \ \eqref{eq:=} \ is solvable for all $\ineqal \in P$}
%\If {\ $\strictli P  \conj  \G$ \ \eqref{eq:>} \ is solvable} 
\label{c:if2}
\Comment see Claims \ref{claim:strictli}, \ref{claim:=}
%in consequence, $\strictli P  \conj  \G$  \ is solvable
\State report solvability of $P$
\Else
\State \textbf{choose any} $\ineqal \in P$ \textbf{such that} \ $\strictli \ineqal \conj \li {P\setminus \ineqal} \conj \G$ \ \eqref{eq:=} \ is non-solvable \label{l:lastif}
\Comment see Claim \ref{claim:samesol}
\smallskip
\State $P \, \gets \, (P \setminus \ineqal) \ \conj \ \tl \ineqal$
\State $\G \,\, \gets \, \G \ \conj \ \lipar \ineqal {=}$
\EndIf 
\EndIf
\Until{solvability or non-solvability of $P$ is reported}
%\medskip
%\State \textbf{Output:} \ $\G$
\end{algorithmic}
\end{algorithm}                                                                                                                                                                  
%\end{minipage}


%\medskip

\para{Transformation step}
The step, defined by the body of the 
\textbf{repeat} loop, proceeds as follows.
If the ordinary system
\begin{align} \label{eq:>=}
%Invoke \ordineqsolvname$(\li P\cup\F)$.
\li P \; \conj \; \G
\end{align}
is non-solvable, 
non-solvability of $P\conj\G$ is reported.
This is correct due to Claim \ref{claim:li}.
%
Otherwise, knowing that \eqref{eq:>=} is solvable,
the algorithm checks, for every $\ineqal\in P$, whether 
the strengthened system 
\begin{align} \label{eq:=}
\strictli \ineqal \; \conj \; \li {P\setminus \ineqal} \; \conj \; \G,
\end{align}
obtained from \eqref{eq:>=} by replacing
the inequality $\li \ineqal$ with $\strictli \ineqal$, is also solvable.
If this is the case, solvability of $P \conj \G$ is reported. 
This is correct due to Claim \ref{claim:strictli} combined with the following one:
%
\begin{claim} \label{claim:=}
Solvability of \eqref{eq:=} for every inequality $\ineqal$ in $P$, implies solvability of
\begin{align}
\label{eq:>}
\strictli P \, \conj \, \G.
\end{align}
%\eqref{eq:>} 
\end{claim}
%
\begin{proof}
Let $m$ be the number of inequalities in $P$, and
suppose that for every inequality $\ineqal$ in $P$, the system \eqref{eq:=}
has a solution, $\vr x_\ineqal$.
All $\vr x_\ineqal$ are thus solutions of \eqref{eq:>=}, and since
the solution set of \eqref{eq:>=} is convex,
the average of all these solutions
$
\frac 1 {m} \cdot \sum_{\ineqal\in P} \vr x_\ineqal
$
is then a solution of $\strictli P \, \conj \, \G$.
\end{proof}
%
Otherwise, we know that some
inequality $\ineqal$ in $P$ is \emph{degenerate}, namely
\eqref{eq:=} is non-solvable.
In other words, the equality $\lipar \ineqal =$
is implied by \eqref{eq:>=}.
The algorithm chooses a degenerate inequality $\ineqal\in P$ 
%
%if the ordinary system
%\begin{align} \label{eq:>}
%\strictli P \; \conj \; \G
%\end{align} 
%is solvable,
%solvability of $P\conj\G$ is reported.
%This is correct due to .
%
and creates a new instance $P' \conj \G'$, where
\begin{align*}
P' \  = \ (P \setminus \ineqal) \; \conj\; \tl \ineqal \qquad\qquad
\G' \  = \ \G \; \conj \; \lipar \ineqal {=}.
\end{align*}
%
In words, $P'$ is obtained from $P$ by replacing $\ineqal$ with $\tl \ineqal$,
and $\G'$ is obtained from $\G$ by adding $\lipar \ineqal =$.
As $\size {\tl \ineqal} < \size \ineqal$, we have $\size {P'} < \size P$, as required.
This completes description of the transformation step.

\para{Correctness}
By Claim \ref{claim:hdtl} we derive:
%
\begin{claim}  \label{claim:samesol}
Systems $P\conj\G$ and $P'\conj\G'$ have the same {\aasol}s.
\end{claim}
%
\begin{proof}
In one direction, 
consider an \aasol $\vr x$ of $P'\conj\G'$.
It is trivially a solution of $\G$. 
Furthermore, being a solution of $\lipar \ineqal =$ and of $\tl \ineqal(n)$ for almost all $n\in\Nat$,
by Claim \ref{claim:hdtl} it is a solution of $\ineqal(n)$ for almost all $n$, and hence an \aasol of $P$.

Conversely, 
consider an \aasol $\vr x$ of $P\conj\G$.
By Claim \ref{claim:li}, it is a solution of $\li P\conj\G$ and hence, 
as $\ineqal$ is degenerate,
%due to Claim \ref{claim:=}, 
also a solution of $\lipar \ineqal =$.
Therefore $\vr x$ is a solution of $\G'$.
Furthermore, being a solution of $\lipar \ineqal =$ and of $\ineqal(n)$ for all sufficiently large
$n\in\Nat$, 
by Claim \ref{claim:hdtl} it is also a solution of $\tl \ineqal(n)$ for all sufficiently large $n\in\Nat$, 
and hence an \aasol of $P'$.
\end{proof}
%

\para{Complexity}
We note that the main loop of the algorithm always terminates, at latest when $P=\emptyset$, as in this case the system
\eqref{eq:=} is vacuously solvable for all $\ineqal\in P$.
%
%there is no difference between \eqref{eq:>=} and \eqref{eq:>}.
%
Solvability of \eqref{eq:>=} in line \ref{c:if1} 
is checked by one solvability test of an ordinary system of inequalities.
Solvability of \eqref{eq:=} in line \ref{c:if2} is also checkable in polynomial time due to the following claim
applied to $Q = \li {P} \; \conj \; \G$:
%
\begin{claim} 
Given an ordinary system $Q$ of linear inequalities and $\ineqal \in Q$, one can check,
in \PTIME,
solvability of \xspace $\ineqal_>  \conj \, (Q\setminus\ineqal)$, where $\ineqal_>$ is the strict
strengthening of $\ineqal$.
\end{claim}
%
\begin{proof}
We invoke ordinary linear programming twice 
%in each case constrained by non-strict inequalities 
(in \PTIME, see e.g.~\cite[Section 8.7]{PSbook82}).
Let $\ineqal$ be of the form $a_1 \cdot x_1 + \ldots + a_k \cdot x_k \geq b$.
If $Q$ is non-solvable, the algorithm reports non-solvability of 
$\ineqal_>  \conj \, (Q\setminus\ineqal)$.
Otherwise, the algorithm computes the supremum $M\in\Rat\cup \set{\infty}$ of the objective function 
\[
S(x_1, \ldots, x_k) = a_1 \cdot x_1 + \ldots + a_k \cdot x_k,
\]
constraint by $Q\setminus\ineqal$, by invoking ordinary linear programming.
By solvability of $Q$ we know that $M\geq b$.
If $M > b$, the algorithm reports solvability, otherwise it reports non-solvability.
\end{proof}


Number of iterations of transformation step is polynomial (as $\size P$ decreases in each iteration) and hence
so is the number of inequalities in $\G$.
In consequence, the number of invocations of ordinary linear programming is polynomial in
each transformation step, and hence polynomial in total, and each its instance
of ordinary linear programming is also polynomial.
Summing up, our decision procedure for \allpolyineqsolvname works in \PTIME.

\smallskip
The proof of Theorem \ref{thm:polyptime} is thus completed.
%
%Let $d$ be the maximal degree of polynomials appearing in $\vr A$ and $\vr t$.
%Therefore $\vr A$ and $\vr t$ can be presented as a sum
%\begin{align*}
%\vr A \ & = \ \vr A_d \cdot n^d + \vr A_{d-1} \cdot n^{d-1} +  \ldots + \vr A_1 \cdot n + \vr A_0 \\
%\vr t \ & = \ \ \vr t_d \cdot n^d  + \ \, \vr t_{d-1} \cdot n^{d-1}  +  \ldots + \ \vr t_1 \cdot n + \ \vr t_0 
%\end{align*}
%where $\vr A_d, \vr A_{d-1}, \ldots, \vr A_1, \vr A_0$ are rational matrices of the same dimension
%as $\vr A$, and likewise $\vr t_d, \vr t_{d-1}, \ldots, \vr t_1, \vr t_0$ are rational vectors of the
%same length as $\vr t$.
\end{proof}

\begin{slremark}
We do not need any explicit bound on the threshold value of $n_0$ guaranteeing that 
every \aasol of $P$ is a solution of $P(n)$ for every integer $n \geq n_0$.
On the other hand, an exponential bound is derivable from our algorithm.
Assuming $P$ has an \aasol, $P$ has also an \aasol $\vr x$ which is at most exponentially large, 
e.g., a solution of an ordinary system \eqref{eq:>}.
Substituting $\vr x$ into $P(n)$ yields a system of univariate polynomial inequalities, and 
one can take as threshold $n_0$ any integer larger than all nonnegative roots of all polynomials appearing
in the system. 
As roots of univariate polynomials are polynomially bounded, we deduce the bound for $n_0$.
\end{slremark}

\begin{slexample}
Recall two polynomially-parametrised inequalities \eqref{eq:ppar} in Example \ref{ex:poly}
in Section \ref{sec:intro}. They have the same head inequality
$
x\geq 1,
$
which is trivially solvable, and hence the algorithm reports solvability after the first iteration.
Both the tail inequalities,
$-x \geq 1$ and
$0 \geq 1$,
are ordinary (non-parametrised).

The following instance $P_0$ admits three iterations of the main loop of the algorithm:
\begin{align*}
n^2 \cdot x \ - \ n^2 \cdot y \  +  \ n \cdot z  & \ \geq \ 0 \\
-n \cdot x \ + \ (n+3) \cdot y \ \ \ \ \ \, & \ \geq \ 0
\end{align*}
The head inequalities of these two inequalities are $x-y \geq 0$ and $-x + y \geq 0$, respectively.
Therefore
the system $\li {P_0}$ is equivalent to $x=y$ and hence solvable, while $\strictli {P_0}$ is not,
and both inequalities in $P_0$ are degenerate.
Supposing the first one is chosen by the algorithm, after the first iteration we get the following systems $P_1$
(left) and $\G_1$ (right): 

\begin{minipage}{0.5\linewidth}
\begin{align*}
n \cdot z  & \ \geq \ 0\\
-n \cdot x \ + \ (n+3) \cdot y  & \ \geq \ 0
\end{align*}
\vspace{-2mm}
\end{minipage}
\begin{minipage}{0.5\linewidth}
\begin{align*}
x \ - \ y  & \ = \ 0 \\
\end{align*}
\vspace{-2mm}
\end{minipage}

\noindent
In the second iteration, 
the system $\li{P_1} \conj \G_1$ (left) is solvable but the system $\strictli{P_1} \conj \G_1$ (right) 
is not:

\begin{minipage}{0.5\linewidth}
\begin{align*}
z  & \ \geq \ 0\\
- x \ + \  y  & \ \geq \ 0 \\
x \ - \ y  & \ = \ 0 
\end{align*}
\vspace{-2mm}
\end{minipage}
\begin{minipage}{0.5\linewidth}
\begin{align*}
z  & \ > \ 0\\
- x \ + \  y  & \ > \ 0 \\
x \ - \ y  & \ = \ 0  
\end{align*}
\vspace{-2mm}
\end{minipage}

\noindent
The algorithm picks up the second inequality in $P_1$, the only degenerate one, and sets
$P_2$ (left) and $\G_2$ (right):

\begin{minipage}{0.5\linewidth}
\begin{align*}
n \cdot z  & \ \geq \ 0\\
3 \cdot y  & \ \geq \ 0
\end{align*}
\vspace{-2mm}
\end{minipage}
\begin{minipage}{0.5\linewidth}
\begin{align*}
x \ - \ y  & \ = \ 0 \\
-x  \ + \ y & \ = \ 0
\end{align*}
\vspace{-2mm}
\end{minipage}

\noindent
In the last third iteration, the system $\strictli {P_2} \conj \G_2$ 
(obtained by replacing the inequality $n\cdot z\geq 0$ by $z > 0$)
is solvable, and hence solvability of $P_0$ is reported.
\end{slexample}

%
%\arka{Is \emph{upward closed} a better name than \emph{monotonic}?}
%\slawek{'monotonic' sounds better in optimisation variant}
%

%One can express \polyineqsolvname in terms of evaluation of an arithmetic formula, as follows.
%Consider real arithmetic, i.e., the first-order theory of $(\R, +, \cdot, 0, 1, \leq)$.
%
%We fix the real arithmetic language: $+, \cdot, \leq, 0, 1$, with standard meanings.
%\probext{}
%{a quantifier-free arithmetic formula}{$\varphi(n, x_1, \ldots, x_k)$}
%{is $\varphi$ true for some $n\in\Nat$ and nonngative}{$x_1, \ldots, x_k\in\R$}

%
%\begin{lemma} \label{lem:1var-iff}
%The following conditions are equivalent for a real arithmetic formula with one free variable $x$:
%\begin{itemize}
%\item The formula is true for all sufficiently large real values of $x$;
%\item The formula is true for all sufficiently large integer values of $x$.
%\end{itemize}
%\end{lemma}
%
%%
%\begin{proof}
%We only need to prove the if direction.
%Due to Lemma~\ref{lem:1var},
%if a real arithmetic formula $\alpha(x)$ is true for all sufficiently large integers, then the set defined by $\alpha$ contains 
%necessarily a right-infinite interval, and hence the formula is true for all sufficiently large reals as well.
%\end{proof}

% !TEX root = lin-prog-main.tex
 


\section{Finitely setwise-supported sets} \label{sec:weak}

In this section
we introduce the novel concept of \emph{setwise-support}, playing a central role in the proofs of Theorems \ref{thm:lp-decid} and \ref{thm:lp-max}.
In short, we replace \emph{pointwise} stabilisers by \emph{setwise} ones. 
%the reduction of \finineqsolvname$(\R)$ to \polyineqsolvname in Section \ref{sec:lp-decid-proof}.
%
%Below, let $S\subseteqfin\A$ be any fixed finite set of atoms.


For any $T\subseteqfin\A$ consider the set of all atom automorphisms
that preserve $T$ as a set only (called \emph{setwise-$T$-automorphisms}):
\[
\Aut {\set{T}} \ = \ \setof{\pi \in \Aut {}}{\pi(T) = T}.%
\footnote{
$\Aut {\set{T}}$ is often called the \emph{setwise stabilizer} of $T$.
}
\]
%\red{The set} $\Aut {\set{T}}$ we call \emph{\red{setwise stabilizer of $T$}}.
Accordingly, we define \emph{setwise-$T$-orbits} as equivalence classes with respect to the action of $\Aut {\set{T}}$:
%\piotrek{change two elements to ``two set with atoms'', in general when i was reading it it was confusing what is what so we should avoid word elements.} 
two 
sets (elements) $x, y$ are in the same setwise-$T$-orbit if $\pi(x) = y$ for some $\pi \in\Aut {\set{T}}$.
We have
\[
\Aut {T} \subseteq \Aut{\set{T}} \subseteq \Aut {},
\]
and hence every equivariant orbit splits into finitely many setwise-$T$-orbits, 
each of which splits in turn into finitely many $T$-orbits.
A set $X$ is \emph{setwise-$T$-supported} if $\pi(X) = X$ for all $\pi \in \Aut {\set{T}}$.
Equivalently, $X$ is a union of setwise-$T$-orbits.
Note that each setwise-$T$-supported set is $T$-supported, but the opposite implication is not true.
When $T$ is irrelevant, we speak of finitely setwise-supported sets.
Finally notice that a setwise-$T$-supported set is not necessarily setwise-$T'$-supported for $T\subseteq T'$,
which distinguishes setwise-support from  standard support.
%\arka{Maybe we should illustrate this difference between support and setwise-support by an example?}
%\piotrek{I think we need to add simple example here. This is one of central definitions in the paper, so we should be sure that reader understands it.}
%\slawek{added example below}

\begin{slexample}
\label{ex:av}
Let $T=\set{\a,\b}\subseteq\A$.
The vector $\vr v$, defined in Example \ref{ex:vect} in Section \ref{sec:problems}, is not setwise-$T$-supported.
Indeed,
\[
\pi(\vr v)(\a, \x) = \vr v(\b, \x) \neq \vr v(\a, \x) 
\]
for any $\x \notin T$ and $\pi \in \Aut{\set{T}}$ that swaps $\a$ and $\b$ but preserves all other atoms.
The \emph{averaged} vector $\vr v'$ defined by
\begin{align*}
\vr v'(\a \x) &= \vr v'(\x \a) = - 1.5 &
\vr v'(\a \b) &= \vr v'(\b \a) = 3  \\ 
\vr v'(\b \x) &= \vr v'(\x \b) = - 1.5
&
\vr v'(\x \g) &= 0,
\end{align*}
for  $\x,\g\in\A\setminus\set{\a,\b}$, is setwise-$T$-supported.
Notice that $\vr v'$, is not setwise-$(T\cup\set{\g})$-supported, for $\g\notin T$.
%\arka{averaged vector $\vr{v}' = \frac{1}{2}\sum_{\pi \in \Aut{\A \setminus T}} \pi(\vr{v})$}\piotrek{avereged is not yet defined}
\end{slexample}

Clearly, with the size of $T$ increasing towards infinity, the number of $T$-orbits 
included in one equivariant orbit may increase towards infinity as well.
The crucial property of setwise-$T$-supported sets is that they do not suffer from this unbounded growth:
the number of setwise-$T$-orbits included in a fixed equivariant orbit is bounded, no matter how large $T$ is.
We will need this property for setwise-$T$-orbits $U\subseteq\otu {\A} n$,  $n\in\Nat$, and
it follows immediately by Lemma~\ref{lem:STorbit}.
%
Intuitively speaking, each such setwise-$T$-orbit is determined by 
a subset $I\subseteq \setto n$ of positions which is filled by arbitrary pairwise different atoms from $T$, 
the remaining positions $\setto n\setminus I$ are filled by arbitrary atoms from $\A\setminus T$ 
(cf.~Lemma \ref{lem:tupleorbit} in Section \ref{sec:atoms}).

\begin{lemma} \label{lem:STorbit}
Let $T\subseteqfin\A$ of size $\size T\geq n$.
Each setwise-$T$-orbit $\O \subseteq \otu {\A} n$ is of the form
%
\begin{align} \label{eq:STorbit}
\begin{aligned}
\O =  \setof{a \in \otu \A n}{\proj n I(a) \in \otu T \ell, \ \  
  \proj n {\setto n \setminus I}(a) \in \otu {\big(\A\setminus T\big)} {n-\ell}},
\end{aligned}
\end{align}
%
for some $I\subseteq \setto n$ of size $\ell$.
\end{lemma}
%
\begin{proof}%[Proof of Lemma \ref{lem:STorbit}]
Consider any tuple $t = (\a_1, \ldots, \a_n) \in \otu {\A} n$.
Let $I = \setof{i \in \setto n}{\a_i \in T}$ denote the positions in $t$ filled by atoms from $T$.
By applying all setwise-$T$-automorphisms to $t$, we obtain all tuples, where positions from $I$ are arbitrarily 
filled by elements of $T$, %(note that $\size T \geq k$), 
and positions outside of $I$ are arbitrarily
filled by elements of $\A\setminus T$.
\end{proof}
%
%The number $\ell$ of positions occupied by elements of $T$ we call
%\emph{$T$-order} of $O$ of the form \eqref{eq:STorbit},
%and denote as $\ord T {O}$.
%\slawek{do we need this?}

%As a direct corollary we have:
%\begin{lemma} \label{lem:STorbits}
%Let $k \in \Nat$ and $T\subseteqfin\A$ be such that $\size T\geq k$ and $S\cap T = \emptyset$.
%% be of size at least $k$, and disjoint from $S$.
%Then $\otu {(\A\setminus S)} k$ splits into $2^k$ many $S,T$-orbits.
%\slawek{no more needed}
%\end{lemma}
%
%The subset $\size I\subseteq \setto k$ of positions corresponding to an $S,T$-orbit $O \subseteq \otu {(\A\setminus S)} k$
%we call $T$-\emph{index} of $O$.
%
%Another crucial property is that the number of $(S\cup T)$-orbits inside an $S,T$-orbit
%grows polynomially with respect to $n = \size T$:
%
%\begin{lemma}
%Let $k, \ell \in \Nat$, let $T\subseteqfin\A$ be of size $n\geq \ell$, and disjoint from $S$,
%and let $O\subseteq \otu {(\A\setminus S)} k$ be an $S,T$-orbit whose $T$-index has size $\ell$.
%Then $O$ splits into $p_\ell(n)$ many $(S\cup T)$-orbits, where
%\[
%p_\ell(n) \ = \ n \cdot (n-1) \cdot \ldots \cdot (n-\ell+1).
%\]
%\end{lemma}
%%
%\begin{proof}
%An $(S\cup T)$-orbit $M\subseteq O$ differs from $O$ only by fixing specific chosen $\ell$ 
%elements of $T$ on positions from $I$.
%There are exactly $p_\ell(n)$ different ways of fixing these elements.
%\end{proof}
%%
%Importantly, the number of $(S\cup T)$-orbits depends on the size $n$ of $T$, but not on $T$ itself.
%Note that $\ell$ is bounded by $k$, and 
%for every fixed $\ell$, the function $p_\ell(n)$ is an univariate polynomial in $n$.

\begin{notation} \label{not:orbsum}
%Recall that $\orbits C$ stands for the set of equivariant orbits included in $C$.
Given a finitary vector $\vr x : C\tofin\R$ and an equivariant orbit $\O\subseteq C$, 
we write
\[
\orbsum {\vr x}(\O) \ = \ \sum_{c\in \O} \vr x(c)
\]
to denote for the sum of $\vr x(c)$ ranging over all $c\in \O$.
This yields the finite \emph{orbit-sum} vector
\[
\orbsum {\vr x} : \orbits C \to \R
\]
mapping the equivariant orbits included in $C$ to $\R$.
\end{notation}

A key observation is that a solvable equivariant system necessarily has a finitely setwise-supported solution:
%
\begin{lemma} \label{lem:weaksol}
If an equivariant system of inequalities $(\vr A, \vr t)$ has a finitary $T$-supported solution 
$\vr x$ then it also has a finitary
setwise-$T$-supported one $\vr y$ such that
$\orbsum{\vr x} = \orbsum{\vr y}$.
% and
%$L(\vr x) = L(\vr y)$ for every equivariant partial linear map $L : \GLin C \to \R$. %\piotrek{why $h$ and not $s$?}
\end{lemma}
%
%\arka{$S$ also has to be linear} 
%\slawek{no, it does not have to!}
%
\begin{proof}
Let $\vr x : C\tofs \R$ be a solution of the system, namely $\innerprod{\vr A}{\vr x}\geq \vr t.$
Let $T = \supp {\vr x}$ and $n= \size T$.
As $(\vr A, \vr t)$ is equivariant,
atom automorphisms 
%that only permute atoms inside $T$ necessarily 
preserve being a solution, namely
for every $\rho \in \Aut {}$, the vector $\rho(\vr x)$ is also a solution:
$
\innerprod {\vr A} \rho(\vr x) \ \geq \  \vr t.
$
Consider $\Aut{\A\setminus T}$, the subgroup of atom automorphisms that only permute $T$ and preserve
all other atoms.
Knowing that the size of $\Aut{\A\setminus T}$ is $n!$, we have
\[
\innerprod {\vr A} \Big(\sum_{\rho \in \Aut{\A\setminus T}}\rho(\vr x)\Big) \ \geq \ n! \cdot \vr t,
\]
and hence the vector $\vr y$ defined by averaging (cf.~Example \ref{ex:av})
\begin{align} \label{eq:defy}
\vr y \ = \ \frac{1}{n!} \ \cdot \ \sum_{\rho \in \Aut{\A\setminus T}} \rho(\vr x)
\end{align}
is also a solution of the system, namely $\innerprod {\vr A} \vr y  \geq  \vr t$.
We notice that for finitary $\vr x$, the vector $\vr y$ is finitary as well.
By the very definition, the averaging \eqref{eq:defy} preserves the orbit-sum:
$\orbsum{\vr x} = \orbsum{\vr y}$.
Furthermore, we claim that the vector $\vr y$ is setwise-$T$-supported.
To prove this, we fix an arbitrary $\pi \in \Aut {\set{T}}$, aiming at showing
that $\pi(\vr y) = \vr y$.
It factors through
$
\pi = \sigma \circ \rho
$
for some $\rho \in \Aut{\A\setminus T}$ and $\sigma \in \Aut T$.
Indeed, $\rho$ acts as $\pi$ on $T$ but is identity elsewhere, while
$\sigma$ acts as $\pi$ outside of $T$ but is identity on $T$.
%
%\arka{identically on $T$ $\mapsto$ identity permutation inside $T$}
%
A crucial but simple observation is that, 
by the very construction of $\vr y$, we have 
%\slawek{emphasized}
\begin{align} \label{eq:1}
\rho(\vr y) = \vr y.
\end{align}
Indeed, as $\vr y$ is defined by averaging over all $\rho'\in\Aut{\A\setminus T}$, 
\[
\rho \Big(\sum_{\rho' \in \Aut{\A\setminus T}}\rho'(\vr x)\Big) \ =
\sum_{\rho' \in \Aut{\A\setminus T}}\rho\circ\rho'(\vr x) \ =
\sum_{\rho' \in \Aut{\A\setminus T}}\rho'(\vr x)
\]
which implies $\rho(\vr y) = \vr y$.
Moreover, as action of atom automorphisms commutes with support, we have
\[
\supp {\rho'(\vr x)} = \rho'(\supp{\vr x})
\]
%\arka{$\supp {\rho'(\vr x)} = \rho'(\supp{\vr x})$?}
for every $\rho' \in \Aut{}$, and therefore 
\[
\supp {\rho'(\vr x)} = \supp{\vr x}
\]
for every $\rho' \in \Aut{\A\setminus T}$. 
Therefore $T$ supports the right-hand side of \eqref{eq:defy}, which means that
$\supp {\vr y} \subseteq T$ and implies
%
%\arka{$\supp{\vr{y}} =
%\cup_{\pi \in \Aut{\A \setminus T}} \supp{\pi(\vr{x})} =
%\cup_{\pi \in \Aut{\A \setminus T}} \pi(\supp{\vr{x}}) =
%\cup_{\pi \in \Aut{\A \setminus T}} \pi(T) = T
%$?}
%
\begin{align} \label{eq:2}
\sigma(\vr y) = \vr y.
%\piotrek{\text{I think this is incorrect, you are not using how $y$ is constructed.}}
\end{align}
By \eqref{eq:1} and \eqref{eq:2} we obtain $\pi(\vr y) = \vr y$, as required.

%\slawek{added}

Finally, 
the equality $\orbsum {\vr x} = \orbsum {\vr y}$ follows directly by \eqref{eq:defy}.
%
%let $L : \GLin C \to \R$. % be an equivariant (partial) linear map.
%By  equivariance of $L$ we get $L(\rho(\vr x)) = L(\vr x)$
%for every $\rho \in \Aut{}$.
%By \eqref{eq:defy} and linearity of $L$ we obtain $L(\vr y) = L(\vr x)$.
\end{proof}

\begin{slexample}
\label{ex:Kirchoff-cont-cont}
Recall the system of inequalities from Examples \ref{ex:Kirchoff} and \ref{ex:Kirchoff-cont}.
Its finitary solutions correspond to finite directed graphs, whose vertices and edges are labeled
by real numbers satisfying constraints \eqref{eq:K1} and \eqref{eq:K2}.
According to Lemma \ref{lem:weaksol}, if such a directed graph existed,  
there would also exist a directed clique, where labels of all vertices are pairwise equal, 
and labels of all edges are pairwise equal as well, which satisfying constraints \eqref{eq:K1} and \eqref{eq:K2}.
In particular, all edges incoming to a vertex would carry the same value as all outgoing edges.
This requirement is clearly contradictory with constraints \eqref{eq:K1} and \eqref{eq:K2},
and hence the system has no finitary solutions.
\end{slexample}

%\arka{I like the description. Maybe we can also add some equations to facilitate reading}
%
%\piotrek{I dont like fragment below. It looks like something not related. I would like to write that later we will need a fact that existence of a single setwise-supported solution implies exsitence of semisupported solutions for all supports bigger than $S$. We capture this fact in the following lemma. I agree with Arka that more general statement of the corollary is better.}
%
%We may apply Lemma \ref{lem:weaksol} to a special case, when
%the solution $\vr x$ is setwise-supported by a set $S$
%obtained by removing one atom from $T$
%($\vr x$ is clearly $T$-supported, but not setwise-$T$-supported),
%in order to derive the following fact useful later:

In the next section we rely on the fact that
existence of a setwise-$S$-supported solution implies existence 
of such a solution for any support larger than $S$.
The fact follows immediately from
Lemma \ref{lem:weaksol}, since every setwise-$S$-supported vector is trivially $T$-supported, for
every superset $T$ of $S$:
%
\begin{corollary} \label{cor:weaksol}
If an equivariant system of inequalities $(\vr A, \vr t)$ has a finitary setwise-$S$-supported solution $\vr x$,
then  for every superset $T$ of $S$ of size $\size{T} = \size S +1$, 
the system $(\vr A, \vr t)$ has a finitary setwise-$T$-supported solution $\vr y$ 
such that $\orbsum {\vr x} = \orbsum {\vr y}$.
\end{corollary}
%
%\begin{proof}
%The averaging construction \eqref{eq:defy} in 
%the proof of Lemma \ref{lem:weaksol},
%applied to $T = S \cup \{\a\}$ for any fresh atom $\a\in\A$, 
%yields a finitary solution $\vr y$ setwise-supported by $T$, and satisfying
%$
%\sum \vr x = \sum \vr y,
%$
%as required.
%\end{proof}


% !TEX root = lin-prog-main.tex


\section{Decidability of real solvability}% orbit-finite linear programming} 
\label{sec:lp-decid-proof}

In this section
we prove Theorem~\ref{thm:lp-decid} by a reduction of \finineqsolvname$(\R)$ to
\polyineqsolvname
(cf.~Example \ref{ex:poly} in Section \ref{sec:intro}).


\subsection{Preliminaries}
\label{sec:prelimproof}

Consider an orbit-finite system of inequalities given by a matrix $\vr A :  B\times C \tofs \Z$ and
a target vector $\vr t : B\tofs \Z$. 
%
%\arka{Replace $\F$ with $\R$?} \slawek{$\Z$}
%

\begin{lemma} \label{lem:wlog}
W.l.o.g.~we can assume that $B$ and $C$ are disjoint unions of equivariant orbits $\otu \A k$, $k\in\Nat$:
\begin{align} \label{eq:eqorbits}
\begin{aligned}
B =  
\otu {\A} {n_1}  \uplus  \ldots \uplus  \otu {\A} {n_s} \qquad\qquad 
C =  \otu {\A} {m_1} \uplus \ldots \uplus \otu {\A} {m_r}
\end{aligned}
\end{align}
(see the figure below),
and that $\vr A$ and $\vr t$ are equivariant. The size blow-up is exponential in
atom dimension, but polynomial when atom dimension is fixed.
\end{lemma}
%
%\TODO{$o_i \mapsto n_i$}
%\TODO{$O \mapsto U$}
%\TODO{$q\mapsto s$}
%\arka{$q$ is also used to denote the polynomials. Maybe we should use some other letter here?}
\vspace{-4mm}
\begin{align*}
& \qquad\  \otu \A {m_1} \, \otu \A {m_2} \cdots \, \otu \A {m_r} \\
\vr A \ = \ \ \ &
\begin{matrix}
\otu \A {n_1} \\ \otu \A {n_2} \\ \cdots \\ \otu \A {n_s}
\end{matrix}
\left[
\begin{matrix}
\ \ \ \ \ \ \ \\ \hline
\\ \hline
\\ \hline
\ \ 
\end{matrix}
\right\rvert
\hspace{-1.25mm}
\left\rvert
\begin{matrix}
\ \ \ \ \ \ \  \\ \hline
\\ \hline
\\ \hline
\ \ 
\end{matrix}
\right\rvert
\hspace{-1.25mm}
\left\rvert
\begin{matrix}
\ \ \ \ \ \ \ \\ \hline
\\ \hline
\\ \hline
\ \ 
\end{matrix}
\right\rvert
\hspace{-1.25mm}
\left\rvert
\begin{matrix}
\ \ \ \ \ \ \ \\ \hline
\\ \hline
\\ \hline
\ \ 
\end{matrix}
\right]
\qquad
\vr t \ = \ \begin{bmatrix}
\ \ \\ \hline
\\ \hline
\\ \hline
\ \ 
\end{bmatrix}
\end{align*}
\vspace{1mm}

\noindent
(The proof is in Section \ref{sec:lp-decid-proof-proofs}.)
Note that this includes the case of finite systems, namely
$n_1 = \ldots = n_s = m_1 = \ldots = m_r = 0$.

\subsection{Idea of the reduction}
%
%Before proceeding with the reduction, we outline its main underlying ideas.
Suppose only finitary \emph{$T$-supported} solutions are sought, for a fixed $T\subseteqfin\A$.
%Tuples $a\in T^*$ we briefly call \emph{$T$-tuples}.
%\arka{Finite $T$-orbits are singleton.}
\finineqsolvname$(\R)$ reduces then to a finite system of inequalities $(\vr A', \vr t')$ obtained
from $(\vr A, \vr t)$ as follows:
%
\begin{enumerate}
\item[(1)] Keep only columns indexed by $T$-tuples (= elements of \emph{finite} $T$-orbits) $c\in C$, discarding all other columns.
\item[(2)]
Pick arbitrary representatives of \emph{all} $T$-orbits included in $B$, and  keep only rows of $\vr A$
and entries of $\vr t$ indexed by the representatives, discarding all others.
% rows.
%\item[(3)] Likewise proceed with the target $\vr t$.
\end{enumerate}
%
The system $(\vr A', \vr t')$ is solvable if and only if the original one $(\vr A, \vr t)$ has
a finitary $T$-supported solution.
Indeed, discarding unknowns as in (1) is justified as a finitary $T$-supported solution of $(\vr A, \vr t)$ 
assigns $0$ to each non-$T$-tuple.
Discarding inequalities as in (2) is also justified.
Indeed, each inequality in the original system is obtained by applying some atom $T$-automorphism to an inequality
in $(\vr A', \vr t')$, while atom $T$-automorphisms preserve $T$-supported solutions of $(\vr A', \vr t')$,
which implies that every $T$-supported solutions of $(\vr A', \vr t')$ 
is also a solution of all inequalities in the original system.
%
%
% defined by the \emph{aggretate} finite $q\times r$ matrix
%$\vr A'$, and the aggregate finite target vector $\vr t' \in \R^q$, defined as follows.
%The value of $\vr A'(i, j)$ is obtained by summing the values $\vr A(b, c)$, where $c$ ranges over all elements
%of $\otu \A {m_j}$,
%and $b\in \otu \A {o_i}$ is chosen arbitrarily:
%\[
%\vr A' (i, j) = \sum_{c\in \otu \A {m_j}} \vr A(b, c).
%\]
%Likewise one defines $\vr t'$.
%The result of aggregation does not dependent on the choice of $b\in \otu \A {o_i}$, since $\vr A$ is equivariant,
%and is well-defined as long as 
%one restricts to \emph{finite}%
%\footnote{We consider finite orbits $\otu \A {m_j}$ for the sake of intuitive explanation, 
%knowing that each such finite orbit is necessarily a singleton ($m_j = 0)$.}
%orbits $\otu \A {m_j}$. 
%%the sum has only finitely-many non-zero elements. 
%%The latter is assured by dropping $S$-orbits $M_j\subseteq C$ not satisfying the condition for some $b\in B$.
%

The above reduction yields no algorithm yet, as
%The first main difficulty is that 
we do not know a priori any bound on size of $T$,
%(cf.~Example~\ref{ex:nofin1}), 
and
% and~\ref{ex:nofin}), and 
the size of $(\vr A', \vr t')$ depends on the number of $T$-orbits and hence grows unboundedly when  $T$ grows.
We overcome this difficulty by using setwise-$T$-orbits instead of $T$-orbits, and relying on Lemmas \ref{lem:weaksol}
and \ref{lem:STorbit}.
The latter one guarantees that the number of setwise-$T$-orbits is constant - independent of $T$.
Once we additionally merge (sum up) all columns indexed by elements of the same setwise-$T$-orbit, we get $\vr A'$ of size independent of $T$.

%\arka{merge $\mapsto$ sum up?}

This still does not yield an algorithm, as entries of $\vr A'$ change
%A second  difficulty is that 
when  $T$ grows.
We however crucially discover %(presented  in Section~\ref{sec:red2poly})
%that allows us to alleviate this problem is 
that the growth of the entries of $\vr A'$ is \emph{polynomial}
in $n = \size T$, for sufficiently large $n$.
Therefore, $\vr A'$ is a matrix of polynomials in one unknown $n$, and solvability of $(\vr A, \vr t)$ is equivalent to solvability
of $(\vr A', \vr t')$ for some value $n\in \Nat$.
As argued in Section \ref{sec:poly}, the latter solvability is decidable.
% expressible in real arithmetic even if $n$ ranges over integers and not reals.

%\arka{Maybe we should explain the above with an example}


\subsection{Reduction of \finineqsolvname$(\R)$ to \allpolyineqsolvname}
%Proof of Lemma~\ref{lem:red2poly}: reduction} 
\label{sec:red2poly}

Let us fix an equivariant system $(\vr A, \vr t)$.
We construct a finite system $P_2$ of polynomially-parametrised 
inequalities such that $(\vr A, \vr t)$ has a finitary solution
 if and only if $P_2(n)$ has a solution for almost all $n\in\Nat$.

Let us denote by $d = \max\set{n_1, \ldots, n_s, m_1, \ldots, m_r}$ the maximal atom dimension of
orbits included in $B$ and $C$.
%
%\piotrek{I think we don't need to assume $n\geq 2d$ here. It is used in Lemma~\ref{lem:poly} for the first time, so on the next page.} 
%\piotrek{Something is in a wrong order below. First we say that we construct a system $P(T) = (\vr A(T), \vr t(T))$ then we describe $B,C$ then we go back to 
%the definition of a system, then we make a comment on orbits of $B$ and $C$. I think if we reorder and first describe semi T orbits of B, C fully, and then we 
%start to define $P(T)$ it will be simpler to read.}
%We construct a finite system $P(T) = (\vr A(T), \vr t(T))$ of inequalities of polynomial size, as a function of $T$.
%Later we will argue that the system depends only on the size $n$ of $T$, and that the dependence
%is expressed by univariate polynomials.
%

Let $T\subseteqfin \A$ be an arbitrary finite subset of atoms. % of size $n = \size T \geq 2d$. 
Both $B$ and $C$ split into setwise-$T$-orbits, refining~\eqref{eq:eqorbits}:
\begin{align} \label{eq:BC}
B = B_1 \uplus \ldots \uplus B_N \qquad
C = C_1 \uplus \ldots \uplus C_{M'}.
\end{align}
Let $C_1, \ldots, C_M$ be the \emph{finite} setwise-$T$-orbits among $C_1, \ldots, C_{M'}$
(clearly, $M$ and $N$ may depend on $T$).
%
Importantly, by Lemma \ref{lem:STorbit}, $N$ and $M$ do not depend on $T$ as long as $\size T\geq d$.
In fact $M = r$, the number of orbits included in $C$, as
by Lemma \ref{lem:STorbit} we deduce:
%
\begin{lemma} \label{lem:Tk}
Assuming $\size T \geq \ell$, the equivariant orbit $\otu {\A} \ell$ includes exactly one finite setwise-$T$-orbit, namely
$\otu T \ell$.
\end{lemma}
%
Our reduction proceeds in two steps:  first, %, the most involved one, 
we derive a finite
polynomially-parametrised system $P_1$, and then
we transform it further to a monotonic system $P_2$.
Monotonicity of $P_2$ guarantees correctness of reduction.

\para{Step 1 (finite polynomially-parametrised system)}
Our construction is parametric in $T$.
Let $b_1, \ldots, b_N$ be arbitrarily chosen representatives of setwise-$T$-orbits included in $B$.
%
Given $\vr A$ and $\vr t$,
we define an $N{\times}M$ matrix $\vr A_1(T)$ and a vector $\vr t_1(T) \in \Int^N$ as follows:
%
\begin{enumerate}
\item[(1)] 
Pick columns of $\vr A(T)$ indexed
by elements of all finite setwise-$T$-orbits included in $C$, and discard other columns; this yields a matrix 
$\vr A'(T)$ with finitely many columns (number thereof depending on $T$).
\item[(2)] 
Merge (sum up) columns of $\vr A'(T)$ indexed by elements of the same setwise-$T$-orbit; this yields a matrix $\vr A''(T)$ with $M$ columns ($M$ independent of $T$). 
\item [(3)]
Pick $N$ rows of $\vr A''$, indexed by $b_1, \ldots, b_N$,
and discard other rows; this yields an $N\times M$ matrix $\vr A_1(T)$.
\item[(4)] Likewise pick the corresponding entries of $\vr t$ and discard others, thus yielding a finite vector $\vr t_1(T)\in\Z^N$.
\end{enumerate}


%\begin{enumerate}
%\item [(1)]
%Pick $N$ rows of $\vr A$, indexed by $b_1, \ldots, b_N$,
%and discard other rows; this yields a matrix $\vr A'(T)$ with $N$ rows.
%\item[(2)] Likewise pick the corresponding entries of $\vr t$ and discard others, thus yielding a finite vector $\vr t(T)\in\Z^N$.
%\item[(3)] 
%Pick columns of $\vr A'(T)$ indexed
%by elements of all finite setwise-$T$-orbits included in $C$, and discard other columns; this yields a finite matrix $\vr A''(T)$.
%\item[(4)] 
%Merge (sum up) columns of $\vr A''(T)$ indexed by elements of the same setwise-$T$-orbit; this yields an $N\times M$ matrix
%$\vr A(T)$. 
%\end{enumerate}

\noindent
For $b\in B$ and $C_j \subseteq C$, $j \in \setto M$, 
we write $\orbsum {\vr A}(b, C_j)$ for the finite sum ranging over elements of $C_j$:
\[
\orbsum {\vr A}(b, C_j) = \sum_{c\in C_j} \vr A(b,c),
\]
%The value of $\vr A(b, C_j)$ is well-defined exactly when $\vr A(b, c)\neq 0$ for only finitely many $c\in C_j$,
%and hence may be ill-defined.  
%An $S,T$-orbit $C_j$ we call \emph{useful} if $\sum \vr A(b, C_j)$ is well-defined for all $b\in B$.
%Let $C_1, \ldots, C_M$ be the useful $S,T$-orbits among $C_1, \ldots, C_{M'}$, for some $M\leq M'$. 
%
which allows us to formally define the $B{\times} M$ matrix $\vr A''(T)$, 
the $N{\times}M$ matrix $\vr A_1(T)$ and the vector $\vr t_1(T) \in \Int^N$:
\begin{align} \label{eq:defAt}
\vr A''(T) (b, j) = \orbsum {\vr A}(b, C_j) 
\qquad\qquad
\vr A_1(T) (i, j) = \vr A''(T)(b_i, j) = \orbsum {\vr A}(b_i, C_j) 
\qquad\qquad
\vr t_1(T) (i) = \vr t(b_i).
\end{align}

\begin{slexample} \label{ex:cont}
We explain how the system \eqref{eq:ppar} in Example \ref{ex:poly} in Section \ref{sec:intro}
is obtained from
the system \eqref{eq:matr} in Example \ref{ex:lp}.
%
Fix a non-empty $T \subseteqfin \A$. 
The set $\A$ includes just one finite setwise-$T$-orbit, namely $T$.
Therefore the matrix $\vr A'(T)$ has $\size T$ columns, $\vr A''(T)$ has just one column,
and the system $(\vr A_1(T), \vr t_1(T))$ has just one unknown.
Furthermore, the set $\A$ includes two setwise-$T$-orbits, the finite one $T$ plus the
infinite one $\A \setminus T$,
and therefore the system $(\vr A_1(T), \vr t_1(T))$ has two inequalities.
Pick arbitrary representatives of the setwise-$T$-orbits, $b_1 \in T$ and $b_2 \in (\A\setminus T)$.
We have
\begin{align*}
\vr A_1(T)(1,1) \  = \ \sum_{c\in T }\vr{A}(b_1, c) = |T| - 1 
\qquad\qquad
\vr A_1(T)(2,1) \  = \  \sum_{c\in T}\vr{A}(b_2, c) = |T|.
\end{align*}
Replacing $|T|$ with $n$ yields the system $(\vr A_1(T), \vr t_1(T))$: % \eqref{eq:ppar}.
\begin{align} \label{eq:pparnmon}
\begin{bmatrix}
n-1 \\
n
\end{bmatrix}
\cdot
x
\ \geq \ 
\begin{bmatrix}
1 \\
1
\end{bmatrix}
\end{align}
which happens to be monotonic.
In general, the system obtained so far needs not be monotonic,
but we will ensure monotonicity in the subsequent step.
%
\end{slexample}


The choice of representatives $b_i$ is irrelevant, and hence
$\vr A_1(T)$ and  $\vr t_1(T)$ are well defined, 
since rows of $\vr A''$ indexed by any two elements of $B$ belonging the same setwise-$T$-orbit
are equal, and likewise the corresponding entries of $\vr t$:
%due to: % the following lemma
%(proof in Appendix):
%
\begin{lemma}\label{lem:equalised row}
If $b, b' \in B$ are in the same setwise-$T$-orbit, then 
$\vr t(b) = \vr t(b')$ and
$\orbsum {\vr A}(b, C_j) = \orbsum {\vr A}(b', C_j)$
for every $j\in \setto M$.
\end{lemma}
%
\begin{proof}%[Proof of Lemma \ref{lem:equalised row}]
Let $\pi \in \Aut{(T)}$ be such that $\pi(b) = b'$.
As $\vr t$ is equivariant, it is necessarily constant on the whole equivariant orbit to which $b$ and $b'$ belong
(cf.~Lemma \ref{lem:1O}), 
and hence $\vr t(b') = \vr{t}(b)$.

For the second point fix $j\in \setto M$.
As $\vr A$ is equivariant, it is constant over the orbit included in $B\times C$ to which $(b,c)$ belongs, 
for every $c\in C$, and hence
$\vr A(b,c) = \vr A(\pi(b), \pi(c))$. This implies
\[
\sum_{c \in C_j} \vr{A}(b,c) = 
\sum_{c \in C_j} \vr A(\pi(b),\pi(c)) = 
\sum_{c \in C_j} \vr A(b',\pi(c)).
\]
Since $\pi$ is a setwise-$T$-automorphism, when restricted to
the setwise-$T$-orbit $C_j$ it is a bijection $C_j \to C_j$,
and hence the two sums below differ only by the order of summation and are thus equal:
\[
\sum_{c \in C_j} \vr A(b',\pi(c)) = 
\sum_{c \in C_j} \vr A(b',c).
\]
The two above equalities imply the claim, namely
$
\sum_{c \in C_j} \vr A(b,c) = 
\sum_{c \in C_j} \vr A(b',c).
$ 
\end{proof}
%

\begin{notation} \label{not:M}
Let $T\subseteqfin \A$.
Due to Lemma \ref{lem:Tk},
the set of finite setwise-$T$-orbits $\set{C_1, \ldots, C_M}$ included in $C$
is in bijection with 
the set $\orbits C = \set{\O_1, \ldots, \O_M}$ of equivariant orbits included in $C$.
W.l.o.g.~assume $C_j\subseteq \O_j$ for $j = 1\ldots M$.
Take any finitary setwise-$T$-supported vector $\vr x : C \tofin \R$.
It is non-zero only inside finite setwise-$T$-orbits $C_j$, which implies
\[
\orbsum {\vr x}(C_j) \ = \ \orbsum {\vr x}(\O_j)
\]
for $j=1\ldots M$  (cf.~Notation \ref{not:orbsum}).
Furthermore, $\vr x$ is constant inside each $C_j$, which allows us to write
$\orbval {\vr x}(C_j)$ (cf.~Notation \ref{not:orbval}).
For notational convenience we slightly relax Notations \ref{not:orbval} and \ref{not:orbsum}
from now on, and treat the orbit-value and orbit-sum vectors as $M$-tuples,
$
\orbval{\vr x}, \ \orbsum{\vr x} \ \in \ \R^M,
$
with obvious meaning $\orbval{\vr x}(j) = \orbval{\vr x}(C_j)$ and $\orbsum{\vr x}(j) = \orbsum{\vr x}(C_j)$.
We note the (obvious) relation between $\orbval{\vr x}$ and $\orbsum{\vr x}$:
\begin{align} \label{eq:overbefore}
\orbsum{\vr x}(j) \ = \ \size{C_j} \cdot \orbval{\vr x}(j).
\end{align}
\end{notation}

The following lemma, being a cornerstone of correctness of the whole reduction,
is now not difficult to prove:
%
\begin{lemma} \label{lem:givenT}
Let $\size T \geq d$ and $\vr x : C\tofin \R$ a finitary setwise-$T$-supported vector.
The following conditions are equivalent:
\begin{itemize}
\item 
$\vr x$ is solution of $(\vr A, \vr t)$;
\item
$\orbval {\vr x}$ is a solution of $P_1(T) = (\vr A_1(T), \vr t_1(T))$.
\end{itemize}
\end{lemma}
%
%\begin{lemma} \label{lem:someT}
%$(\vr A, \vr t)$ has a weakly $S$-supported nonnegative rational solution if, and only if
%$P(T) = (\vr A(T), \vr t(T))$ has a nonnegative rational solution for some $T\subseteqfin \A\setminus S$.
%\end{lemma}
%
\begin{proof}
Take any setwise-$T$-supported vector $\vr x : C\tofin \R$, and
let $\vr x'$ be the restriction of $\vr x$ to $C' = C_1\uplus\ldots\uplus C_M$.
We argue that the following four conditions are equivalent, which implies the claim:

\begin{enumerate}
\item 
$\vr x$ is solution of $(\vr A, \vr t)$;
\item 
$\vr x'$ is solution of $(\vr A'(T), \vr t)$;
\item 
$\orbval{\vr x}$ is solution of $(\vr A''(T), \vr t)$;
\item
$\orbval {\vr x}$ is a solution of $(\vr A_1(T), \vr t_1(T))$.
\end{enumerate}

First, as $\vr x$ is finitary,
we have $\vr x(c)=0$ for all $c\notin C'$, and hence 
$\vr{A}'(T) \cdot \vr{x}' = \vr{A} \cdot \vr{x}$.  
This implies equivalence of (1) and (2).
Second, as $\vr A''$ is obtained from $\vr A'$ by summing columns over a setwise-$T$-orbit 
where the vector $\vr x$, being setwise-$T$-supported, is constant,
we have $\vr{A}''(T) \cdot \orbval{\vr{x}} = \vr{A}'(T) \cdot \vr{x}'$.
This implies equivalence of (2) and (3).
Finally, (3) implies (4) as $(\vr A_1(T), \vr t_1(T)$ is obtained from $(\vr A''(T), \vr t)$ by removing 
inequalities.
For the reverse implication, we recall that Lemma \ref{lem:equalised row} shows that
$\vr A''(b, j) = \vr A''(b_i, j)$ and $\vr t(b) = \vr t(b_i)$ for every $i\in\setto N$ and $b\in B_i$, 
and therefore $\vr A''(T)$ contains the same inequalities as $\vr A_1(T)$.
In consequence, (4) implies (3).
%
%In one direction,
%% assuming existence of a finitary $T$-supported real solution of $(\vr A, \vr t)$,
%we take a finitary setwise-$T$-supported solution $\vr x : C \tofin \R$ of $(\vr A, \vr t)$.
%%(as guaranteed by Lemma \ref{lem:weaksol}).
%%
%As $\vr x$ is finitary,
%we have $\vr x(c)=0$ for all $c\notin C'=C_1\uplus\ldots\uplus C_M$.
%Therefore
%the restriction $\vr x'$ of $\vr x$ to $C'$ is a solution of $(\vr A'(T), \vr t)$, since
%\begin{align} \label{eq:AAA}
%\vr{A}'(T) \cdot \vr{x}' = \vr{A} \cdot \vr{x}.
%\end{align}
%Furthermore, $\vr x'$ is constant on every $C_j$.
%% (the value on elements of $C_j$ we denote as $\vr x''(C_j)$).
%Therefore, the vector 
%\begin{align} \label{eq:''}
%\orbval {\vr x} = \orbval{\vr x}'
%\end{align}
%%$\vr x : \setto M \to \R$ defined by $\vr x(j) = \vr x''(C_j)$
%is a solution of $(\vr A''(T), \vr t)$, since the definition \eqref{eq:defAt} implies
%\begin{align} \label{eq:AA}
%\vr{A}''(T) \cdot \orbval{\vr{x}} = \vr{A}'(T) \cdot \vr{x}'.
%\end{align}
%Finally, $\orbval{\vr x}$ is automatically a solution of the system $(\vr A(T), \vr t(T))$, 
%as the system is obtained by discarding inequalities.
%
%
%%\arka{We delete rows to get $\vr{A}'$. So this can not be true.}
%%\slawek{$\vr{A}''(T)$, $\vr{A}'(T)$ and $\vr{A}(T)$ have the same row index.}
%%\arka{No. Look at point (1) of the reduction}
%%
%%\arka{Should we add the fact that
%%\[
%%(\vr{A}''(T) \cdot \vr{x}')(i) = (\vr{A} \cdot \vr{x})(b)
%%\]
%%for any $i \in \{1,\dots,N\}$ and $b \in B_i$?
%%And the same for last paragraph as well.
%%}
%
%In the opposite direction, take any setwise-$T$-supported vector $\vr x : C\tofin \R$ such that 
%$\orbval{\vr x} : \setto M \to \R$ is a solution of $(\vr A(T), \vr t(T))$.
%As $\vr x$ is setwise-$T$-supported, it necessarily
%maps each element of each finite setwise-$T$-orbit $C_j$ to $\orbval{\vr x}(j)$.
%%By defininion, this vector is setwise-$T$-supported. % and hence also $T$-supported; 
%In order to prove that $\vr x$ is also a solution of $(\vr A''(T), \vr t)$ we recall that each inequality
%in this system is obtained by applying some setwise-$T$-automorphism $\pi$
%%\red{an automorphism $\pi$ in the setwise stabilizer of $T$}  
%to some inequality $e$ in $(\vr A'(T), \vr t(T))$.
%As $\vr x$ solves $e$, $\pi(\vr x)$ solves $\pi(e)$ by equivariance of $(\vr A, \vr t)$.
%Finally, as $\pi(\vr x) = \vr x$, we deduce that $\vr x'$ solves $\pi(e)$.
%As $\pi\in\Aut{\set T}$
%%\red{$\pi$, the element of the setwise stabilizer of $T$,} 
%was chosen arbitrarily, $\vr x$ is a solution of $(\vr A, \vr t)$.
%
%The equalities
%\eqref{eq:AAA}, \eqref{eq:''} and \eqref{eq:AA} hold again,
%hence $\vr x$ and is a solution of $(\vr A'(T), \vr t(T))$.
\end{proof}
%
The function $T\mapsto P_1(T)$ is equivariant, i.e., invariant under action of atom automorphisms.
In consequence, the entries of $\vr A_1(T)$ and $\vr t_1(T)$ do not depend on the set $T$ itself, but only
on its size $\size T$. 
Indeed, if $\size T = \size{T'}$ then $\pi(T) = T'$ for some atom automorphism $\pi$, and hence
$\pi(P_1(T)) = P_1(T')$.
Since the system $P_1(T)$ is atom-less 
%(\arka{Maybe here we should recall the atom-less representation of semi $T$-orbits?})
we have also $\pi(P_1(T)) = P_1(T)$,
which implies 
$ P_1(T) = P_1(T'). $
%\[ (\vr A(T), \vr t(T)) = (\vr A(T'), \vr t(T')). \]
We may thus meaningfully write $P_1(\size T) = (\vr A(\size T), \vr t(\size T))$, i.e.,
$P_1(n) = (\vr A_1(n), \vr t_1(n))$ for $n\in\Nat$
 (cf.~Example \ref{ex:cont}).
% instead of $P(T)$, where $n=\size T$.
%

%The proof of Lemma \ref{lem:givenT} actually demonstrates the 
%more informative version thereof, formulated here for future reference:
%\begin{corollary} \label{cor:givenT}
%Assume $\size T = n \geq d$.
%Let $\vr x : C\tofin \R$ be setwise-$T$-supported, and $\vr y \in \R^M$ such that
%\begin{align} \label{eq:givenT}
%\vr y(j) \ = \ \vr x(C_j)
%\end{align}
%for $j = 1\ldots M$.
%Then $\vr x$ is a solution of $(\vr A, \vr t)$ if and only if $\vr y$ is a solution of $P(n)$.
%\end{corollary}

\smallskip

We argue that the dependence on $\size T$ is polynomial, as long as $\size T\geq 2d$:
%
\begin{lemma} \label{lem:poly}
There are univariate polynomials $p_{ij}(n) \in \Int[n]$ 
%with integer coefficients 
such that
$\vr A_1(n)(i,j) = p_{ij}(n)$ for $n \geq 2d$.
\end{lemma}
%
%\arka{Replace $\Int$ with $\Rat$ or $\R$?} 
%
\begin{proof} 
Let $n=\size T$.
Fix a setwise-$T$-orbits $B_i\subseteq B$ and a finite setwise-$T$-orbit $C_j\subseteq C$.
Each of them is included in a unique equivariant orbit, say:
\[
B_i \subseteq B' =  \otu {\A} p \qquad
C_j \subseteq C' = \otu {\A} \ell
\] 
(cf.~the partitions~\eqref{eq:eqorbits}).
Recall Lemma~\ref{lem:STorbit}:  $B_i$ is determined by the subset 
$I \subseteq \setto p$ of positions where atoms of $T$ appear in tuples belonging to $B_i$.
Let $m = \size{I}$.
On the other hand $C_j = \otu T \ell$ (cf.~Lemma~\ref{lem:Tk}).
% is induced by the subset $J\subseteq\setto \ell$ of positions where atoms of $T$
%appear in tuples belonging to $C_j$.
%
%Fix an arbitrary tuple $b\in B_i$.
%
%We know that $C_j$ is useful and hence the sum $\sum \vr A(b, C_j)$ is well-defined,
%i.e., 
%\begin{claim} \label{claim:nzeroC}
%The set $\setof{c\in C_j}{\vr A(b, c)\neq 0}$ is finite.
%\end{claim}
%
%Let $R = \supp {b_i}$ be the set of atoms appearing in $b_i$.
Note that $m = \size{T\cap \supp{b_i}}$.

\smallskip

We are going to demonstrate that the value $\orbsum{\vr A}(b_i, C_j)$ is
polynomially depending on $n=\size T$.
%\piotrek{$q_{v,w}$ are in wikipedia 
%''In a similar manner, the number of arrangements of k items from n objects is sometimes called a partial permutation or a k-permutation. It can be written as 
%$nPk$ (which reads "n permute k"), and is equal to the number $n ( n - 1 )\ldots ( n - k + 1 )$ (also written as $n ! / ( n - k ) !$.[11][12]''
%}
%\piotrek{I am not sure if we should use the name $k$-permutations or not. To think about it.}
We will use the polynomials 
$\otun n w$
of degree $w$, for $w\leq d$, defined by 
\begin{align} \label{eq:poly2}
\begin{aligned}
\otun n w & = n \cdot (n-1) \cdot \ldots \cdot (n-w+1).
%q_{v,w}(n) & = (n{-}v) \cdot (n{-}v{-}1) \cdot \ldots \cdot (n{-}v{-}w{+}1).
%\\
%& = \frac {(n-v)!} {(n-v-w)!} % = p(n).
\end{aligned}
\end{align}
In the special case of $w=0$, we put $\otun n w = 1$.
The value $\otun n w$ can be interpreted as follows:
%
\begin{claim} \label{claim:polyint}
For $n\geq w$, $\otun n w$ is equal to the number of 
arrangements of $w$ items chosen from  $n$ objects into a sequence.
%possible ways of filling $w$ slots with pairwise distinct element chosen from a collection of size $n-v$.
\end{claim}

Denote by $\mathcal D$ the set of equivariant orbits $\O \subseteq B' \times C'$.
% that contain at least one pair $(b_i, c)$, for some $c\in C_j$, i.e., for which $\O(b_i)$ is nonempty.
For $\O\in \mathcal D$, we put $\O(b_i, C_j) := \setof{c\in C_j}{(b_i,c)\in \O}$.
%
%\piotrek{Do we need that in $\mathcal D$ are only orbits for which $\O(b)$ is nonempty? Maybe we can 
%use all orbits then it is simpler to parse i think.}
%\slawek{I suppose it is simpler as it is now}
%
As $\vr A$ is equivariant, the value $\vr A(b_i, c)$ depends only on the orbit to which $(b_i, c)$ belongs.
We write $\vr A(\O)$, for $\O\in\mathcal D$, and get:
%
\begin{claim} \label{claim:sum}
$\orbsum {\vr A}(b_i, C_j) = \sum_{\O\in \mathcal D} \vr A(\O) \cdot \size {\O(b_i, C_j)}$.
\end{claim}


By Lemma \ref{lem:orbitpair} in Section \ref{sec:atoms}, orbits $\O\subseteq B' \times C'$ are in 
one-to-one correspondence
with partial injections $\iota : \setto p \to \setto \ell$.
We write $\O_\iota$ for the orbit corresponding to $\iota$.
%Consequently, orbits $\O\in\mathcal D$ are in
%one-to-one correspondence
%with partial injections $\iota$ satisfying
%$\dom {\iota} \subseteq I$, where 
Let $\dom \iota = \setof{x}{\iota(x) \text{ is defined}}$ denote the domain of $\iota$.
%
\begin{claim} \label{claim:Uiotaempty}
%$\mathcal D = \setof{\O_\iota}{\dom{\iota}\subseteq I}$.
$U_\iota(b_i, C_j)\neq \emptyset$ if and only if $\dom{\iota}\subseteq I$. 
\end{claim}
%
\noindent
Indeed, recall again Lemma \ref{lem:orbitpair}  which yields
$
\O_\iota(b_i, C_j) = \setof{ c \in C_j }{\prettyforall{x, y}{b_i(x) = c(y) \iff \iota(x) = y}}.
$
If $\dom{\iota}\subseteq I$, the set $\O_\iota(b_i, C_j)$ contains tuples $c\in C_j$ with fixed values
on positions $J = \setof{\iota(x)}{x\in\dom \iota}$, namely
\begin{align} \label{eq:iotam}
b_i(x) = c(\iota(x)),
\end{align}
and arbitrary other atoms from $T$ elsewhere, and therefore is nonempty.
If there is $x \in \dom{\iota}\setminus I$ then $b_i(x)\notin T$ and therefore
no $c\in C_j$ satisfies \eqref{eq:iotam}.
Claim \ref{claim:Uiotaempty} is thus proved.
%
\begin{claim} \label{claim:polOC}
Let $k = \size{\dom{\iota}}$ be
the number of pairs related by $\iota$.
If $\O_\iota(b_i, C_j)\neq \emptyset$ then $\size {\O_\iota(b_i, C_j)}  = \otun {(n-m)} {\ell - k}$. %  $q_{m, \ell-k}(n)$.
\end{claim}
%
\noindent
According to \eqref{eq:iotam}, tuples $c\in \O_\iota(b_i)$ have fixed values on $k$ positions in $J$.
The remaining $\ell - k$ positions in tuples $c\in \O_\iota(b_i,C_j)$ are filled arbitrarily
using $n-m$ atoms from $T\setminus \supp{b_i}$.
%
Due to the assumption that $n\geq 2d$, we have $n-m \geq d \geq  \ell - k$, and therefore
using Claim \ref{claim:polyint} (for $w=\ell-k$) we deduce 
$\size {\O_\iota(b_i,C_j)} = \otun {(n-m)} {\ell-k}$,  % $q_{m, \ell-k}(n)$, 
thus proving Claim \ref{claim:polOC}.

\smallskip

Once $b_i\in B_i$  and $\O\in\mathcal D$  
are fixed, the values $k, \ell$ and $m$ are fixed too, and
the formula of Claim~\ref{claim:polOC}
%~\eqref{eq:poly}
is an univariate polynomial of degree $\ell-k$.
%
%We observe, finally, that the definition of $p(n)$ does not depend on the choice of $O$ - we obtain
%the same polynomial $p(n)$ for each $(S\cup R)$-orbit $O$ such that $O\cap C_i$ is finite and nonempty.
%This yields the formula:
%\[
%p_{ij}(n) = \size {\setof{O}{O\cap C_i \text{ is finite and nonempty}}} \cdot p(n).
%\]
%
%we have thus shown:
%%
%\begin{claim}
%$\size{O \cap C_j}$ is polynomial in $n = \size T$ but is independent from $T$ itself.
%\end{claim}
%
%
The formula of  Claim~\ref{claim:sum} yields the required polynomial%
\footnote{
This confirms, in particular,  that $\vr A(T)(i,j)$ is independent from the actual set $T$, 
and only depend on its size $n=\size T$.}
 $\vr A(n)(i,j) = p_{ij}(n)$ and hence the proof of Lemma~\ref{lem:poly} is completed.
\end{proof}
%

Relying on Lemma \ref{lem:poly} we get a polynomially-parametrised system 
$P_1(n) = (\vr A_1(n), \vr t_1(n))$.

\para{Step 2 (monotonicity)}
The system $P_1$ constructed so far, does \emph{not} have to be monotonic in general.
As an immediate corollary of Lemma~\ref{lem:givenT} and Corollary \ref{cor:weaksol}, we only know that
If $P_1(n)$ has a solution for $n\geq d$, then
$P_1(n+1)$ has a (potentially different) solution.
We slightly modify the system $P_1$
in order to achieve monotonicity.

Before defining formally the new system 
$P_2(n) = ({\vr A_2}(n), {\vr t_2}(n))$,
we point to our objective: we aim at replacing the orbit-value vector $\orbval {\vr x}$ 
in Lemma \ref{lem:givenT} by the orbit-sum vector
$\orbsum {\vr x}$, as in Lemma \ref{lem:givenTT} below.
In other terms, we want the solutions $\vr y_1$ and ${\vr y_2}$ of $P_1(n)$ and
$P_2(n)$, respectively, 
%
%For $j = 1 \ldots M$ let $E_j$ be the equivariant orbit including $C_j$.
%Let $T\subseteqfin \A$ and $n=\size T$.
%Recall Corollary \ref{cor:givenT} determining a solution $\vr y$ of $P(n)$
%from a solution $\vr x$ of $(\vr A, \vr t)$.
%%that, 
%%according to the proof of Lemma \ref{lem:givenT}, given a finitary setwise-$T$-supported 
%%solution $\vr x : C\tofin \R$ of $(\vr A, \vr t)$ 
%%(which is by definition constant over every setwise-$T$-orbit $C_j$),
%%one obtains a solution $\vr y$ of $P(n)$ by 
%%%
%%\begin{align} \label{eq:y}
%%\vr y(j) = \vr x(C_j) = \vr x(c) \quad (\text{for any } c\in C_j)
%%\end{align}
%%%
%%for $j = 1 \ldots M$.
%Instead, a solution of the modified system $\overline P$ will be obtained by
%summing $\vr x(c)$ for all $c\in C_j$:
%%
%\begin{align} \label{eq:overy}
%\overline{\vr y}(j) = \sum \vr x(C_j) = \sum_{c\in C_j} \vr x(c)
%\end{align}
%%
%(recall the orbit-summation mapping $\sum \vr x$ from Section \ref{sec:weak}).
%Therefore, the values
%$\vr y(j)$ and $\overline{\vr y}(j)$ 
%%
differ on position $j$ by the multiplicative factor of $\size{C_j}$, 
namely
%
\begin{align} \label{eq:over}
{\vr y_2}(j) \ = \ \size{C_j} \cdot \vr y_1(j)
\end{align}
%
for $j = 1, \ldots, M$ (cf.~\eqref{eq:overbefore}).
The size $\size {C_j}$ of the setwise-$T$-orbit $C_j$, where $\size T = n$, is equal to
\begin{align} \label{eq:orbitsize}
\size{C_j} \ = \ \otun n {e_j},   % q_{0, e_j}(n),
\end{align}
where $C_j \subseteq \otu \A {e_j}$, i.e., 
$e_j$ is the atom dimension of the equivariant orbit including $C_j$, assuming
$\size T \geq e_j$.
These considerations lead to the following formal definition of $P_2$:
\begin{align} \label{eq:defP}
\begin{aligned}
{\vr A_2}(n)(i,j) \ & = \  \vr A_1(n)(i,j) \cdot
\frac{\otun n d}{\otun n {e_j}} 
% \frac{q_{0,d}\,(n)}{q_{0,e_j}(n)}
\qquad\qquad\qquad
%
{\vr t_2}(i) \ & = \ \vr t_1(i) \cdot \otun n d  % q_{0,d}(n)
\end{aligned}
\end{align}
where $\vr A_1(n)(i,j) = p_{ij}(n)$.
%
We rely on the following fact:
%
\begin{claim} \label{claim:qeq}
$\otun n w \cdot \otun {(n-w)} u \ = \ \otun n {w+u}$. 
%$q_{0, w}(n) \cdot q_{w, u}(n) \ = \ q_{0, w+u}(n)$.
\end{claim}
%
\noindent
By the claim, all coefficients in \eqref{eq:defP} are polynomials, namely:
$
{\vr A_2}(n)(i,j) \ = \  p_{ij}(n) \cdot \otun {(n-e_j)} {d-e_j},  % q_{e_j,d-e_j}(n),
$
since $e_j \leq d$.
It remains to conclude that the systems $P_1(n)$ and $P_2(n)$ have the same 
solutions modulo \eqref{eq:over}:
%
\begin{lemma} \label{lem:modulo}
Let $n\geq d$ and let 
$\vr y_1, {\vr y_2}\in\R^M$ satisfy 
$
{\vr y_2}(j) \ = \  \otun n {e_j} % q_{0, e_j}(n) 
\cdot \vr y_1(j)
$
for $j = 1, \ldots, M$.
%\eqref{eq:over}.
Then
$\vr y_1$ is a solution of $P_1(n)$ if and only if \xspace ${\vr y_2}$ is a solution of $P_2(n)$.
\end{lemma}
%
Combination of \eqref{eq:overbefore}, \eqref{eq:orbitsize},
and Lemmas \ref{lem:givenT} and \ref{lem:modulo} yields:
\begin{lemma} \label{lem:givenTT}
Let $\size T = n \geq d$ and $\vr x : C\tofin \R$ a finitary setwise-$T$-supported vector.
The following conditions are equivalent:
\begin{itemize}
\item 
$\vr x$ is solution of $(\vr A, \vr t)$;
\item
$\orbsum {\vr x}$ is a solution of $P_2(n)$.
\end{itemize}
\end{lemma}
%
\begin{slexample} \label{ex:contcont}
The system \eqref{eq:ppar} in Example \ref{ex:poly} is obtained from \eqref{eq:pparnmon}
in Example \ref{ex:cont} by applying the definition \eqref{eq:defP}.
Indeed, $M = d = e_1 = 1$ and hence $\vr A$ stays unchanged, while
the right-hand side vector $\vr t$ gets multiplied by $\otun n 1  = n$.
\end{slexample}

%By Lemmas \ref{lem:weaksol}, \ref{lem:givenT} and \ref{lem:modulo},
%for every $T\subseteqfin \A$ of size $\size T = n \geq 2d$, the following conditions are equivalent:
%\begin{enumerate}
%\item
%$(\vr A, \vr t)$ has a  finitary $T$-supported solution;
%\item
%$(\vr A, \vr t)$ has a  finitary setwise-$T$-supported solution;
%\item 
%$\overline P(n)$ has a solution.
%\end{enumerate}
%%
%(The constraint $n\geq 2d$ is inherited from assumption in Lemma \ref{lem:poly}.)
%%
%Relying on equivalence of conditions 2) and 3) above, 
%together with Corollary \ref{cor:weaksol},
%we deduce:

\begin{lemma}[Monotonicity] \label{lem:monoto}
Let $n\geq d$. Every solution of $P_2(n)$ is also a solution 
of $P_2(n+1)$.
\end{lemma}
%
\begin{proof}
Suppose ${\vr y}$ is a solution of $P_2(n)$, for $n\geq d$.
Let $T\subseteqfin\A$ be any subset of atoms of size $\size T = n$.
Let $\vr x$ be the finitary setwise-$T$-supported vector uniquely determined by
%
\begin{align} \label{eq:overyone}
\orbsum {\vr x} \ = \  {\vr y}.
\end{align}
%
By Lemma \ref{lem:givenTT}, $\vr x$ is a solution of $(\vr A, \vr t)$.
We apply Corollary \ref{cor:weaksol} to obtain another finitary solution 
$\vr x'$ of $(\vr A, \vr t)$, 
setwise-$T'$-supported by some $T'$ of size $\size T' = n+1$,
and having the same orbit-summation mapping:
%
\begin{align}\label{eq:sums}
\orbsum {\vr x} \ = \  \orbsum {(\vr x')}.
\end{align}
%
Equalities \eqref{eq:overyone} and \eqref{eq:sums} imply $ \orbsum {(\vr x')} = {\vr y}$.
By Lemma \ref{lem:givenTT} again, $ \orbsum {(\vr x')} = {\vr y}$
is a solution of $P_2(n+1)$, as required.
\end{proof}

%
%(restricting to $n\geq 2d$ is necessary as this is assumed in Lemma \ref{lem:poly}, but it is void due to monotonicity).
%
%\arka{We only get monotonicity for $n \geq 2d$. So the above statement should be omitted.}
%

%\para{Step 3 (monotonicity)}

Combining Lemmas \ref{lem:weaksol}, \ref{lem:poly}, \ref{lem:givenTT} and \ref{lem:monoto}
we derive correctness of reduction
(the constraint $n\geq 2d$ is inherited from the assumption in Lemma \ref{lem:poly}):
%
\begin{corollary} \label{cor:iff}
The following conditions are equivalent:
\begin{itemize}
\item
$(\vr A, \vr t)$ has a finitary solution,
\item
$P_2(n)$ has a solution for some integer $n \geq 2d$,
\item
$P_2(n)$ has a solution for almost all $n \in \Nat$.
\end{itemize}
\end{corollary}
%
%
%In order to drop the constraint, 
%we substitute $n+2d$ in place of $n$ to all polynomials appearing in 
%$P_2$.
%This substitution yields
%the final system $P_3$ such that:
%\begin{claim} \label{cor:QP}
%$P_3(n)$ has the same solutions as $P_2(n+2d)$  for  every $n\in\Nat$. 
%\end{claim}
%%
%Therefore Corollary \ref{cor:fin} and Lemma \ref{lem:overP}, respectively, rewrite to:
%%
%\begin{corollary} \label{cor:iff}
%$(\vr A, \vr t)$ has a  finitary solution if and only if
%$P_3(n)$ has a solution for some $n\in\Nat$.
%\end{corollary}
%
%\noindent
%Furthermore,
%Lemma \ref{lem:overP} implies:
% equivalence of conditions 2) and 3) above, together with Corollary \ref{cor:weaksol}, we deduce:
%
%\begin{corollary} \label{cor:monotonic}
%$P_3$ is monotonic.
%\end{corollary}
%

Reduction of \finineqsolvname$(\R)$ to \allpolyineqsolvname is thus completed.

\para{Complexity}
It remains to argue that $P_2$ is computable from $(\vr A, \vr t)$, and estimate the 
computational complexity.
Computability of $P_2$ follows immediately from computability of $P_1$, which we focus now on:
%
\begin{lemma} \label{lem:comp}
The system $P_1$ is computable from $(\vr A, \vr t)$. % in polynomial time.
\end{lemma}
\begin{proof}
Indeed, it is enough to range over representations of setwise-$T$-orbits $B_i$ and $C_j$ of $B$ and $C$, respectively
(such representations are given by Lemma~\ref{lem:STorbit}), and for each
pair of such orbits proceed with computations outlined in the proof of Lemma~\ref{lem:poly}, applied
to an arbitrarily chosen representative $b_i\in B_i$. 
\end{proof}

By Corollary \ref{cor:iff} and Lemma \ref{lem:comp},
\finineqsolvname$(\R)$  reduces to \allpolyineqsolvname.


\smallskip

%

%\smallskip


%\begin{slremark}
%As a corollary of Lemma~\ref{lem:givenT}, together with the effective quantifier elimination of
%real arithmetic \cite{?}, we deduce the following bound on the smallest size of support of a solution: 
%whenever an orbit-finite system of linear inequalities of atom dimension $d$, supported by $S$,
%has a solution then it necessarily has a solution, whose support is 
%larger than $\size S$ by a quantity at most doubly exponential in $d$. 
%\end{slremark}

%\arka{I think this example should either be later, or should have more explanation.}
%\slawek{Move it to the end of sect?}


Concerning computational complexity,
the number of setwise-$T$-orbits included in an equivariant orbit $\otu \A \ell$ is exponential in $\ell$
(cf.~Lemma \ref{lem:STorbit}).
That is why 
the size of  $P_2$ may be exponential in atom dimension $d$ of $(\vr A, \vr t)$.
On the other hand, the size of $P_2$ is only polynomial
(actually, linear) in the number of orbits included in $B\times C$.
In consequence, for fixed atom dimension we get a polynomial-time reduction and hence,
relying on Theorem \ref{thm:polyptime},
the decision procedure for \finineqsolvname$(\R)$ in \PTIME.
Without fixing atom dimension, we get an exponential-time reduction and hence
the decision procedure is in \EXPTIME.
%\smallskip

The same complexity bounds apply to the algorithm for the optimisation problem presented in Section \ref{sec:max}.




%\input{lin_prog_complicated_example}

% !TEX root = lin-prog-main.tex

\section{Optimisation problems}
\label{sec:max}


In this section we prove Theorem~\ref{thm:lp-max}:
we introduce a maximisation variant of \polyineqsolvname and routinely adapt
the decision procedure of Section \ref{sec:poly}, as well as
the reduction of Section \ref{sec:red2poly},
to the maximisation setting.



%\compprobext{\finineqsolvname$(\F)$}
%{an orbit-finite system of linear inequalities $(\vr A, \vr t)$,}{and a score function $S$}{}
%{the supremum of $\setofbeg{S(\vr x)}{\vr x \text{ solves } (\vr A, \vr t)}$}


\subsection{Polynomially-parametrised maximisation problem}

We consider a maximisation problem, whose instance 
$(P,S)$ consists of a finite system of polynomially-parametrised 
inequalities $P$ as in \eqref{eq:polyeq}, and an \emph{ordinary} (non-parametrised)
objective function $S$ given by a linear map
\begin{align*} %\label{eq:score}
S(x_1, \ldots, x_k) = 
a_1 \cdot x_1 + \ldots + a_k \cdot x_k.
\end{align*}
%
%$S(n,\vec x) \leq S(n+1, \vec x)$, for every $n\in\Nat$ and $\vec x\in\R^k$.
%
%The instance $(P, S)$ is \emph{monotonic} if the system $P$ is so.
%In this section we exclusively concentrate on such instances.
%
%
As in Section \ref{sec:poly-mon}, % (c.f.~Notation \ref{not:sol}),
by an \emph{\aasol} of a system $P$ we mean in this section a solution of $P(n)$ 
for almost all $n\in\Nat$.  % (all sufficiently large $n\in\Nat$).
%
We define the \emph{supremum} of a monotonic instance $(P, S)$ 
as
\[
\sup (P, S) \ := \ \sup \setof{S(\vr x)}{\vr x \text{ is an \aasol of } P},
\]
under a proviso that $\sup (P,S) = -\infty$ if $P$ has no {\aasol}s.
Referring to the standard terminology, we can say that the system is \emph{infeasible} if
$sup (P,S) = -\infty$, and it is \emph{unbounded} if $\sup (P, S) = \infty$.
%
%CONVENTION not needed; In the following we apply the convention that 
%$-\infty < r < +\infty$ for all $r\in\R$.
%
Interestingly, the supremum can not be irrational
(see Corollary \ref{cor:limQ} below).
% (the proof is in Appendix):
%


%For the sake of formulating a computation problem, we use the following representation of real numbers:
%a number $r\in\R$ is \emph{represented} by a real arithmetic formula $\varphi(x)$ with one free variable if
%the formula defines the singleton set $\set{r}$.%
%\footnote{
%Using quantifier elimination, the formula can be further transformed to the standard  representation
%as a pair $(p, i)$, where $p$ is an univariate polynomial and $i$ a positive integer, such that $r$ is the $i$th 
%real root of $p$.
%%\slawek{remove it?}\piotrek{Yes I think it is irrelevant for the content of the paper.}
%}

In this section we study the problem of computing the supremum of monotonic instances:

\compprobext{\allpolyineqmaxname} %\polyineqmaxname}
{An instance $(P, S)$}
%{A finite system of polynomially-parametrised}{inequalities $P$, a polynomially parametrized}{linear map $S$} 
{The supremum of $(P, S)$}

\noindent
The problem generalises ordinary (non-parametrised) linear programming, and can be solved
similarly to \allpolyineqsolvname (of which it is a strengthening):
%In the sequel we concentrate on monotonic instances, and argue that on such instances
%the problem is still in \PTIME, thus strengthening \PTIME complexity of monotonic
%\polyineqsolvname (cf.~Theorem \ref{thm:polyptime}):
%
\begin{theorem} \label{thm:lp-max-poly}
%For monotonic instances, 
\allpolyineqmaxname is in \PTIME.
\end{theorem}
%
\begin{proof}
Let $(P_0, S)$ be an instance.
The algorithm is essentially the same as Algorithm \ref{alg:1} for \allpolyineqsolvname 
%(that subsumes monotonic \polyineqsolvname) 
in the proof of Theorem \ref{thm:polyptime}, and proceeds
by iterating the transformation step until either unsolvability or solvability is reported.
Recall that solution set is preserved by the transformation step (Claim \ref{claim:samesol}).
If unsolvability is reported, the algorithm returns $-\infty$.
If solvability is reported --- let $P\cup \G$ be the system examined in the last iteration
---
the decision procedure computes and returns $\sup (\li P \conj \G, S)$, 
the supremum of $S$ constrained by the ordinary system of inequalities $\li  P \conj \G$,
by invoking any \PTIME procedure for ordinary linear programming.
%\cite{CohenLS19} \slawek{what should we cite?}). 

Correctness follows by the two claims formulated below. 
First, since solution set is preserved by the transformation step, we have:
%
\begin{claim} \label{claim:proof1}
$\sup (P_0, S) = \sup (P \cup \G, S)$.
\end{claim}
%
\noindent
Second, the supremum does not change if
the polynomially-parametrised constraints $P$ are replaced by the 
overapproximation $\li P$:
%
\begin{claim} \label{claim:proof2}
$\sup (P\cup\G, S) =  \sup (\li P \cup \G, S)$. 
\end{claim}
%
\noindent
%We prove both inequalities, using Claims \ref{claim:li} and \ref{claim:strictli}.
%
For the claim it is enough to prove the inequality
$
\sup (\strictli P\cup\G, S) \geq  \sup (\li P \cup \G, S)
$
as, according to Claims \ref{claim:li} and \ref{claim:strictli}, we have
$\sup (\strictli P \cup \G, S) \leq \sup (P\cup\G, S) \leq  \sup (\li P \cup \G, S)$.
Take any solution $\vr y$ of $\li P\cup\G$, and any solution 
$\vr x$ of $\strictli P\cup\G$ (we rely here on solvability of the latter system).
For every $k\in\Nat$, the vector
$\vr x_k = \frac{\vr x + k \vr y}{k+1}$ is a solution of $\strictli P\cup\G$, and $S(\vr x_k)$
tends to $S(\vr y)$ when $k$ tends to $\infty$.
Hence 
$
\sup (\strictli P\cup\G, S) \geq  \sup (\li P \cup \G, S),
$
as required.
%
%Indeed, one inequality $\sup (P\cup\G, S) \leq  \sup (\li P \cup \G, S)$ follows immediately by 
%Claim \ref{claim:li} from Section \ref{sec:poly}.
%For proving the opposite inequality we first deduce 
%$\sup (\strictli P\cup\G, S) \leq  \sup (P \cup \G, S)$, using Claim \ref{claim:strictli}.
%Furthermore, relying on solvability of $\strictli P\cup\G$ (as solvability of $P_0$ has been reported
%by the algorithm) we observe the equality
%$\sup (\strictli P\cup\G, S) =  \sup (\li P \cup \G, S)$.
%Therefore $\sup (\li P\cup\G, S) \leq  \sup (P \cup \G, S)$, as required.
%
%\slawek{is this clear enough?}
%\arka{No. We have to say that this is safe since $\strictli P \cup \G$ is solvable. Hence we are guranteed that the supremum of both systems are the same. Moreover, we can check if the supremum is actually achieved by enlarging the system by one equation $S(\vr{x}) = \sup$}
%\piotrek{Maybe we can write as follows: Let Set(Eq) be the set of solutions of Eq. Then
%$Set(\lipar{P}{>} \cup \G)\subseteq Set( P \cup \G)\subseteq Set(\li{P} \cup \G)$ the first containment is because of Lemma~\ref{claim:strictli} the second because of Lemma~\ref{claim:li}. Because of the choice for $P\cup \G$ we have that $Set(\lipar{P}{>} \cup \G)$ is not empty so $sup(\lipar{P}{>} \cup \G,S)= sup(\li{P}  \cup \G,S)$ for any $S$, and this  proves the claim.}
\end{proof}

By Claims \ref{claim:proof1} and \ref{claim:proof2} in the proof of Theorem \ref{thm:lp-max-poly},
the supremum of a monotonic instance, if not  $-\infty$ nor $\infty$,
is equal to the supremum of an ordinary linear program and hence is rational:
\begin{corollary} \label{cor:limQ}
The supremum of a monotonic instance $(P,S)$ belongs to $\Rat\cup\set{-\infty,+\infty}$. 
\end{corollary}

\begin{slremark}
As illustrated in Example \ref{ex:poly} in Section \ref{sec:intro}, 
the objective function $S$ may not achieve its supremum over the {\aasol}s of $P$.
Once the supremum $s\in \Rat$ is computed, 
one can easily check if $S$ achieves its supremum, by adding to the system an equation
$S(x_1, \ldots, x_k) = s$ and checking if the system is still solvable.
\end{slremark}

\subsection{Reduction of \finineqmaxname$(\R)$ to  \allpolyineqmaxname}

%\slawek{check this section, it is new}

We only sketch the reduction as it amounts to a slight adaptation of the reduction
 of Section \ref{sec:red2poly}.
The input of \finineqmaxname$(\R)$ consists of a system $(\vr A, \vr t)$ and 
an integer vector $\vr s : C \tofs \Z$ representing
the objective function
\[
S(\vr x) = \innerprod{\vr s}{\vr x},
\]
and we ask for the supremum of values $S(\vr x)$, for $\vr x$ ranging over finitary solutions of $(\vr A, \vr t)$.
This value we denote as $\sup(\vr A, \vr t, \vr s)$.
%
%\piotrek{This looks odd, maybe we can introduce this fact in the text, or claim it and then say that additionally to to Lemma \ref{lem:wlog} we have to prove following Lemma.}
In addition to Lemma \ref{lem:wlog} we show (the proof is in Section \ref{sec:lp-decid-proof-proofs}):
% (in Appendix):
%
\begin{lemma} \label{lem:wlogmax}
W.l.o.g.~we may assume that $\vr s$ is equivariant.
\end{lemma}
%
We proceed by adapting the reduction of Section \ref{sec:red2poly}:
given an instance $(\vr A, \vr t, \vr s)$ of \finineqmaxname$(\R)$ we compute a monotonic instance
 $(P_2, S')$ of \polyineqmaxname, where
the finite system $P_2(n) = (\vr A_2(n), \vr t_2(n))$ of polynomially-parametrised inequalities 
is exactly as in Section \ref{sec:red2poly}, and
the objective function is
%
\begin{align} \label{eq:newobj}
S'(x_1, \ldots, x_k) = a_1 \cdot x_1 + \ldots + a_M \cdot x_M,
\end{align}
%
where $a_j = \orbval{\vr s}(C_j)$ for $j=1\ldots M$
(recall Notations \ref{not:orbval} and \ref{not:M}).
More concisely, the vector $\vr a = a_1 \ldots a_M$ is defined as
$\vr a = \orbval {\vr s}$.
%\piotrek{$\orbval {\vr s}$ is computed using division i think so it will be rational function not a polynomial here?, am i right?, but then i am not sure if we can speak about monotonicity.}
%Monotonicity of the instance $(P_2, S')$ is immediate due to Corollary \ref{cor:monotonic}.
%
%In other terms,
%as follows.
%As $\vr s$ is equivariant, it necessarily assigns the same value $\vr s(C_j)$
%to all elements of $C_j$.
%Assuming $\size T = n\geq 2d$ ($n\geq d$ is sufficient here),
%the finite setwise-$T$-orbit  $C_j= \otu T \ell$ contains $q_{0, \ell}(n)$ elements, 
%which allows us to define $p_j(n)$ as
%\[
%p_j(n) =  \vr s(C_j) \cdot q_{0, \ell}(n)
%\]
%(cf.~Claims \ref{claim:sum} and \ref{claim:polOC}, and Example \ref{ex:cont} in Section \ref{sec:red2poly}).
%\begin{claim} %\label{lem:instmon}
%The instance $(Q, S')$ is monotonic.
%\end{claim}
%
%\begin{proof}%[Proof of Lemma \ref{lem:instmon}]
%%Indeed, due to Lemmas \ref{lem:weakobj} and Lemma \ref{lem:givenTobj}.
%%\piotrek{This is its proof, there is nothing more to be said.}
%%
%Indeed, suppose $T \subseteq T'$ differ by one element, 
%and let $\size T = n$ (therefore $\size {T'} = n+1$).
%%\piotrek{we don't need $n$ it is better to write $|T|$ i think}.
%By  Lemma \ref{lem:givenTobj} the supremum $s_n$ is achieved on a finitary $T$-supported solution, which is also
%$T'$-supported. By Lemma \ref{lem:givenTobj} again, $s_{n+1}$ is the maximal value achieved on a finitary $T'$-supported
%solution, which implies 
%%\piotrek{$s_n \leq s_{\size{T'}}$}
%$s_n \leq s_{n+1}$.
%\end{proof}
%
We apply Lemmas \ref{lem:weaksol} and \ref{lem:givenTT}
%and Corollary \ref{cor:QP}
to obtain:
%
\begin{lemma}  \label{lem:givenTobj}
$\sup (\vr A, \vr t, \vr s) = \sup (P_2, S')$.
\end{lemma}
%
\begin{proof}%[Proof of Lemma \ref{lem:givenTobj}]
Let $\vr x : C\tofin\R$. By equivariance of $S$ and the definition of $S'$ we have the equality
$
S(\vr x) \ = \ S'(\orbsum{\vr x}),
$
that is,
the value of the objective function $S(\vr x)$ depends only on the orbit-sum vector 
$\orbsum{\vr x} : \orbits C \to \R$.
As Lemmas \ref{lem:weaksol} and \ref{lem:givenTT} preserve orbit-sum, we deduce that
for every $T\subseteqfin\A$ of size $\size T =n \geq 2d$,
the values of $S$ on finitary $T$-supported solutions of $(\vr A, \vr t)$ are the same as
the values of $S'$ on solutions of $P_2(n)$.
By Lemma \ref{lem:monoto}, the solutions of $P_2(n)$ for some $n\geq 2d$ are exactly the same as
the {\aasol}s of $P_2$.
%Finally, by Corollary \ref{cor:QP} the values of $S'$ on solutions of $P_2(n+2d)$ are the same
%as the values of $S'$ on solutions of $P_3(n)$.
In consequence, the two suprema are equal.
%the supremum of all values of $S$ on finitary $T$-supported real solutions of 
%$(\vr A, \vr t)$ is the same as the supremum of values of $S'$ on solutions of $Q(n)$
%for $n\in\Nat$.
%
% $(\vr A, \vr t, \vr s)$ is the same as the supremum of $(Q, S')$.
%
%we get the claim.
%the value 
%
%each value of $S$ on a finitary $T$-supported solution of $(\vr A, \vr t)$
%is also achieved on some finitary setwise-$T$-supported solution.
%%\piotrek{We have changed Lemma \ref{lem:weaksol}  so it covers this.}
%%%\end{claim}
%%%
%%Indeed, the averaging construction of $\vr y$ from $\vr x$ in \eqref{eq:defy} in the proof of Lemma \ref{lem:weaksol} 
%%preserves the value of the objective function,
%%$S(\vr y) = S(\vr x)$, due to equivariance of $S$.
%This allows us to restrict to finitely setwise-supported solutions only.
%
%We observe that the 
%the mutual transformations between finitary setwise-$T$-supported real solutions $\vr x$ of $(\vr A, \vr t)$ 
%and real solutions $\vr x'$ of $P'(n)$, used in the proof of Lemma \ref{lem:givenT},
%preserve the value of the objective functions:
%$S(\vr x) = S'(\vr x')$.
\end{proof}
%
%\noindent
%%(Proofs in Appendix.)
%This completes the reduction of \finineqmaxname$(\R)$ to monotonic  \polyineqmaxname.


%\arka{I think we can write shortly here how this wraps up the proof of theorem \ref{thm:lp-max}} 



\begin{slexample}
To illustrate the reduction, consider the modification of the system in Example \ref{ex:Kirchoff}:
% where the constraint \eqref{eq:K2} is changed to:
%
\begin{align}  \label{eq:K22}
\sum_{\a \in \A} \a \ \geq \  1
\qquad\qquad\qquad\qquad
\sum_{\b \in \A} \a\b \ - \ \a \ - \ 2 \cdot \sum_{\b \in \A} \b\a  \ \geq \ 0 \qquad\qquad (\a \in\A).
\end{align}
%
It enforces, for each vertex $\a\in\A$, the sum of values assigned to all outgoing edges
to be larger than \emph{double} the sum of values assigned to all ingoing edges, plus the value assigned
to the vertex $\a$.
The indexing sets $B=\A\cup\set{*}$ and $C=\A\cup\otu \A 2$ %\eqref{eq:BCex} 
and the shape of the matrix \eqref{eq:Kirchoff-matr} are the same.
We identify the singletons $\set{*} = \otu \A 0$.
%
We consider maximisation of triple the sum of values assigned to edges:
$S(\vr x) = \innerprod {\vr s}{\vr x}$, where $\vr s = 3\cdot \constvr 1 {\otu \A 2}$, or
$
S(\vr x) \ = \ 3 \cdot
\sum_{\a\b\in\otu\A 2} \vr x (\a\b).
$ 

According to Lemma \ref{lem:Tk},
the set $C$ includes exactly 2 finite setwise-$T$-orbits, namely
$T\subseteq \A$ and $\otu T 2 \subseteq \otu \A 2$, and therefore
the system computed by the reduction has 2 unknowns, $x_1$ and $x_2$.
%
By Lemma \ref{lem:STorbit}, 
for any nonempty $T \subseteqfin \A$,
the set $B$ includes 3 setwise-$T$-orbits,
namely $T$, $\A\setminus T$ and $\set{*}$,
%
and therefore the system $P_1$ computed in the first step has 3 inequalities:
%
\begin{align}
\label{eq:ukl}
\begin{aligned}
-x_1 \ - \, (n-1) \cdot x_2 & \ \geq \  0 \qquad\qquad && (T) \\
0 & \ \geq \  0 && (\A\setminus T) \\
n\cdot x_1  & \ \geq \ 1 && (\set{*})
\end{aligned}
\end{align}
%
% as shown in \eqref{eq:ukl} and \eqref{eq:ukltwo} below.
%
%
For instance, the coefficient $-(n-1)$ in the first inequality arises as:
%
\begin{align*}
 \vr A(U_\text{out}) \cdot \size{U_\text{out}(\a, \otu T 2)} \ + \
\vr A(U_\text{in}) \cdot  \size{U_\text{in}(\a, \otu T 2)}  \  = \  
 1 \cdot (n-1) -2 \cdot (n-1) \ = \ -(n-1)
\end{align*}
%
(cf.~Claim \ref{claim:sum}),
for some arbitrary $\a\in T$  and the following two orbits included in $\A \times \otu \A 2$:
\begin{align*}
& U_\text{out} = \setof{\!(\a, \a\b)\!}{\!\b\neq\a\!},
\ \ 
U_\text{in} = \setof{\!(\a, \b\a)\!}{\!\b\neq\a\!}.
\end{align*}
%
%We omit $\vr A(U_3) \cdot \size{U_3(\a\b)}$ above, as $A(U_3) = 0$ for the orbit
%%
%\begin{align*}
%& U_3 = \setof{\!(\a\b,\c)\!}{\!\a\b\in\otu \A 2, \ \c\notin\set{\a,\b}\!}.
%\end{align*}
%
Likewise, the coefficient $n$ in the last inequality arises as
$
\vr A(U) \cdot \size{O(*, T)} = 1 \cdot n = n,
$
for the orbit
$
U = \set{*}\times\A.
$
According to \eqref{eq:defP},
the system $P_2$ is obtained from \eqref{eq:ukl} by multiplying all occurrences of
$x_1$ by $\otun {(n-1)} 1  = n-1$, and by multiplying all right-hand sides by
$\otun n 2  = n(n-1)$ (the trivial second inequality is omitted):
%
\begin{align}
\label{eq:ukltwo}
\begin{aligned}
-(n-1)\cdot x_1 \ - \, (n-1) \cdot x_2 & \ \geq \  0  \\
n(n-1) \cdot x_1  & \ \geq \ n(n-1)
\end{aligned}
\end{align}
%
Finally, the objective function produced by the reduction, as in \eqref{eq:newobj}, is
$
S'(x_1, x_2) \ = \ \vr s({\otu \A 2}) \cdot x_2 \ = \   3 \cdot x_2.
$
%\[
%S'(x_2, x_1) \ = \ S({\otu \A 2}) \cdot q_{0,2}(n) \cdot x_2 \ = \  3\cdot n(n-1) \cdot x_2
%\]
It achieves $-3$ as its supremum,
%  large values, for every $n> 1$ grows, 
as the system \eqref{eq:ukltwo} is equivalent to the ordinary system (its head):
\[
x_1 \geq 1 \qquad\qquad x_2 \leq - x_1.
\]
%are only restricted by $x_2 \leq 4(n-1) x_1 \leq 4$.
%
%where $n = |T|$.
%Notice that the objective function becomes $\frac{n(n-1)x}{2}$.
%For every $n \geq 1$ clearly we can achieve the value $\frac{(n-1)}{2}$ of the objective function.
%This implies that we can achieve arbitrarily high value of the objective function.
For every $n\geq 2$, 
the optimal solution $x_1 = 1$, $x_2 = -1$ corresponds, via the constructions of Section \ref{sec:red2poly},
to a clique of $n$ vertices where each vertex is assigned $\frac 1 n$, 
and each edge is assigned 
$-\frac 1 {n(n-1)}$.
\end{slexample}


% !TEX root = lin-prog-main.tex


\section{Undecidability of integer solvability} \label{sec:ilp-undecid-proof}

\newcommand{\zeroinstr}{\text{\sc zero}}
\newcommand{\updinstr}{\text{\sc upd}}
We prove Theorem~\ref{thm:ilp-undecid} by showing undecidability of \finineqsolvname$(\Int)$.
We proceed by reduction from the reachability problem of counter machines.

We conveniently define a \emph{$d$-counter machine}
as a finite set of instructions $I$, where each instruction is a function
\[
i : \set{1 \ldots d} \to \Int \cup \set{\zeroinstr}
\]
that specifies, for each counter $k\in\set{1,\ldots, d}$, either the additive update of $k$ (if $i(k) \in \Int$)
or the zero-test of $k$ (if $i(k) = \zeroinstr$).
Configurations of $M$ are nonnegative vectors $c\in \Nat^d$, and
each instruction induces steps between configurations: $c \trans {i} c'$ if
$c'(k) = c(k) + i(k)$ whenever $i(k)\in\Nat$, and $c'(k) = c(k) = 0$ whenever $i(k)=\zeroinstr$.
A run of $M$ is defined as a finite sequence of steps
\begin{align} \label{eq:run}
c_0 \trans{i_1} c_1 \trans{i_2} \ldots \trans{i_n} c_n .
\end{align}
The reachability problem asks, given a machine $M$ and two its configurations, a source $c_0$ and
a target $c_f$, if $M$ admits a run from $c_0$ to $c_f$.
The problem is undecidable, as counter machines can easily simulate classical Minsky machines.%
\footnote{A $d$-counter machine resembles a vector addition system with zero tests.
A Minsky machine with $n$ states and $k$ counters can be simulated by an $(n+k)$-counter machine,
by encoding control states into additional counters.
%and has undecidable reachability already for $d=5$.
} 

For $k\in\set{1,\ldots,d}$ we denote by $\zeroinstr(k) = \setof{i\in I}{i(k)=\zeroinstr}$ 
the set of instructions that zero-test counter $k$, and
$\updinstr(k) = \setof{i\in I}{i(k)\in\Int}$ 
the set of instructions that update counter $k$.

Given a $d$-counter machine $M$ and two configurations $c_0, c_f$, 
we construct an orbit-finite system of inequalities $S = (\vr A, \vr t)$ 
such that $M$ admits a run from $c_0$ to $c_f$ if and only if 
$S$ has a finitary nonnegative integer solution.
(Nonegativeness is \emph{enforced} by adding inequalities  $x\geq 0$ for all unknowns $x$.)
%
We describe construction of $S$ gradually, on the way giving intuitive explanations 
and sketching the proof of the if direction.

The system $S$ has unknowns $e_{\a \b}$ indexed by pairs of distinct atoms $\a\b\in\otu {\A} 2$, and
contains the following inequalities:
%
\begin{align}\label{eq:graph constraint}
e_{\a \b} &\leq 1 \qquad\qquad (\a\b\in\otu \A 2).
\end{align}
%
Therefore, in every solution the unknowns $e_{\a\b}$ define a directed graph $G$, where atoms are vertices, 
$e_{\a\b}=1$ encodes an edge from $\a$ to $\b$ and $e_{\a\b}=0$ encodes a non-edge.
In case of a finitary solution, the graph $G$ is finite (when atoms with no adjacent edges are dropped).
%
Let us fix two distinct atoms $\iota, \zeta\in \A$.
The system $S$ contains the following further equations and inequalities:
%
\begin{align}\label{eq:degree constraint one}
\sum_{\b\neq \a} e_{\b\a} \ =  \ \sum_{\b\neq \a} e_{\a\b} \leq 1 \qquad\qquad
(\a\in\A\setminus\set{\iota,\zeta}) 
\end{align}
%
enforcing that in-degree of every vertex, except for $\iota$ and $\zeta$, is the same as its out-degree, and equal 0 or 1, and also
\begin{align}\label{eq:degree constraint two}
\sum_{\b\neq \iota} e_{\b\iota} = 0 \qquad\qquad \sum_{\b\neq \iota} e_{\iota\b} = 1 \qquad\qquad
\sum_{\b\neq \zeta} e_{\b\zeta} = 1 \qquad\qquad \sum_{\b\neq \zeta} e_{\zeta\b} = 0 
\end{align}
enforcing that in-degree of $\iota$ and out-degree of $\zeta$ are 0, 
while out-degree of $\iota$ and in-degree of $\zeta$ are 1.
Thus atoms split into three categories: inner nodes (with in- and out-degree equal 1), 
end nodes ($\iota$ and $\zeta$) and non-nodes (with in- and out-degree equal 0).
Therefore, the graph $G$ defined by a finitary solution consists of a directed path from $\iota$ to $\zeta$ 
plus a number of vertex disjoint directed cycles.
The path will be used below to encode a run of $M$:
each edge, intuitively speaking, will be assigned a configuration of $M$, while each inner node will be assigned an instruction of $M$.

The system $S$ has also unknowns $t_{i \alpha}$ indexed by instructions $i\in I$ of $M$ and atoms
$\a\in\A$, and the following equations:
%
\begin{align}\label{eq:instr constraint}
\sum_{i \in I} t_{i \a} &= \sum_{\b\neq \a} e_{\a\b} \qquad\qquad (\a\in\A\setminus\set{\iota, \zeta}).
\end{align}
%
Therefore in every finitary solution, for each inner node $\a$ of the above-defined graph $G$,
there is exactly one instruction $i\in I$ such that $t_{i \a}$ equals 1
(intuitively, this instruction $i$ is \emph{assigned} to node $\a$), 
and $t_{i\a}$ equals 0 for all other instructions.
(This applies to \emph{all} inner nodes of $G$, both those on the path as well as those on cycles.)
For non-nodes $\a$, all $t_{i\a}$ are necessarily equal 0.
Note that the values of unknowns $t_{i\iota}$ and $t_{i\zeta}$ are unrestricted, as they are irrelevant.

Finally, the system $S$ contains unknowns $c_{\a\b\c k}$ indexed by $\a\b\c\in\otu \A 3$ and 
$k\in\setto d$.
The following inequalities:
%
\begin{align}\label{eq:cleqe}
c_{\a\b\c k} & \leq e_{\a\b} \qquad\qquad (\a\b\c\in\otu \A 3, k \in\set{1\ldots d})
\end{align}
%
enforce that, whatever atom $\c$ is, the value of unknown $c_{\a\b\c k}$ may be 0 or 1 when $\a\b$ is an edge
(i.e., when $e_{\a\b}=1$), but $c_{\a\b\c k}$ is forcedly 0 when $\a\b$ is a non-edge (i.e., when $e_{\a\b}=0$).
The underlying intuition is that for each $k\in\setto d$, we represent the $k$th coordinate of the configuration \emph{assigned} to the edge $\a\b$ by the 
(necessarily finite) sum
%
\begin{align} \label{eq:encconf}
\sum_{\c\notin\set{\a, \b}} c_{\a\b\c k}.
\end{align}
%
(In particular, configurations assigned to non-edges are necessarily zero on all coordinates.)
In agreement with this intuition,
we add to $S$ the requirement that the configuration assigned to the edge 
outgoing from $\iota$ is the source $c_0$,
and the configuration assigned to the edge incoming to $\zeta$ is the target $c_f$:
%
\begin{align*}%\label{eq:init final}
\sum_{\b,\c\neq \iota} c_{\iota\b\c k}  = c_0(k)  \qquad\qquad  %(k \in \set{1,\ldots, d}) \\ % \notag \\
\sum_{\b,\c\neq \zeta} c_{\b\zeta\c k}  = c_f(k)  \qquad\qquad (k \in \set{1,\ldots, d}).
\end{align*}
%\arka{Since $k$ is not an atom, I think we should separate $\a,\b,\c$ and $k$ and write something like $c_{(\a\b\c,k)}$.}
%
Furthermore, in order to enforce correctness of encoding of a run of $M$,
we add to $S$ equations that relate, intuitively speaking, two consecutive configurations.
Recall that, due to \eqref{eq:degree constraint one}--\eqref{eq:degree constraint two} 
and \eqref{eq:cleqe}, for every $\a\in\A\setminus\set{\iota,\zeta}$, 
unknowns $c_{\b\a\c k}$ may be positive for at most one 
$\b\in\A$; likewise unknowns $c_{\a\b \c k}$.
We add to $S$ the following equations:
%
\begin{align}\label{eq:update constraint}
\sum_{\b,\c\neq \a} c_{\b\a\c k} \ + \sum_{i \in \updinstr(k)} i(k) \cdot t_{i \a}
\ = \sum_{\b,\c\neq \a} c_{\a\b\c k} %\notag
\qquad\qquad
(\a\in\A\setminus\set{\iota, \zeta}, k \in \set{1\ldots d}).
\end{align}
%\piotrek{I would add here some explanation, because it can be confusing that we sum all configurations with $\a$ on the second component on the left hand side.
%Recall that as variables $e_{\d\a}$ defines a graph with in-degree $1$ there is only one $\b$ for which variables $c_{\b\a\c k}$ can be not equal $0$ and only 
%one $\b$ for which variables $c_{\a\b\c k}$ can be not equal $0$.} 
%
These equalities say that for every inner node or non-node $\a$ (i.e., every atom except for the end nodes $\iota$ and $\zeta$),
on every coordinate $k$, the configuration incoming to $\a$ differs from the configuration
outgoing from $\a$ exactly by the sum
\[
\sum_{i \in \updinstr(k)} i(k) \cdot t_{i \a}
\]
ranging over those instructions $i$ of $M$ that update counter $k$.
%We can safely further restrict the range of $i$ to those instructions for which $t_{i \a} \neq 0$.
Remembering that for each $\a$ there is at most one instruction $i$ satisfying $t_{i \a} \neq 0$,
we get that the configurations differ on coordinate $k$ by exactly $i(k)$ 
(if $i$ updates counter $k$) or the configuration are equal on coordinate $k$
(if $i$ zero-tests counter $k$, or there is no instruction $i$ such that $t_{i \a} \neq 0$).

In order to deal with zero tests, we add to $S$ not just the inequalities~\eqref{eq:cleqe},
but the following strengthening thereof:
%
\begin{align}\label{eq:zero test}
c_{\a\b\c k} + \!\!\!\! \sum_{i\in\zeroinstr(k)} \!\!\! t_{i\a} & \leq e_{\a\b} 
\qquad\qquad (\a\b\c\in\otu \A 3, k \in\set{1\ldots d}).
\end{align}
%
In consequence, for every edge $\a\b$, if the instruction $i$ assigned to $\a$ updates counter $k$,
%or no instruction is executed at $\a$, 
\eqref{eq:zero test} does not restrict further the  $k$th coordinate of the configuration assigned to $\a\b$.
But if the instruction $i$ assigned to $\a$ zero-tests counter $k$, the sum
\[
\sum_{i\in\zeroinstr(k)} t_{i\a}
\]
equals $1$ and therefore the $k$th coordinate of the configuration assigned to $\a\b$, 
encoded by~\eqref{eq:encconf}, is necessarily 0
(the same applies also to the configuration incoming to $\a$, due to inequalities \eqref{eq:update constraint} below). 
The above considerations apply to \emph{all} edges of $G$, both those on the path as well as those on cycles.
As a further consequence, for a  non-edge $\a\b$, the configuration assigned at $\a\b$, encoded by~\eqref{eq:encconf}, 
is necessarily the zero configuration.


The construction of $S$ is thus completed, and it remains to argue towards its correctness:
%
\begin{lemma} \label{lem:undecid-corr}
$M$ admits a run from $c_0$ to $c_f$ if and only if 
$S$ has a finitary nonnegative integer solution.
\end{lemma}
%
\begin{proof}%[Proof of Lemma \ref{lem:undecid-corr}]
For the `if' direction, given a finitary nonnegative integer solution of $S$, we consider
the graph $G$ determined by values of unknowns $e_{\a\b}$, as discussed in the course of construction, 
consisting of inner nodes and two end nodes,
and having the form of a finite directed path plus (possibly) a number of directed cycles. 
By the construction of $S$, each edge of $G$ has assigned a configuration of $M$,
and each inner node has assigned an instruction of $M$, so that the configuration on the edge outgoing
from an inner node is exactly the result of executing its instruction on the configuration assigned to the incoming edge.
(As above, this applies to \emph{all} inner nodes and edges of $G$, both those on the path as well as those on cycles.)
Ignoring the cycles of $G$, we conclude that the sequence of configurations and
instructions along the path of $G$ is a run of $M$ from $c_0$ to $c_f$.

For the `only if' direction, given a run of $M$ from $c_0$ to $c_n$ as in~\eqref{eq:run},
one constructs a solution of $S$ in the form of a sole path involving end nodes $\a_0 = \iota, \a_{n+1} = \zeta$,
$n$ inner nodes $\a_1, \ldots, \a_n$, and $n+1$ edges $\a_j \a_{j+1}$. 
Thus unknowns $e_{\a_j, \a_{j+1}}$ are equal 1.
The values of unknowns $t_{i \a_j}$ are determined by instructions $i_j$ used in the run,
and the values of the unknowns $c_{\a_j \a_{j+1} \c k}$, for sufficiently many fresh atoms $\c$, are determined
by configurations $c_j$. 
All other unknowns are equal 0.
\end{proof}


\begin{slremark}
The proof does not adapt to \finnonnegsolvname$(\Z)$.
Indeed, the standard way of transforming inequalities into equations involves 
adding an infinite set of additional unknowns, that might be all non-zero.
\end{slremark}


% !TEX root = lin-prog-main.tex


\section{Conclusions} \label{sec:conc}

As two main contributions, we show two contrasting results:
decidability of orbit-finite linear programming, and 
undecidability of orbit-finite integer linear programming.
For decidability, we invent a novel concept of setwise-$T$-orbit, and provide a reduction
to a finite but polynomially-parametrised linear programming.
%
In addition to the decidability of the latter problem,
%relying on encoding into real arithmetic\piotrek{Do we rely on this encoding?}, 
we show that it can be solved in \EXPTIME, and even in \PTIME for every fixed atom dimension. 
We thus match, in case of fixed atom dimension, the complexity of classical linear programming.

We consider non-strict inequalities for presentation only, 
and our decision procedures
may be straightforwardly adapted to to mixed systems of strict and non-strict inequalities.

We leave a number of intriguing open questions, all of them except the last one referring to
linear programming:

\begin{question}
In this paper we only consider finitely supported solutions.
We do not know the decidability status of %orbit-finite (integer) 
linear programming when this restriction is dropped
(like in \cite{KKOT15}).
It is decidable for finitary inequalities, where existence of a solution implies existence of an equivariant one \cite{szymon-oral}.
\end{question}

\begin{question}
We exclusively consider equality atoms, and
extension to richer structures seems highly non-trivial.
In the important case of \emph{ordered atoms},
we are currently only able to prove decidability of 
\solvname$(\F)$, for any commutative ring $\F$.
\end{question}

\begin{question}
It is very natural to ask if the classical duality of linear programs extends to the orbit-finite setting.
According to our initial observations this is indeed the case, under the restriction that either 
(v) vertical vectors of a matrix and the target vector are finitary, or
(h) horizontal vectors of a matrix and objective function are finitary. 
Whenever the primary program satisfies one of the conditions (v), (h), the dual one satisfies the other one.
\end{question}

\begin{question}
Solution sets of orbit-finite systems are not always finitely generated.
Therefore an interesting question arises if one can compute a representation of solution sets
that would enable testing for equality or inclusion of such sets?
%
For instance, the solution set of the system of inequalities \,
$\sum_{\a \in \A\setminus\set{\b}} \a  \ \geq \ \b$ \ \ $(\b \in \A)$, 
in matrix form
%
\begin{align*} 
%& \qquad \begin{matrix} \ \cdots &  \a  & \cdots\ \end{matrix} \\
& 
%\begin{matrix}
%\vdots \\
%\b \\
%\vdots \\
%\end{matrix}
%\
\begin{bmatrix}
\ -1 & 1 & 1 & \cdots \ \\
\ 1 & -1 & 1 & \cdots \ \\
\ 1 & 1 & -1 & \cdots \ \\
 \vdots     & \vdots & \vdots & \ddots \   
\end{bmatrix}
\cdot
%\begin{bmatrix}
%\vdots \\
%\a \\
%\vdots \\
%\\
\ \vr x %\\
%\\
%\end{bmatrix}
\quad
\geq
\quad
\begin{bmatrix}
\, 0 \, \\
\, 0 \, \\
\, 0 \, \\
 \vdots 
\end{bmatrix}
\end{align*}
is not equal to the cone generated by (=non-negative linear combinations of) an orbit-finite set.
%
%Indeed, it is easy to see that solutions are exactly all those finitary vectors $\vr x$ where 
%$\sum_{\a\in\A} \vr x(a) \geq 2 \cdot \max_{a\in\A} \vr x$.
%Therefore the values 
%
\end{question}

\begin{question}
We would be happy to know if our general \EXPTIME upper complexity bound is tight.
\end{question}

\begin{question}
Concerning integer linear programming,
an intriguing research task is to identify the decidability borderline.
For instance, we suspect decidability in case when all inequalities are finitary.
The reduction proving undecidability produces a system of atom dimension $3$, and it is unclear if the dimension can be lowered
to $2$.
In case of atom dimension 1 we suspect decidability (along the lines of \cite{HLT17}).
\end{question}




%%
%% The acknowledgments section is defined using the "acks" environment
%% (and NOT an unnumbered section). This ensures the proper
%% identification of the section in the article metadata, and the
%% consistent spelling of the heading.

\begin{acks}
We are grateful to Miko{\l}aj Boja{\'n}czyk, Asia Fijalkow, Lorenzo Clemente and Damian Niwi{\'n}ski 
for fruitful discussions.

We acknowledge partial support of 
NCN Opus grant 2019/35/B/ST6/02322 (Arka Ghosh and S{\l}awomir Lasota),
of NCN Preludium grant 2022/45/N/ST6/03242 (Arka Ghosh),
and of
ERC Starting grant INFSYS, agreement no.~950398 (Piotr Hofman and S{\l}awomir Lasota).
\end{acks}

%%
%% The next two lines define the bibliography style to be used, and
%% the bibliography file.


%%
%% If your work has an appendix, this is the place to put it.

\appendix



% !TEX root = lin-prog-main.tex


%\subsection{Spans}


%We introduce notation useful in proving Theorem \ref{thm:equiv-problems}.
%For a subset $P\subseteq \GLin B$ and $\F\in \set{\R,\Int}$, we define $\Span {\F} P\subseteq \GLin B$ as the set of 
%all linear $\F$-combinations of vectors from $P$:
%%, forming a subspace of $\GLin B$:
%\begin{align*} %\label{eq:span}
%\begin{aligned}
% \Span {\F} P \ = \ \ & \setofbeg{q_1 \cdot \vr p_1 + \ldots + q_k \cdot \vr p_k}{k\geq 0,}\\
%& \ \setofend{\ q_1,\ldots,q_k\in \F, \ \vr p_1, \ldots, \vr p_k\in P}.
%\end{aligned}
%\end{align*}
%%
%Given a matrix $\vr A \in \GLin {B\times C}$ with rows $B$ and columns $C$, we can define a partial operation of 
%multiplication of $\vr A$ by a vector $\vr v\in\GLin C$ in an expected way:
%\[
%(\mult {\vr A}{\vr v})(b) = \innerprod {\vr A(b, \_)} {\vr v}
%\]
%for every $b\in B$.
%The result $\mult {\vr A} {\vr v} \in \GLin B$ is well-defined if $\innerprod {\vr A(b, \_)}{\vr v}$ is well-defined for all $b\in B$.
%%Considering the matrix, via its column-indexed representation $\vr A : C \to \GLin B$, as an indexed set of column vectors
%%from $\GLin B$, 
%The multiplication $\mult {\vr A}{\vr v}$ can be also seen as an \emph{orbit-finite} linear combination of 
%column vectors $\vr A(\_, c)$, for $c\in C$, with coefficients given by $\vr v$. 
%This allows us to define the \emph{span} of $\vr A$ seen as a $C$-indexed orbit-finite set of 
%vectors $\vr A(\_, c) \in \GLin B$:
%%
%\begin{align*} %\label{eq:Gspan}
%\GSpan {\F} {\vr A}  :=  \setof{\mult{\vr A}{\vr v}}{\vr v : C\tofs \F, \ \mult{\vr A}{\vr v} \text{ well-def.}}.
%\end{align*}
%%
%Therefore, a system of inequalities $(\vr A, \vr t)$ has a solution if 
%$\GSpan {\F} {\vr A}$ contains some vector $\vr u \geq \vr t$.
%When $\vr v$ is finitary, well-definedness is vacuous, and we may define:
%By restricting to finitary vectors $\vr v$ we do not need to assume well-definedness, as it vacuously holds.
%The finitary span of $\vr A$ is thus defined as:
%\[
%\Span {\F} {\vr A}  :=  \setof{\mult{\vr A}{\vr v}}{\vr v : C\tofin \F}  =  \Span {\F} P,
%\]
%%We intentionally overload the notation, as  $\Span {} {\vr A} = $ $\Span {} P$, 
%for $P = \setof{\vr A(\_, c)}{c\in C}$ the set of column vectors of $\vr A$.


%\subsection{Proof of Theorem~\ref{thm:equiv-problems}}

%The proof of mutual
%reductions between \ineqsolvname$(\F)$ and \nonnegsolvname$(\F)$
%amounts to lifting of standard arguments from finite to orbit-finite systems, 
%and checking that all constructed objects are finitely supported. 
%We include the reductions here mostly in order to get acquainted with orbit-finite systems.
%One of the remaining two reductions build on results of \cite{GHL22}.
%

%\para{Reduction of \ineqsolvname$(\F)$ to \nonnegsolvname$(\F)$}
%
%Consider an orbit-finite system of linear inequalities supported by $S$,
%\begin{equation}\label{eq: ineq}
%\vr{A} \cdot \vr{x} \geq \vr{t},
%\end{equation}
%where $\vr A : B \times C \tofs \F$ and $\vr t : B\tofs \F$.
%We construct the new system of linear equations $(\vr A', \vr t)$,
%with the same target vector $\vr t$, and
%with matrix is of the form
%$
%\vr A' : B \times (C \uplus C \uplus B) \tofs \F,
%$ 
%such that
%the original system has a finitely-supported solution if and only if the new system
%has a non-negative finitely-supported one.
%In the new system, we double each variable $x$ into $x_+$ and $x_-$, and we add a fresh variable per each equation.
%The matrix $\vr A'$ of the new system is a composition of $\vr A$, $-\vr A$, and opposite of the identity diagonal matrix
%$B\times B\tofs \F$:

%\vspace{-4mm}
%\begin{align*}
%\vr A' \ = \ 
%\left[
%\begin{matrix}
%& &  \phantom{1} \\  & \ \ \ \vr A & \phantom{\ddots}    \\     &  &    \phantom{1}
%\end{matrix}
%\right\rvert
%\begin{matrix}
%  & &   \\  & \ \ - \vr A &  \phantom{\ddots}  \\     &  &   
%\end{matrix}
%\left\rvert
%\begin{matrix}
%\; -1 & &   \\  & \ddots &    \\     &  & -1 \;  
%\end{matrix}
%\right]
%\end{align*}
%%
%%\piotrek{I think it should be $-1$ instead of $1$, as we have to subtract from the real solution to get to $t$}
%$\vr A'$ is thus supported by $S$.
%Correspondingly, solutions of the new system split into 3 vectors $\vr x_+, \vr x_- : C\tofs \F$ and $\vr y : B\tofs \F$ satisfying
%% 
%\begin{equation}\label{eq: non-neg eq equiv to ineq}
%\vr{A} \cdot \vr{x}_+ - \vr{A} \cdot \vr{x}_- - \vr{y} = \vr{t}.
%\end{equation}
%%
%Each such nonnegative solution, supported by $S'$, yields a solution $\vr{x}_+ - \vr{x}_-$ of the original system \eqref{eq: ineq},
%also supported by $S'$.
%In the opposite direction, a solution of the original system supported by $S'$
%determines a non-negative solution of (\ref{eq: non-neg eq equiv to ineq}):
%%
%\begin{align*}
%\vr{x}_+(c) & =
%\begin{cases}
%\vr{x}(c) & \text{ if }\vr{x}(c) \geq 0, \\
%0         & \text{ otherwise;}
%\end{cases} \\
%\vr{x}_-(c) & =
%\begin{cases}
%-\vr{x}(c) & \text{ if }\vr{x}(c) < 0, \\
%0         & \text{ otherwise;}
%\end{cases} \\
%\vr{y}(c) & = (\vr{A}\cdot \vr{x} - \vr t)(c).
%\end{align*}
%%
%This solution is supported by $S\cup S'$.
%
%\para{Reduction of \nonnegsolvname$(\F)$ to \ineqsolvname$(\F)$}
%%
%Consider an orbit-finite system of linear equations supported by $S$:
%%
%\begin{align*}%\label{eq: eq}
%\vr{A} \cdot \vr{x} = \vr{t}.
%\end{align*}
%%
%Its nonnegative solutions are exactly solutions of the following system of linear inequalities, also supported by $S$:
%\begin{align*}%\label{eq: ineq equiv to eq}
%\vr{A} \cdot \vr{x} \geq \vr{t} \qquad \vr{A} \cdot \vr{x} \leq \vr{t} \qquad \vr{x} \geq 0.
%\end{align*}
%%
%
%\begin{slremark}
%The above reductions preserve row-finiteness.
%\end{slremark}
%
%\para{Reduction of \finineqsolvname$(\F)$ to \ineqsolvname$(\F)$}
%%
%Given a system of inequalities $\vr{A} \cdot \vr{x} \geq \vr{t}$ supported by $S$, where
%$\vr A : B\times C\tofs \F$, we construct a new system of inequalities
%$\vr{A}' \cdot \vr{x}' \geq \vr{t}'$ by
%extending the column index $C$ by one additional variable $y$ and extending the system by one inequality:
%
%\vspace{-2mm}
%\begin{align*}
%\vr A' \ = \ 
%\left[
%\begin{matrix}
%  & & \phantom{0}  \\  & \ \ \vr A \ \ \phantom{\vdots} &     \\     &  & \phantom{0} \\ \hline
%  1 & \cdots & 1
%\end{matrix}
%\right\rvert
%\!\!\!
%\left\rvert
%\begin{matrix}
%  0 \\     \vdots \\   0 \\       \hline -1
%\end{matrix}
%\right]
%\qquad
%\vr t' \ = \ \begin{bmatrix}
%\phantom{0} \\ \ \vr t \phantom{\vdots} \\ \phantom{0} \\ \hline 0 
%\end{bmatrix}
%\end{align*}
%%
%The new system is supported by $S$, and its solutions have the form $\vr x' = (\vr x, y)$, where
%\[
%\vr{A} \cdot \vr{x} \geq \vr{t} \qquad \sum_{c \in C} \vr{x}(c) \geq y.
%\]
%%By Claim~\ref{claim:?}, a
%Any such finitely supported solution is necessarily finitary,
%and therefore new system has a finitely-supported solution if and only if the 
%the original system has a finitary solution.
%%
%%\begin{claimproof}
%%Let $\vr{x}$ and $y$ be such that $\vr{x}$ is finitely supported and
%%\[
%%\vr{A} \cdot \vr{x} \geq \vr{t} ,\ \sum_{c \in C} \vr{x}(c) \geq y\ .
%%\]
%%Let $U$ be an infinite $\supp{\vr{x}}$-orbit of $C$.
%%$sum_{c \in C} \vr{x}(c)$ is well-defined only if $sum_{c \in U} \vr{x}(c)$ is well-defined.
%%Since $U$ is a $\supp{\vr{x}}$-orbit, there exists $r_U$ such that $\vr{x}(c) = r_U$ for all $c \in U$.
%%Then $sum_{c \in U} \vr{x}(c)$ is well-defined only if $r_U = 0$.
%%Since $U$ was an arbitrary infinite $\supp{\vr{x}}$-orbit of $C$,
%%this implies $\vr{x}$ is non-zero only on finite $\supp{\vr{x}}$-orbits.
%%In other words, $\vr{x}$ is finite.
%%\end{claimproof}
%%%
%%Using the above claim we can say that the system $\vr{A} \cdot \vr{x} \geq \vr{t}$ has a finite solution if and only if (\ref{eq: with eq forcing fin}) has a finitely supported solution.
%
%
%\para{Reduction of \ineqsolvname$(\F)$ to \finineqsolvname$(\F)$}
%%
%
%We rely on the following result of~\cite{GHL22}%
%\footnote{The result, as shown in~\cite{GHL22}, holds for any commutative ring $\F$.}:
%%
%\begin{claim}[\cite{GHL22} Claim~20]
%Let $\F \in\set{\Int,\R}$.
%Given an $S$-supported orbit-finite matrix $\vr{A}$ one can effectively construct an $S$-supported orbit-finite matrix $\spread{\vr{A}}$
%such that $\GSpan{\F}{\vr{A}} = \Span{\F}{\spread{\vr{A}}}$.
%\end{claim}
%%
%In consequence, the system $(\vr{A}, \vr{t})$ of inequalities has a finitely-supported solution if and only
%of the system $(\spread{\vr A}, \vr{t})$ of inequalities has a finitary one.
%\qed
%
%
%\subsection{Proof of Theorem~\ref{thm:finnegsolv-decid}}
%In case of $\F = \R$, by a direct reduction of \finnonnegsolvname$(\F)$ to \nonnegsolvname$(\F)$
%and Theorem~\ref{thm:lp-decid}.
%In case of $\F = \Int$, by 
%\TODO{todo}
%
%\qed
%
%




\section{Missing proofs}

%\TODO{todo}

We start by introducing notation 
useful in proving Theorems \ref{thm:equiv-problems} and \ref{thm:equiv-max}.
For subsets $P\subseteq \GLin B$ and $\F\subseteq \R$, we define $\Span {\F} P\subseteq \GLin B$ as the set of 
all linear $\F$-combinations of vectors from $P$:
%, forming a subspace of $\GLin B$:
\begin{align*} %\label{eq:span}
\begin{aligned}
 \Span {\F} P \ = \ \  \setof{q_1 \cdot \vr p_1 + \ldots + q_k \cdot \vr p_k}{k\geq 0,
\ q_1,\ldots,q_k\in \F, \ \vr p_1, \ldots, \vr p_k\in P}.
\end{aligned}
\end{align*}
%
Recall that given a matrix $\vr A \in \GLin {B\times C}$ with rows $B$ and columns $C$, we can define a partial operation of 
multiplication of $\vr A$ by a vector $\vr v\in\GLin C$ in an expected way:
\[
(\mult {\vr A}{\vr v})(b) = \innerprod {\vr A(b, \_)} {\vr v}
\]
for every $b\in B$.
The result $\mult {\vr A} {\vr v} \in \GLin B$ is well-defined if $\innerprod {\vr A(b, \_)}{\vr v}$ is well-defined for all $b\in B$.
%Considering the matrix, via its column-indexed representation $\vr A : C \to \GLin B$, as an indexed set of column vectors
%from $\GLin B$, 
For $c\in C$ we denote by $\vr A(\_, c) \in \GLin B$ the corresponding (column) vector.
The multiplication $\mult {\vr A}{\vr v}$ can be also seen as an \emph{orbit-finite} linear combination of 
column vectors $\vr A(\_, c)$, for $c\in C$, with coefficients given by $\vr v$. 
This allows us to define the \emph{span} of $\vr A$ seen as a $C$-indexed orbit-finite set of 
vectors $\vr A(\_, c) \in \GLin B$:
%
\begin{align*} %\label{eq:Gspan}
\GSpan {\F} {\vr A}  :=  \setof{\mult{\vr A}{\vr v}}{\vr v : C\tofs \F, \ \mult{\vr A}{\vr v} \text{ well-def.}}.
\end{align*}
%
Therefore, a system of inequalities $(\vr A, \vr t)$ has a solution if 
$\GSpan {\F} {\vr A}$ contains some vector $\vr u \geq \vr t$.
When $\vr v$ is finitary, well-definedness is vacuous, and we may define:
%By restricting to finitary vectors $\vr v$ we do not need to assume well-definedness, as it vacuously holds.
%The finitary span of $\vr A$ is thus defined as:
\[
\Span {\F} {\vr A}  :=  \setof{\mult{\vr A}{\vr v}}{\vr v : C\tofin \F}  =  \Span {\F} P
\]
%We intentionally overload the notation, as  $\Span {} {\vr A} = $ $\Span {} P$, 
for $P = \setof{\vr A(\_, c)}{c\in C}$ the set of column vectors of $\vr A$.
Therefore, a system of inequalities $(\vr A, \vr t)$ has a finitary solution if 
$\Span {\F} {\vr A}$ contains some vector $\vr u \geq \vr t$.


\subsection{Proof of Theorems \ref{thm:equiv-problems} and \ref{thm:equiv-max} 
(Section \ref{sec:problems})}
\label{sec:problems-proofs-one}

%
Recall that we consider supremum of a maximisation problem to be $-\infty$ if the constraints in the problem are infeasible.
Therefore proving that two maximisation problems have the same supremum also proves that the underlying systems of inequalities are equisolvable.
In consequence, Theorem \ref{thm:equiv-max} implies \ref{thm:equiv-problems}, and hence
we concentrate
in the sequel on proving the former one.


The proof of mutual
reductions between \ineqmaxname$(\F)$ and \nonnegmaxname$(\F)$
amounts to lifting of standard arguments from finite to orbit-finite systems, 
and checking that all constructed objects are finitely supported. 
We include the reductions here mostly in order to get acquainted with orbit-finite systems.
One of the remaining two reductions builds on results of \cite{GHL22}.
%

\para{Reduction of \ineqmaxname$(\F)$ to \nonnegmaxname$(\F)$}
%
Consider an instance $(\vr{A},\vr{t},\vr{s})$ of \ineqmaxname$(\F)$, supported by $S$,
where $\vr A : B \times C \tofs \F$, $\vr t : B\tofs \F$ and $\vr{s} : C\tofs \F$.
We construct an instance $(\vr A', \vr t, \vr{s}')$ of \nonnegmaxname$(\F)$,
with the same target vector $\vr t$, and
$
\vr A' : B \times (C \uplus C \uplus B) \tofs \F,
$
$\vr{s}' : (C \uplus C \uplus B) \tofs \F$,
such that
\begin{align*} %\label{eq:supsup}
\textnormal{supremum}(\vr A', \vr t, \vr{s}')
=
\textnormal{supremum}(\vr A, \vr t, \vr{s}).
\end{align*}
In the new system, we double each variable $x$ into $x_+$ and $x_-$, and we add a fresh variable per each equation.
The matrix $\vr A'$ of the new system is a composition of $\vr A$, $-\vr A$, and the diagonal matrix 
$B \times B\tofs \F$ with $-1$ in the diagonal:

\vspace{-2mm}
\begin{align*}
\vr A' \ = \ 
\left[
\begin{matrix}
  & &  \phantom{1} \\  & \ \ \ \vr A & \phantom{\ddots}    \\     &  &    \phantom{1}
\end{matrix}
\right\rvert
\begin{matrix}
  & &   \\  & \ \ - \vr A &  \phantom{\ddots}  \\     &  &   
\end{matrix}
\left\rvert
\begin{matrix}
\; -1 & &   \\  & \ddots &    \\     &  & -1 \;  
\end{matrix}
\right]
\end{align*}
%
%\piotrek{I think it should be $-1$ instead of $1$, as we have to subtract from the real solution to get to $t$}

\noindent
Similarly, $\vr{s}'$ is defined as the composition of $\vr{s}$, $-\vr{s}$ and the zero vector $B \tofs \F$:

\vspace{-4mm}
\begin{align*}
\vr s' \ = \ 
\left[
\begin{matrix}
  & \vr s & 
\end{matrix}
\right\rvert
\begin{matrix}
  & - \vr s &   
\end{matrix}
\left\rvert
\begin{matrix}
\; 0 & \cdots & 0 \;  
\end{matrix}
\right]
\end{align*}
$\vr A'$ and $\vr{s}'$ are thus supported by $S$.

Any vector 
$
\vr{x}' :(C \uplus C \uplus B) \tofs \F
$
can be written as
\[
\vr x' \ = \ (\vr{x}_+ \rvert \vr{x}_- \rvert \vr{y}),
\]
where $\vr{x}_+,\vr{x}_- : C \tofs \F$ and $\vr{y} : B \tofs \F$.
%
If any such non-negative vector $\vr x'$ satisfies the above constructed system of constraints, i.e. if we have
%
\begin{equation}\label{eq: non-neg max equiv to ineq}
\vr{A}' \cdot (\vr{x}_+ \rvert \vr{x}_- \rvert \vr{y}) = \vr{t},
\end{equation}
%
then then vector ${\vr{x}_+ - \vr{x}_-}$, supported by $\supp{\vr{x}'}$,
is a solution of $(\vr A, \vr t)$, namely
\[
\vr{A} \cdot (\vr{x}_+ - \vr{x}_-) \geq
\vr{A} \cdot (\vr{x}_+ - \vr{x}_-) - \vr{y} =
\vr{A}' \cdot (\vr{x}_+ \rvert \vr{x}_- \rvert \vr{y}) = \vr{t}.
\]
Furthermore, by the very definition of $\vr s'$ we have
%
\begin{align} \label{eq:ss}
\vr{s}' \cdot (\vr{x}_+ \rvert \vr{x}_- \rvert \vr{y}) = \vr{s} \cdot (\vr{x}_+ - \vr{x}_-),
\end{align}
%
which implies
$
\textnormal{supremum}(\vr A', \vr t, \vr{s}')
\leq
\textnormal{supremum}(\vr A, \vr t, \vr{s}).
$

In the opposite direction, given a finitely supported vector $\vr{x}$ such that 
$\vr{A} \cdot \vr{x} \geq \vr{t}$, we define a non-negative vector 
$\vr x' = (\vr{x}_+ \rvert \vr{x}_- \rvert \vr{y})$ supported by $\supp{\vr{x}}\cup S$ as follows:

\begin{minipage}{0.35\linewidth}
\begin{align*}
\vr{x}_+(c) & =
\begin{cases}
\vr{x}(c) & \text{ if }\vr{x}(c) \geq 0, \\
0         & \text{ otherwise;}
\end{cases} 
\end{align*}
\end{minipage}
\begin{minipage}{0.35\linewidth}
\begin{align*}
\vr{x}_-(c) & =
\begin{cases}
-\vr{x}(c) & \text{ if }\vr{x}(c) < 0, \\
0         & \text{ otherwise;}
\end{cases} 
\end{align*}
\end{minipage}
\begin{minipage}{0.3\linewidth}
\begin{align*}
\vr{y}(c) & = (\vr{A}\cdot \vr{x} - \vr t)(c). \\
\ 
\end{align*}
\end{minipage}
%
\smallskip

\noindent
Then $\vr{x} = \vr{x}_+ - \vr{x}_-$ and
$
\vr{A}' \cdot (\vr{x}_+ \rvert \vr{x}_- \rvert \vr{y}) =
\vr{A} \cdot \vr{x} - \vr{y} = \vr{t}.
$
The equality \eqref{eq:ss} holds again,
%\[
%\vr{s}' \cdot (\vr{x}_+ \rvert \vr{x}_- \rvert \vr{y}) =
%\vr{s} \cdot (\vr{x}_+ - \vr{x}_-) \ .
%\]
which implies 
$
\textnormal{supremum}(\vr A, \vr t, \vr{s})
\leq
\textnormal{supremum}(\vr A', \vr t, \vr{s}').
$

\para{Reduction of \nonnegmaxname$(\F)$ to \ineqmaxname$(\F)$}
%
For any orbit-finite system of linear equations supported by $S$:
%
\begin{align*}%\label{eq: eq}
\vr{A} \cdot \vr{x} = \vr{t},
\end{align*}
%
its nonnegative solutions are exactly solutions of the following system of linear inequalities, also supported by $S$:
\begin{align*}%\label{eq: ineq equiv to eq}
\vr{A} \cdot \vr{x} \geq \vr{t} \qquad\qquad 
\vr{A} \cdot \vr{x} \leq \vr{t} \qquad\qquad 
\vr{x} \geq 0.
\end{align*}
%
This implies an easy reduction from \nonnegmaxname$(\F)$ to \ineqmaxname$(\F)$.

\begin{slremark}
The above two reductions preserve row-finiteness, i.e., transform a system of finite equations (inequalities) to a system of finite inequalities (equations), or vice versa.
\end{slremark}

\para{Reduction of \finineqmaxname$(\F)$ to \ineqmaxname$(\F)$}
%
Consider an instance $(\vr{A},\vr{t},\vr{s})$ of \finineqmaxname$(\F)$ supported by $S$,
where $\vr A : B\times C\tofs \F$.
We construct an instance $(\vr{A}',\vr{t}',\vr{s}')$ of \ineqmaxname$(\F)$ as follows.
The new system of inequalities
$\vr{A}' \cdot \vr{x}' \geq \vr{t}'$ is obtained by
extending the column index $C$ by one additional variable $y$ and extending the system by one inequality:

\vspace{-2mm}
\begin{align*}
\vr A' \ = \ 
\left[
\begin{matrix}
  & & \phantom{0}  \\  & \ \ \vr A \ \ \phantom{\vdots} &     \\     &  & \phantom{0} \\ \hline
  1 & \cdots & 1
\end{matrix}
\right\rvert
\hspace{-1.25mm}
\left\rvert
\begin{matrix}
  0 \\     \vdots \\   0 \\       \hline -1
\end{matrix}
\right]
\qquad\qquad
\vr t' \ = \ \begin{bmatrix}
\phantom{0} \\ \ \vr t \phantom{\vdots} \\ \phantom{0} \\ \hline 0 
\end{bmatrix}
\end{align*}
%
and the new objective function $\vr s'$ is defined as expected:

\vspace{-4mm}
\begin{align*}
\vr s' \ = \ 
\left[ 
\begin{matrix}
& \;   \vr s  \; & 
\end{matrix}
\rvert
\; 0 \;  
\right].
\end{align*}
%
The so constructed instance is supported by $S$, and its solutions have the form $\vr x' = (\vr x, y)$, where
\[
\vr{A} \cdot \vr{x} \geq \vr{t} \qquad\qquad \sum_{c \in C} \vr{x}(c) \geq y.
\]
%By Claim~\ref{claim:?}, a
Any such finitely supported solution is necessarily finitary.
This implies
$
\textnormal{supremum}(\vr{A},\vr{t},\vr{s}) =
\textnormal{supremum}(\vr{A}',\vr{t}',\vr{s}).
$
%where $(\vr{A}',\vr{t}',\vr{s})$ is an instance of \ineqmaxname$(\F)$.
%
%\begin{claimproof}
%Let $\vr{x}$ and $y$ be such that $\vr{x}$ is finitely supported and
%\[
%\vr{A} \cdot \vr{x} \geq \vr{t} ,\ \sum_{c \in C} \vr{x}(c) \geq y\ .
%\]
%Let $U$ be an infinite $\supp{\vr{x}}$-orbit of $C$.
%$sum_{c \in C} \vr{x}(c)$ is well-defined only if $sum_{c \in U} \vr{x}(c)$ is well-defined.
%Since $U$ is a $\supp{\vr{x}}$-orbit, there exists $r_U$ such that $\vr{x}(c) = r_U$ for all $c \in U$.
%Then $sum_{c \in U} \vr{x}(c)$ is well-defined only if $r_U = 0$.
%Since $U$ was an arbitrary infinite $\supp{\vr{x}}$-orbit of $C$,
%this implies $\vr{x}$ is non-zero only on finite $\supp{\vr{x}}$-orbits.
%In other words, $\vr{x}$ is finite.
%\end{claimproof}
%%
%Using the above claim we can say that the system $\vr{A} \cdot \vr{x} \geq \vr{t}$ has a finite solution if and only if (\ref{eq: with eq forcing fin}) has a finitely supported solution.


\para{Reduction of \ineqmaxname$(\F)$ to \finineqmaxname$(\F)$}
%

We rely on the following result of~\cite{GHL22}%
\footnote{The result, as shown in~\cite{GHL22}, holds for any commutative ring $\F$.}:
%
\begin{claim}[\cite{GHL22} Claim~20]
Let $\F \in\set{\Int,\R}$.
Given an $S$-supported orbit-finite matrix $\vr{M}$ one can effectively construct an $S$-supported orbit-finite matrix $\spread{\vr{M}}$
such that $\GSpan{\F}{\vr{M}} = \Span{\F}{\spread{\vr{M}}}$.
\end{claim}
%
Consider an instance $(\vr{A},\vr{t},\vr{s})$ of \ineqmaxname$(\F)$, and
apply the above claim to the matrix $\vr M$ (left) 
in order to get the matrix $\spread{\vr{M}}$ (right),
\[
\vr{M} \ = \ 
\left[
\begin{matrix}
            &        &  \\
            & \vr A  &  \\
            &        &  \\
              \hline
\phantom{1} & \vr{s} & \phantom{1}
\end{matrix}
\right]
%\]
\qquad\qquad\qquad\qquad
%\[
\spread{\vr{M}} \ = \ 
\left[
\begin{matrix}
            &                  &  \\
            & \vr {A}' &  \\
            &                  &  \\
              \hline
\phantom{1} & \vr{s}'  & \phantom{1}
\end{matrix}
\right]
\]
such that
\begin{align} \label{eq:spany}
\GSpan{\F}{\vr{M}} =
\Span{\F}{\spread{\vr{M}}}.
\end{align}
This yields 
an instance  $(\vr{A}',\vr{t},\vr{s}')$  of \finineqmaxname$(\F)$, supported 
by $\supp{\vr{A},\vr{t},\vr{s}}$.

By the equality \eqref{eq:spany},
for every $r \in \R$ we have the following:
there exists a finitely supported vector $\vr{x}$ such that 
$\mult{\vr{A}}  {\vr{x}} \geq \vr{t}$ and $\innerprod{\vr{s}}  {\vr{x}} = r$ 
if and only if
there exists a finitary vector $\vr{x}'$ such that $\mult{\vr{A}'}  {\vr{x}'} \geq \vr{t}$ and 
$\innerprod {\vr{s}'}  {\vr{x}'} = r$.
In consequence,
\[
\text{supremum}(\vr{A},\vr{t},\vr{s}) =
\text{supremum}(\vr{A}',\vr{t},\vr{s}').
\]

Theorem \ref{thm:equiv-max} is thus proved.
%\end{proof}


\subsection{Proof of Theorem~\ref{thm:finnegsolv-decid} (Section \ref{sec:problems})}
\label{sec:problems-proofs-two}

We consider cases $\F = \R$ and $\F = \Z$ separately. 

\para{Case $\F = \R$}

Decidability of \finnonnegsolvname$(\R)$ follows by a direct reduction of 
\finnonnegsolvname$(\R)$ to \nonnegsolvname$(\R)$
(similar to the reduction of \finineqmaxname$(\F)$ to \ineqmaxname$(\F)$)
and Theorem~\ref{thm:lp-decid}.

\para{Case $\F = \Z$}

Decidability of \finnonnegsolvname$(\Int)$ follows by results of \cite{GHL22} and \cite{HR21}.

Let $(\vr{A},\vr{t})$ be an instance of \mbox{\finnonnegsolvname$(\Int)$}, where
$\vr A : B\times C \tofs \Z$, and
%Let $B$ and $C$ respectively be the indexing set of rows and columns of $\vr{A}$.
consider the set of column vectors 
\[
P = \setof{\vr{A}(\_,c)}{c \in C} \subseteq \GLin B
\]
of $\vr{A}$.
Then the system of equations $\vr{A} \cdot \vr{x} = \vr{t}$ has a finite 
non-negative integer solution if and only if 
\begin{align}\label{eq:tspan}
\vr{t} \in \Span{\Nat}{P}.
\end{align}
We rely on Theorem 3.3 of \cite{GHL22} which says that $\GLin B$ has an orbit-finite basis.
Let $\widehat{B} \subseteq \GLin B$ be such a basis.
This implies that there exists a linear isomorphism 
$
\varphi : \GLin B \to\Lin{\widehat{B}}.
$
In consequence, \eqref{eq:tspan} is equivalent to
\begin{align} \label{eq:tspan2}
\varphi(\vr{t}) \in \Span{\Nat}{\varphi(P)}.
\end{align}
By Remark 11.16 of \cite{HR21} we can compute a finite set of vectors 
$\set{\vr{t}'_1,\dots,\vr{t}'_k}\subseteq \Lin {\widehat B}$ and an orbit-finite subset 
$P'\subseteq\varphi(P)$ such 
that \eqref{eq:tspan2}
%$\varphi(\vr{t}) \in \Span{\Nat}{\varphi(P)}$ if and only if 
holds if and only if 
\begin{align} \label{eq:tspan3}
\vr{t}'_i \in \Span{\Int}{P'}
\end{align} 
for some $i\in\setto k$.
The question \eqref{eq:tspan3} is nothing but finitary integer solvability of an orbit-finite system of equations,
which is decidable using Theorem 6.1 of \cite{GHL22}.
%\end{proof}


%\subsection{Proofs missing in Section \ref{sec:weak}}
%
%\smallskip
%
%\subsection{Proofs missing in Section \ref{sec:poly}}
%
%\smallskip




\subsection{Proof of Theorem  \ref{thm:polydecid} (Section \ref{sec:poly})}
%{Proofs missing in Section \ref{sec:poly}}
\label{sec:poly-proofs}

%\begin{proof}
%[Proof of Theorem  \ref{thm:polydecid} (Section \ref{sec:poly})]
We show decidability of \polyineqsolvname %(cf.~Theorem~\ref{thm:polydecid}) 
by 
encoding the problem into real arithmetic, i.e., first-order theory of $(\R, +, \cdot, 0, 1, \leq)$.
%
We say that a real arithmetic formula $\varphi(x_1, \ldots, x_k)$ with free variables $x_1, \ldots, x_k$
\emph{defines} the set of all valuations $\set{x_1, \ldots, x_k}\to\R$ satisfying it. 
%
When the order of free variables is fixed, we naturally identify the set defined by $\varphi$ with a subset of $\R^k$.
%
\begin{claim} \label{claim:1var}
Every real arithmetic formula $\varphi(x)$ with one free variable, 
%possibly with constants, 
defines
%the set $\setof{x\in\R}{\psi(x) \text{ holds}}$ defined by $\psi(x)$ is
a finite union of (possibly infinite) disjoint intervals.
\end{claim}
\begin{proof}
%\subsection{Proof of Lemma \ref{lem:1var}}
%\slawek{to appendix}
By quantifier elimination~\cite{Tar51}, the formula $\varphi(x)$ is equivalent to a 
quantifier-free formula $\overline\varphi(x)$ with constants, namely $\varphi(x)$ and $\overline\varphi(x)$ define the same set.
Therefore $\overline\varphi(x)$ is a Boolean combination of inequalities
$
p(x) \geq 0,
$
for univariate polynomials $p\in\polyring \R x$, and validity of $\overline\varphi(x)$
depends only on the sign of $p(x)$, for (finitely many) polynomials that appear in $\overline\varphi(x)$.
This implies the claim.
%\qed
\end{proof}

Consider a fixed system $P$ of polynomially-parametrised inequalities over unknowns $x_1, \ldots, x_k$, and
let $n$ range over \emph{reals}, not just over nonnegative integers.
For each $n\in\R$, we  get the system $P(n)$ of linear inequalities with \emph{real} coefficients.  
Let 
%
\begin{align} \label{eq:sigmaP}
\sigma_P(n, x_1,\ldots, x_k)
\end{align}
%
 be the conjunction of inequalities in $P$, each of the form~\eqref{eq:polyeq};
it is thus a quantifier-free real arithmetic formula which says
that a  tuple $\vr x = x_1, \ldots, x_k$ is a solution of  $P(n)$. 
The existential real arithmetic formula $\psi(n) \ \equiv \ \prettyexists{\vr x}{\sigma_P(n, \vr x) }$ 
with one free variable $n$,
%\[
%\psi(n) \ \equiv \ \prettyexists{x_1 \ldots x_k} {%x_1 \geq 0 \land \ldots \land x_k\geq 0 \land 
%\sigma_P(n, x_1,\ldots,x_k) },
%\]
%which 
says that $P(n)$ has a real solution.
%
Thus 
\polyineqsolvname has positive answer exactly when $\psi(n)$ is true for some $n\in\Nat$.

Evaluating real arithmetic formulas of fixed quantifier alternation depth is doable in \EXPTIME~\cite{Ben-OrKR86}, 
\cite[Theorem 14.16]{ra-book}.
%
In order to decide \polyineqsolvname, the algorithm evaluates the closed formula
\begin{align*} %\label{eq:unb-form}
\prettyexists{\tilde n}{\prettyforall{n}{n >\tilde n {\implies} \psi(n)}}
%
%\prettyforlargeall{\!}{n}{\psi(n)}.
\end{align*}
%~\eqref{eq:unb-form} 
and answers positively if the formula is true.
Otherwise, we know that the set $D$ defined by $\psi$, being a finite union of intervals (cf.~Claim~\ref{claim:1var}), 
is bounded from above.
The algorithm computes an integer upper bound $m_0$ of $D$, by evaluating closed existential formulas
\[
\varphi_m \ \equiv \ \prettyexists n  {n > m \ \land \ \psi(n)},
\]
for increasing nonnegative integer constants $m = 0, 1, \ldots$,
until $\varphi_m$ eventually evaluates to false.
Finally, the algorithm evaluates the formula $\psi(m)$ for all nonnegative integers $m$ between $0$ and $m_0$,
and answers positively if $\psi(m)$ is true for some such $m$; otherwise the algorithm answers negatively.
%\TODO{what is the complexity in general case? Is there an exponential bound on $m_0$? This would yield \PSPACE}
%\end{proof}



\subsection{Proofs of Lemmas \ref{lem:wlog} and \ref{lem:wlogmax} (Sections \ref{sec:lp-decid-proof} and \ref{sec:max})}
%{Proofs missing in Sections \ref{sec:lp-decid-proof} and \ref{sec:max}}
\label{sec:lp-decid-proof-proofs}


%\begin{proof}
%[Proofs of Lemmas \ref{lem:wlog} and \ref{lem:wlogmax} (Sections \ref{sec:lp-decid-proof} and \ref{sec:max}, resp.)]
%
We sketch the proofs only, as they amount to a slightly tedious but entirely standard exercise in sets with atoms.

Consider an instance $(\vr{A},\vr{t},\vr{s})$ of the maximisation problem \finineqmaxname$(\R)$.
Let $S = \supp{\vr{A},\vr{t},\vr{s}}$, and let
$\vr A : B\times C\tofs \Z$.
Thus the row and column index sets $B$ and $C$ are necessarily supported by $S$.
We want to effectively transform the instance into another one $(\spread{\vr{A}},\spread{\vr{t}},\spread{\vr{s}})$, 
where the row and column index sets are 
disjoint unions of sets of the form $\otu \A \ell$ (non-repeating tuples of atoms of a fixed length), as in
\eqref{eq:eqorbits} in Section \ref{sec:prelimproof}.
%
%We want to show that there is another equivariant orbit-finite maximisation problem $(\spread{\vr{A}},\spread{\vr{t}},\spread{\vr{s}})$,
%where rows and columns of $\spread{\vr{A}}$ are indexed by equivariant straight sets (an equivariant straight orbit-finite set is a finite union of sets of the form $\A^{(k)}$)
Moreover, the transformation should preserve the supremum:
\begin{align} \label{eq:supsupprove}
\text{supremum}(\vr{A},\vr{t},\vr{s}) =
\text{supremum}(\spread{\vr{A}},\spread{\vr{t}},\spread{\vr{s}}).
\end{align}

Recall that we consider supremum of a maximisation problem to be $-\infty$ if the constraints are infeasible.
Therefore proving that two maximisation problems have the same supremum also proves that the underlying systems of inequalities are equisolvable.

%\arka{Should I say supremum or result of an optimisation problem}
%\slawek{supremum}

%We rely on the following result (which can be shown similarly as Lemma 6.2 in \cite{atombook};
%Lemma 6.3 therein for an effective version of this result):
%%
%\begin{lemma}\label{lem:rep}
%Let $S\subseteqfin\A$.
%For every $S$-orbit $U$  there are some $k, n \in \Nat$ and an $S$-supported 
%surjective function $$f : \otu {(\A \setminus S)} k \to U$$ such that
%for all $x \in U$, $|f^{-1}(x)| = n$.
%%%\todo{$f_X$ probably}\arka{Updated}
%%is finite and independent of $x$ (from now on we will denote this number by $n_U$).
%\end{lemma}
%%
%%\arka{Why not $=n$?}
%%
%%\noindent
%In the sequel we write $n(U)$ for the constant $n$, in order to indicate which $S$-orbit $U$ is considered.
%
%\arka{Suggestion starts here}
%

We proceed in two steps.
First we show that the row and column index sets $B$ and $C$
may be assumed to be disjoint unions of sets of the form $\otu {(\A\setminus S)} \ell$.
As mentioned in Section \ref{sec:results},
$B$ and $C$ are assumed to be given as finite union of $S$-orbits of the form $(\A \setminus S)^{(n)}/_G$ where $n \in \Nat$ and $G$ is a subgroup of $S_n$, the group of all permutations of the set $\set{1,\dots,n}$.
Consider the partition of $B$ and $C$ into $S$-orbits:
\[
B = B_1 \uplus \dots \uplus B_k
\qquad\qquad
C = C_1 \uplus \dots \uplus C_{\ell},
\]
where
\[
B_i = (\A \setminus S)^{(p_i)}/_{G_i}
\qquad \text{ and }
\qquad
C_j = (\A \setminus S)^{(q_j)}/_{H_j}
\]
for some $p_i,q_j \in \Nat$ and subgroups $G_i$ and $H_j$ of respectively $S_{p_i}$ and $S_{p_j}$.
Let $f_i$ and $g_j$ be the quotient maps
\[
f_i : (\A \setminus S)^{(p_i)} \to B_i \qquad
g_j : (\A \setminus S)^{(q_j)} \to C_j.
\]
Notice that for every $i,j$ and $x \in B_i$ and $y \in C_j$
\begin{equation}\label{eq:const size preimage}
|f_i^{-1}(x)| = |G_i| \qquad\qquad
|g_j^{-1}(y)| = |H_j|.
\end{equation}
%
%\arka{Suggestion ends here}
%
%
%We proceed in two steps. 
%As the first one we show that the row and column index sets $B$ and $C$
%may be assumed to be disjoint unions of of sets of the form $\otu {(\A\setminus S)} \ell$.
%Towards this, consider the partition of $B$ and $C$ into $S$-orbits:
%\[
%B = B_1 \uplus \dots \uplus B_k
%\qquad
%C = C_1 \uplus \dots \uplus C_{\ell},
%\]
%and apply Lemma \ref{lem:rep} to the $S$-orbits $B_i$ and $C_j$ to get:
%\[
%f_i : (\A \setminus S)^{(p_i)} \to B_i \qquad
%g_j : (\A \setminus S)^{(q_j)} \to C_j
%\]
%for $i \in\setto k$, $j\in \setto \ell$ and $p_i, \ldots, p_i, q_1, \ldots, q_j\in\Nat$.
%
%\arka{From the end of the suggestion till this point can be removed}
%
We put:
\begin{align} \label{eq:BCdef}
\begin{aligned}
B' \  = \ (\A \setminus S)^{(p_1)} \uplus \dots \uplus (\A \setminus S)^{(p_k)} \qquad\qquad
C' \  = \ (\A \setminus S)^{(q_1)} \uplus \dots \uplus (\A \setminus S)^{(q_{\ell})}
\end{aligned}
\end{align}
and define maps $f : B' \to B$ and $g: C' \to C$ by disjoint unions of $f_1, \ldots, f_k$ and
$g_1, \ldots, g_\ell$, respectively:
\[
f = f_1 \uplus \dots \uplus f_k
\qquad\qquad
g = g_1 \uplus \dots \uplus g_{\ell}.
\]
Both the maps are surjective.
We write $(f, g): B'\times C' \to B\times C$ for the product of the two maps.
Finally, we
define a matrix $\vr{A}' : (B' \times C') \tofs \Int$ and vectors $\vr{t}' \in \GLin {B'}$ and
$\vr{s}' \in \GLin{C}'$ by pre-composing with the above-defined maps:
\begin{align} \label{eq:precom}
\vr A' = \vr A \circ (f, g) 
%\vr{A}'(b,c) = \vr{A}(f(b),g(c))
\qquad\qquad
\vr{t}'      = \vr{t} \circ f  \qquad\qquad
\vr{s}'      = \vr{s} \circ g.
\end{align}
%
%We claim:
%
\begin{lemma}\label{lem:canon S supp}
$
\textnormal{supremum}(\vr{A},\vr{t},\vr{s}) =
\textnormal{supremum}(\vr{A}',\vr{t}',\vr{s}').
$
\end{lemma}
%
\begin{proof}
Define two functions $F : \GLin{C} \to \GLin{C'}$ and $G : \GLin{C'} \to \GLin{C}$ as follows:
\begin{align*}
F(\vr{x}) \  : \  c' \mapsto\  \frac {\vr{x}(g(c'))} {|H_i|}, \text{ where }g(c') \in C_i 
\qquad\qquad
G(\vr{x}') \  : \ c\ \mapsto \sum_{g(c') = c} \vr{x}'(c').
\end{align*}
%
%
%\arka{Replace $n(C_i)$ with $|H_j|$}
%
Both $F$ and $G$ are supported by $S$.
By the very definition of $F$ and $G$, together with \eqref{eq:const size preimage}, we deduce
the following two facts, assuming either
$\vr x' = F(\vr x)$ or $\vr x = G(\vr x')$, where $\vr x\in \GLin C$ and $\vr x' \in \GLin{C'}$.
First,
the value of $\mult {\vr{A}}  {\vr{x}}$ is well-defined if and only if
the value of $\mult{\vr{A}'} {\vr{x}'}$ is so, and in such case
%
\[
\mult {\vr{A}}  {\vr{x}} \geq \vr{t} \iff
\mult{\vr{A}'} {\vr{x}'} \geq \vr{t}'.
\]
Second,
the value of $\innerprod {\vr{s}}  {\vr{x}}$ is well defined if and only if 
the value of $\innerprod {\vr{s}'} {\vr{x}'}$ is so, and in such case
$
\innerprod {\vr{s}}  {\vr{x}} = \innerprod {\vr{s}'} {\vr{x}'}.
$
The two facts prove the lemma.
%
%\begin{enumerate}
%\item For any $\vr{x} \in \lin{C}$ we have $\vr{A} \cdot \vr{x} \geq \vr{t}$
%if and only if
%$\vr{A}' \cdot \varphi(\vr{x}) \geq \vr{t}'$.
%
%\item For any $\vr{x} \in \lin{C}$, $\vr{s} \cdot \vr{x}$ is well defined if and only if $\vr{s}' \cdot \varphi(\vr{x})$ is well defined and in that case they are equal.
%
%\item For any $\vr{x}' \in \lin{C'}$ we have $\vr{A} \cdot \psi(\vr{x}') \geq \vr{t}$
%if and only if
%$\vr{A}' \cdot \vr{x}' \geq \vr{t}'$.
%
%\item For any $\vr{x}' \in \lin{C'}$, $\vr{s} \cdot \psi(\vr{x})$ is well defined if and only if $\vr{s}' \cdot \vr{x}'$ is well defined and in that case they are equal.
%
%\end{enumerate}
%%
%This implies that $\varphi(\lin{C}) \subseteq \lin{C'}$ and $\psi(\lin{C'}) \subseteq \lin{C}$.
%These, combined with the above observations give us the following.
\end{proof}

%The above-described transformation is itself equivariant (it does not depend on $S$), and therefore
The instance $(\vr{A}',\vr{t}',\vr{s}')$ is supported by $S$.

%\arka{I think the statement above is confusing.
%Probably readers will understand if we simply mention that $(\vr{A}',\vr{t}',\vr{s}')$ is supported by $S$.}


As the second step we transform the instance $(\vr{A}',\vr{t}',\vr{s}')$ further so that
the row and column index sets $B$ and $C$ are disjoint unions of sets of the form $\otu \A \ell$.
Let $h : \A \to \A \setminus S$ be an arbitrarily chosen bijection.
Since atoms from $S$ do not appear in tuples belonging to $B'$ or $C'$, the map $h$
induces two further bijective maps 
\[
f: \spread{B} \to B' \qquad\qquad g : \spread{C} \to C',
\]
where
\[
\spread{B} = \A^{(p_1)} \uplus \dots \uplus \A^{(p_k)}
\qquad\qquad
\spread{C} = \A^{(q_1)} \uplus \dots \uplus \A^{(q_{\ell})}
\]
(cf.~\eqref{eq:BCdef}).
%
%\arka{Should we mention that these are bijections?}
%
%We overload the notation and denote these two maps by $H$ as well:
%\[
%H : \spread{B} \to B'
%\qquad
%H : \spread{C} \to C'.
%\]
We define a matrix $\spread{\vr{A}} : \spread B \times \spread C \tofs \Z$ 
and two vectors $\spread{\vr{t}} : \spread B \tofs \Z$ and 
$\spread{\vr{s}} : \spread C \tofs \Z$ by pre-composing with the two above-defined maps,
similarly as in \eqref{eq:precom}:
\[
\spread {\vr A} = \vr A' \circ (f,g) \qquad\qquad
\spread {\vr t} = \vr t' \circ f \qquad\qquad
\spread {\vr s} = \vr s' \circ g.
%\spread{\vr{A}}(b,c) = \vr{A}'(\cE^{-1}(b),\cE^{-1}(c)),
%\qquad
%\spread{\vr{t}}      = \vr{t}' \circ \cE^{-1} \ ,
%\spread{\vr{s}}      = \vr{s}' \circ \cE^{-1} \ .
\]
Knowing that $\vr A'$, $\vr t'$ and $\vr s'$ are all supported by $S$, we deduce that
the so defined instance $(\spread{\vr{A}},\spread{\vr{t}}, \spread{\vr s})$ is equivariant and independent from
the choice of the bijection $h : \A \to \A\setminus S$.
The size blowup is exponential only in atom dimension of $(\vr A, \vr t, \vr s)$, and hence polynomial
when atom dimension if fixed.

\begin{lemma}\label{lem:equivariant sys}
$
\textnormal{supremum}(\vr{A}',\vr{t}',\vr{s}') =
\textnormal{supremum}(\spread{\vr{A}},\spread{\vr{t}},\spread{\vr{s}}).
$
\end{lemma}
\begin{proof}
Similarly as before,
assuming $\vr x' = g(\spread{\vr x})$ for some vectors $\vr x' \in \GLin{C'}$ and $\spread{\vr x} \in \GLin{\spread C}$,
we deduce the following two facts.
First,
the value of $\mult {\vr{A}'}  {\vr{x}'}$ is well-defined if and only if
the value of $\mult{\spread{\vr{A}}} {\spread{\vr{x}}}$ is so, and in such case
%
\[
\mult {\vr{A}'}  {\vr{x}'} \geq \vr{t'} \iff
\mult{\spread{\vr{A}}} \spread{{\vr{x}}} \geq \spread{\vr{t}}.
\]
Second,
the value of $\innerprod {\vr{s}'}  {\vr{x}'}$ is well defined if and only if 
the value of $\innerprod {\spread{\vr{s}}} {\spread{\vr{x}}}$ is so, and in such case
$
\innerprod {\vr{s}'}  {\vr{x}'} = \innerprod {\spread{\vr{s}}} {\spread{\vr{x}}}.
$
The two facts prove the lemma.
%
%%
%\begin{enumerate}
%\item For any $\vr{x}' \in \lin{C'}$ we have $\vr{A}' \cdot \vr{x}' \geq \vr{t}'$
%if and only if
%$\spread{\vr{A}} \cdot \cE(\vr{x}') \geq \spread{\vr{t}}$.
%
%\item For any $\vr{x}' \in \lin{C}$, $\vr{s}' \cdot \vr{x}'$ is well defined if and only if $\spread{\vr{s}} \cdot \cE(\vr{x}')$ is well defined and in that 
%
%\item For any $\spread{\vr{x}} \in \lin{\spread{C}}$ we have $\vr{A}' \cdot \cE^{-1}(\spread{\vr{x}}) \geq \vr{t}'$
%if and only if
%$\spread{\vr{A}} \cdot \spread{\vr{x}} \geq \spread{\vr{t}}$.
%
%\item For any $\spread{\vr{x}} \in \lin{\spread{C}}$, $\vr{s}' \cdot \cE^{-1}(\spread{\vr{x}})$ is well defined if and only if $\spread{\vr{s}} \cdot 
%\spread{\vr{x}}$ is well defined and in that case they are equal.
%\end{enumerate}
%
\end{proof}

The last two lemmas prove Lemmas \ref{lem:wlog} and \ref{lem:wlogmax}.
%
%\end{proof}


